% Barycentres et lignes de niveau — Mathématiques 1ère TI — Cours signé OnBuch+
% Programme officiel MINESEC. Compiler avec : tectonic N12-barycentres-lignes-niveau.tex
\newcommand{\DOCMATIERE}{Mathématiques}
\newcommand{\DOCNIVEAU}{1ère TI}
\newcommand{\DOCMODULE}{Configurations et transformations élémentaires du plan}
\newcommand{\DOCLECON}{N12}
\newcommand{\DOCTITRE}{Barycentres et lignes de niveau}
% OnBuch+ — préambule « cours signés » ENRICHI (compilé avec tectonic / XeLaTeX)
% Système visuel « Pop » d'OnBuch+ : fond crème (jamais blanc), bordures encre
% épaisses, ombres « dures » décalées, titres Archivo Black en capitales.
% Version MATHÉMATIQUES : graphes (pgfplots/tikz), boîtes pédagogiques
% (définition, propriété, méthode, attention, le savais-tu, exercice résolu,
% exercice, TP, bilan), page de garde et signature OnBuch+ ancrée en bas.
%
% Macros attendues dans main.tex AVANT \input{preamble} :
%   \DOCMATIERE  \DOCNIVEAU  \DOCMODULE  \DOCLECON (numéro)  \DOCTITRE
\documentclass[11pt]{article}

\usepackage{fontspec}
\usepackage[french]{babel}
\usepackage[a4paper,margin=2cm,top=3.4cm,bottom=2.8cm,headheight=1.4cm,footskip=1.3cm]{geometry}
\usepackage{amsmath,amssymb,amsfonts}
\usepackage{mathtools} % \xmapsto, \coloneqq... souvent utilisés par les modèles
\usepackage{cancel} % simplifications barrées
% \abs{z} (module, valeur absolue) et \norm{u} (norme), utilisés par les modèles
\providecommand{\abs}{}\let\abs\relax\DeclarePairedDelimiter{\abs}{\lvert}{\rvert}
\providecommand{\norm}{}\let\norm\relax\DeclarePairedDelimiter{\norm}{\lVert}{\rVert}
\usepackage[table]{xcolor} % [table] : fournit aussi \rowcolors (alternance de lignes)
\usepackage{pagecolor}
\usepackage{tikz}
\usetikzlibrary{arrows.meta,positioning,calc,shapes.geometric,shapes.misc,shapes.arrows,decorations.pathmorphing,decorations.pathreplacing,decorations.markings,patterns,shadows,mindmap,trees,fit,backgrounds,matrix,intersections,angles,quotes}
\usepackage{pgfplots}
\usepackage{pgf-pie} % diagrammes circulaires \pie (statistiques)
\pgfplotsset{compat=1.18}
\usepgfplotslibrary{fillbetween,groupplots} % aires entre courbes (calcul intégral)
\usepackage{enumitem}
\usepackage{fancyhdr}
% [most] charge toutes les bibliothèques tcolorbox (listings, minted, raster,
% documentation...) et gonfle inutilement la mémoire TeX sur un document long
% (~300 boîtes) : on ne charge que ce qui est réellement utilisé.
\usepackage[skins,breakable]{tcolorbox}
\usepackage{graphicx}
\usepackage{colortbl}
\usepackage{booktabs}
\usepackage{tabularx}
\usepackage{diagbox} % case d'en-tête diagonale des tableaux à double entrée
% Colonnes extensibles centrée / gauche / droite (C, L, R), souvent utilisées par les modèles
\newcolumntype{C}{>{\centering\arraybackslash}X}
\newcolumntype{L}{>{\raggedright\arraybackslash}X}
\newcolumntype{R}{>{\raggedleft\arraybackslash}X}
\usepackage{array}
\usepackage{multicol}
\usepackage{siunitx}
% parse-numbers=false : accepte aussi \SI{2,5 \times 10^{-3}}{...}
\sisetup{locale=FR,output-decimal-marker={,},group-minimum-digits=4,parse-numbers=false}
\DeclareSIUnit\ohm{\ensuremath{\symup{\Omega}}} % U+2126 absent de la police
\usepackage{qrcode}
% \degree : commande fréquemment utilisée par les modèles pour un angle ou une
% température en dehors de \SI{}{} (non fournie par les paquets chargés).
\providecommand{\degree}{^{\circ}}
% \qed (fin de démonstration), utilisé par les modèles sans amsthm
\providecommand{\qed}{\hfill$\square$}
% \barre : double barre des valeurs interdites dans un tableau de variations (tkz-tab)
\providecommand{\barre}{\ensuremath{\|}}
% environnement proof (amsthm non chargé) utilisé par les modèles
\ifdefined\proof\else\newenvironment{proof}[1][Démonstration]{\par\noindent\textit{#1.}\ }{\qed\par}\fi
\usepackage{lastpage}
\usepackage{hyperref}
\hypersetup{hidelinks}

% ============================================================
% Palette « Pop » OnBuch+ — mêmes hex que la classe P (plus_ui.dart)
% ============================================================
\definecolor{poporange}{HTML}{F59321}
\definecolor{popdark}{HTML}{E07A0C}
\definecolor{popcream}{HTML}{FFF6E8}
\definecolor{popink}{HTML}{1C1714}
\definecolor{poppurple}{HTML}{7A5AE0}
\definecolor{popgreen}{HTML}{2E9E6A}
\definecolor{poppink}{HTML}{E0457B}
\definecolor{popblue}{HTML}{2F7DE1}
\definecolor{popgold}{HTML}{C98A12}
\definecolor{popmuted}{HTML}{8A7F76}
\definecolor{popline}{HTML}{EADFD2}
\definecolor{popgray}{HTML}{8A7F76}
\colorlet{popgrey}{popgray}
% Alias tolérants (le modèle nomme parfois les couleurs différemment)
\colorlet{popwhite}{white}
\colorlet{popblack}{popink}
\colorlet{popred}{poppink}
\colorlet{popyellow}{popgold}
% Teintes claires pour fonds de tableaux / zones
\colorlet{popblueL}{popblue!10!white}
\colorlet{poporangeL}{poporange!14!white}
\colorlet{popgreenL}{popgreen!12!white}
\colorlet{poppurpleL}{poppurple!12!white}
\colorlet{poppinkL}{poppink!12!white}
\colorlet{popgoldL}{popgold!14!white}

\pagecolor{popcream}

% ============================================================
% Typographie
% ============================================================
\defaultfontfeatures{Ligatures=TeX}
\setmainfont{PlusJakartaSans-Medium.ttf}[Path=../../../fonts/,BoldFont=PlusJakartaSans-Bold.ttf]
% Maths en sans-serif (Fira Math) avec chiffres et lettres droites de Plus
% Jakarta, pour que les formules se fondent dans le texte.
\usepackage[math-style=ISO,bold-style=ISO]{unicode-math}
\setmathfont{FiraMath-Regular.otf}[Path=../../../fonts/]
\setmathfont{PlusJakartaSans-Medium.ttf}[Path=../../../fonts/,range={up/num,up/{Latin,latin},\mathup}]
\usepackage{icomma} % virgule décimale sans espace en mode math
% Caractères mathématiques Unicode absents des polices de texte (Plus Jakarta,
% Archivo Black) : rendus par leur équivalent mathématique, en texte comme en
% formule — sinon ils s'affichent en carré vide (ex. « Arithmétique dans ℤ »).
\usepackage{newunicodechar}
\newunicodechar{ℝ}{\ensuremath{\mathbb{R}}}
\newunicodechar{ℤ}{\ensuremath{\mathbb{Z}}}
\newunicodechar{ℂ}{\ensuremath{\mathbb{C}}}
\newunicodechar{ℕ}{\ensuremath{\mathbb{N}}}
\newunicodechar{ℚ}{\ensuremath{\mathbb{Q}}}
\newunicodechar{𝔻}{\ensuremath{\mathbb{D}}}
\newunicodechar{α}{\ensuremath{\alpha}}
\newunicodechar{β}{\ensuremath{\beta}}
\newunicodechar{γ}{\ensuremath{\gamma}}
\newunicodechar{δ}{\ensuremath{\delta}}
\newunicodechar{ε}{\ensuremath{\varepsilon}}
\newunicodechar{θ}{\ensuremath{\theta}}
\newunicodechar{λ}{\ensuremath{\lambda}}
\newunicodechar{μ}{\ensuremath{\mu}}
\newunicodechar{σ}{\ensuremath{\sigma}}
\newunicodechar{τ}{\ensuremath{\tau}}
\newunicodechar{φ}{\ensuremath{\varphi}}
\newunicodechar{ψ}{\ensuremath{\psi}}
\newunicodechar{ω}{\ensuremath{\omega}}
\newunicodechar{Σ}{\ensuremath{\Sigma}}
% majuscules produites par \MakeUppercase dans les titres (α-aminés -> Α)
\newunicodechar{Α}{\ensuremath{\alpha}}
\newunicodechar{Β}{\ensuremath{\beta}}
\newunicodechar{Γ}{\ensuremath{\Gamma}}
\newunicodechar{Δ}{\ensuremath{\Delta}}
\newunicodechar{Ε}{\ensuremath{\varepsilon}}
\newunicodechar{Θ}{\ensuremath{\Theta}}
\newunicodechar{Λ}{\ensuremath{\Lambda}}
\newunicodechar{Μ}{\ensuremath{\mu}}
\newunicodechar{Τ}{\ensuremath{\tau}}
\newunicodechar{Φ}{\ensuremath{\Phi}}
\newunicodechar{Ψ}{\ensuremath{\Psi}}
\newunicodechar{Ω}{\ensuremath{\Omega}}
\newunicodechar{↦}{\ensuremath{\mapsto}}
\newunicodechar{∈}{\ensuremath{\in}}
\newunicodechar{∉}{\ensuremath{\notin}}
\newunicodechar{∧}{\ensuremath{\wedge}}
\newunicodechar{⇔}{\ensuremath{\Leftrightarrow}}
\newunicodechar{⟺}{\ensuremath{\Longleftrightarrow}}
\newunicodechar{⇒}{\ensuremath{\Rightarrow}}
\newunicodechar{⟹}{\ensuremath{\Longrightarrow}}
\newunicodechar{∘}{\ensuremath{\circ}}
\newunicodechar{∪}{\ensuremath{\cup}}
\newunicodechar{∩}{\ensuremath{\cap}}
\newunicodechar{≡}{\ensuremath{\equiv}}
\newunicodechar{⊂}{\ensuremath{\subset}}
\newunicodechar{⊥}{\ensuremath{\perp}}
\newunicodechar{∃}{\ensuremath{\exists}}
\newunicodechar{∀}{\ensuremath{\forall}}
\newunicodechar{∛}{\ensuremath{\sqrt[3]{\,}}}
\newunicodechar{∜}{\ensuremath{\sqrt[4]{\,}}}
\newunicodechar{⃗}{\ensuremath{{}^{\to}}}
\newunicodechar{ⁿ}{\ifmmode^{n}\else\textsuperscript{\ensuremath{n}}\fi}
\newunicodechar{ˣ}{\ifmmode^{x}\else\textsuperscript{\ensuremath{x}}\fi}
\newunicodechar{ᵇ}{\ifmmode^{b}\else\textsuperscript{\ensuremath{b}}\fi}
\newunicodechar{ʳ}{\ifmmode^{r}\else\textsuperscript{\ensuremath{r}}\fi}
\newunicodechar{ᵘ}{\ifmmode^{u}\else\textsuperscript{\ensuremath{u}}\fi}
\newunicodechar{ʸ}{\ifmmode^{y}\else\textsuperscript{\ensuremath{y}}\fi}
\newunicodechar{ᵃ}{\ifmmode^{a}\else\textsuperscript{\ensuremath{a}}\fi}
\newunicodechar{ᵉ}{\ifmmode^{e}\else\textsuperscript{\ensuremath{e}}\fi}
\newunicodechar{ᵏ}{\ifmmode^{k}\else\textsuperscript{\ensuremath{k}}\fi}
\newunicodechar{ᵖ}{\ifmmode^{p}\else\textsuperscript{\ensuremath{p}}\fi}
\newunicodechar{ᴺ}{\ifmmode^{N}\else\textsuperscript{\ensuremath{N}}\fi}
\newunicodechar{ᶜ}{\ifmmode^{c}\else\textsuperscript{\ensuremath{c}}\fi}
\newunicodechar{ʰ}{\ifmmode^{h}\else\textsuperscript{\ensuremath{h}}\fi}
\newunicodechar{ᵠ}{\ifmmode^{q}\else\textsuperscript{\ensuremath{q}}\fi}
\newunicodechar{ₙ}{\ifmmode_{n}\else\textsubscript{\ensuremath{n}}\fi}
\newunicodechar{ₐ}{\ifmmode_{a}\else\textsubscript{\ensuremath{a}}\fi}
\newunicodechar{ᵢ}{\ifmmode_{i}\else\textsubscript{\ensuremath{i}}\fi}
\newunicodechar{ᵣ}{\ifmmode_{r}\else\textsubscript{\ensuremath{r}}\fi}
\newunicodechar{ₖ}{\ifmmode_{k}\else\textsubscript{\ensuremath{k}}\fi}
\newunicodechar{ᵦ}{\ifmmode_{\beta}\else\textsubscript{\ensuremath{\beta}}\fi}
\newunicodechar{ₓ}{\ifmmode_{x}\else\textsubscript{\ensuremath{x}}\fi}
\newfontfamily\popdisplay{ArchivoBlack-Regular.ttf}[Path=../../../fonts/]
\newfontfamily\poplogo{SpaceGrotesk-Bold.ttf}[Path=../../../fonts/]
\newfontfamily\popsemi{PlusJakartaSans-SemiBold.ttf}[Path=../../../fonts/]
\newfontfamily\popxbold{PlusJakartaSans-ExtraBold.ttf}[Path=../../../fonts/]
\newfontfamily\popxboldls{PlusJakartaSans-ExtraBold.ttf}[Path=../../../fonts/,LetterSpace=3.0]

\setlist[enumerate]{leftmargin=1.7em,itemsep=5pt,topsep=4pt}
\setlist[itemize]{leftmargin=1.7em,itemsep=4pt,topsep=4pt,label={\color{poporange}\textbullet}}
\setlength{\parindent}{0pt}
\setlength{\parskip}{5pt}
\renewcommand{\arraystretch}{1.25}

% pgfplots : style Pop par défaut
\pgfplotsset{
  every axis/.append style={axis line style={popink,line width=1pt},tick style={popink},
    grid style={popline,line width=0.5pt},label style={font=\small},tick label style={font=\footnotesize}},
  popcurve/.style={poporange,line width=1.8pt,no markers,smooth},
}
% Même style disponible dans un \draw TikZ simple (hors pgfplots)
\tikzset{popcurve/.style={poporange,line width=1.8pt}}

% ============================================================
% Pill / squiggle
% ============================================================
\newcommand{\poppill}[3]{%
  \tikz[baseline=-0.55ex]\node[draw=popink,line width=1.2pt,rounded corners=2.6pt,%
    fill=#2,inner xsep=6.5pt,inner ysep=3pt]{%
    \popxboldls\fontsize{7}{7}\selectfont\color{#3}\MakeUppercase{#1}};%
}
\newcommand{\popsquiggle}[1]{%
  \tikz[baseline=-0.4ex]{
    \draw[#1,line width=2.6pt,line cap=round]
      plot [smooth] coordinates {(0,0.09) (0.34,0.01) (0.68,0.21) (1.02,0.01) (1.36,0.10)};
  }%
}
% Mot-clé surligné (marqueur orange sous le texte)
\newcommand{\cle}[1]{{\popxbold\color{popdark}#1}}

% ============================================================
% En-tête / pied
% ============================================================
\pagestyle{fancy}
\fancyhf{}
\renewcommand{\headrulewidth}{0pt}
\renewcommand{\footrulewidth}{0pt}
\lhead{\raisebox{-0.15ex}{\poplogo\fontsize{13}{13}\selectfont\color{popink}OnBuch\color{poporange}+}}
\rhead{\poppill{\DOCMATIERE\ \textbullet\ \DOCNIVEAU\ \textbullet\ Leçon \DOCLECON}{white}{popink}}
\lfoot{\popxbold\fontsize{7.6}{7.6}\selectfont\color{popmuted}\MakeUppercase{Cours signé OnBuch+}\ \textbullet\ {\popsemi\DOCTITRE}}
\rfoot{\tikz[baseline=-0.6ex]\node[rounded corners=8.5pt,fill=popink,minimum height=17pt,minimum width=17pt,inner xsep=5pt,inner sep=0]{\popxbold\fontsize{8}{8}\selectfont\color{popcream}\thepage\,/\,\pageref*{LastPage}};}

% ============================================================
% Titres
% ============================================================
\newcommand{\chaptitre}[2]{%
  \begin{tcolorbox}[enhanced,colback=poporange,colframe=popink,boxrule=2.6pt,arc=11pt,
      drop shadow=popink,left=15pt,right=15pt,top=12pt,bottom=11pt,before skip=3pt,after skip=18pt]
    \poppill{#2}{popink}{popcream}\par\vspace{9pt}
    {\raggedright\hyphenpenalty=10000\exhyphenpenalty=10000\popdisplay\fontsize{22}{24}\selectfont\color{popcream}\MakeUppercase{#1}\par}
  \end{tcolorbox}
}
% Section numérotée : pastille orange avec le numéro + titre Archivo
\newcounter{popsec}
\newcommand{\coursec}[1]{%
  \stepcounter{popsec}\setcounter{popsub}{0}%
  \par\vspace{16pt}\noindent
  \begin{minipage}{\linewidth}
  \tikz[baseline=-0.7ex]\node[draw=popink,line width=1.6pt,fill=poporange,rounded corners=4pt,minimum size=22pt,inner sep=0]{\popdisplay\fontsize{12}{12}\selectfont\color{popcream}\thepopsec};%
  \hspace{8pt}{\popdisplay\fontsize{15}{17}\selectfont\color{popink}\MakeUppercase{#1}}\par
  \vspace{4pt}\noindent{\color{poporange}\rule{36pt}{3.6pt}}
  \end{minipage}\par\nopagebreak\vspace{8pt}
}
\newcounter{popsub}
\newcommand{\courssub}[1]{%
  \stepcounter{popsub}%
  \par\vspace{9pt}\noindent{\popxbold\fontsize{11.5}{13}\selectfont\color{popink}{\color{poporange}\thepopsec.\thepopsub}\hspace{6pt}#1}\par\nopagebreak\vspace{4pt}
}

% ============================================================
% Boîtes pédagogiques — même recette : fond blanc, bordure épaisse colorée,
% ombre dure encre, badge plein. Titre optionnel [#1].
% ============================================================
\tcbset{popbase/.style={breakable,enhanced,colback=white,boxrule=2.3pt,arc=9pt,
  shadow={3pt}{-3pt}{0pt}{popink},left=14pt,right=14pt,top=11pt,bottom=11pt,before skip=11pt,after skip=15pt,
  fontupper=\popsemi\fontsize{10.6}{14.5}\selectfont\color{popink},
  fontlower=\popsemi\fontsize{10.6}{14.5}\selectfont\color{popink}}}
% Coche dessinée (le glyphe ✓ n'existe pas dans les polices du document)
\newcommand{\popcheck}[1]{\tikz[baseline=-0.5ex]\draw[#1,line width=1.6pt,line cap=round,line join=round](-0.13,0.0)--(-0.04,-0.09)--(0.13,0.1);}
\newcommand{\popboxhead}[3]{\poppill{#1}{#2}{white}\ifx\relax#3\relax\else\hspace{6pt}{\popxbold\fontsize{10.6}{12}\selectfont\color{popink}#3}\fi\par\vspace{7pt}}

\newtcolorbox{aretenir}[1][]{popbase,colframe=popgreen,before upper={\popboxhead{À retenir}{popgreen}{#1}}}
\newtcolorbox{exemplebox}[1][]{popbase,colframe=poporange,before upper={\popboxhead{Exemple}{poporange}{#1}}}
\newtcolorbox{definition}[1][]{popbase,colframe=popblue,before upper={\popboxhead{Définition}{popblue}{#1}}}
\newtcolorbox{propriete}[1][]{popbase,colframe=poppurple,before upper={\popboxhead{Propriété}{poppurple}{#1}}}
\newtcolorbox{methode}[1][]{popbase,colframe=popgold,before upper={\popboxhead{Méthode}{popgold}{#1}}}
\newtcolorbox{attention}[1][]{popbase,colframe=poppink,before upper={\popboxhead{Attention}{poppink}{#1}}}
\newtcolorbox{experience}[1][]{popbase,colframe=popblue,colback=popblueL,before upper={\popboxhead{Expérience}{popblue}{#1}}}
% « Le savais-tu ? » : carte violette pleine (contexte, culture, Cameroun)
\newtcolorbox{savaistu}[1][]{popbase,colback=poppurple,colframe=popink,
  fontupper=\popsemi\fontsize{10.4}{14.2}\selectfont\color{white},
  before upper={\poppill{Le savais-tu\,?}{white}{poppurple}\ifx\relax#1\relax\else\hspace{6pt}{\popxbold\color{white}#1}\fi\par\vspace{7pt}}}
% Exercice résolu : énoncé en haut, \tcblower puis la solution sur fond crème
\newcounter{popexo}\newcounter{popexr}
\newtcolorbox{exoresolu}[1][]{popbase,colframe=poporange,colbacklower=popcream,
  segmentation style={popink,line width=1.2pt,dashed},
  before upper={\stepcounter{popexr}\popboxhead{Exercice résolu \thepopexr}{poporange}{#1}},
  before lower={\poppill{Solution}{popink}{popcream}\par\vspace{6pt}}}
% Exercice (non résolu en place) — la correction est dans la partie Corrigés
\newtcolorbox{exercice}[1][]{popbase,colframe=popink,
  before upper={\stepcounter{popexo}\popboxhead{Exercice \thepopexo}{popink}{#1}}}
% Corrigé
\newtcolorbox{corrige}[1][]{popbase,colframe=popgreen,colback=popgreenL,
  before upper={\popboxhead{Corrigé}{popgreen}{#1}}}
% Objectifs / prérequis / bilan
\newtcolorbox{objectifs}{popbase,colframe=popink,colback=poporangeL,
  before upper={\popboxhead{Objectifs de la leçon}{popink}{}}}
\newtcolorbox{prerequis}{popbase,colframe=popink,colback=popblueL,
  before upper={\popboxhead{Prérequis}{popblue}{}}}
\newtcolorbox{bilan}{popbase,colframe=popink,colback=white,boxrule=2.8pt,
  before upper={\popboxhead{Fiche bilan}{poporange}{L'essentiel à connaître pour l'examen}}}
% Figure encadrée avec légende
\newtcolorbox{popfigure}[1][]{enhanced,breakable=false,colback=white,colframe=popline,boxrule=1.4pt,arc=8pt,
  left=10pt,right=10pt,top=10pt,bottom=8pt,before skip=10pt,after skip=12pt,halign=center,
  lower separated=false,halign lower=center,
  fontlower=\popsemi\fontsize{9}{11.5}\selectfont\color{popmuted},
  #1}
\newcommand{\legende}[1]{\par\vspace{6pt}{\popsemi\fontsize{9}{11.5}\selectfont\color{popmuted}#1}}

% ============================================================
% Signature OnBuch+
% ============================================================
\newcommand{\poplogoinline}[1]{{\poplogo\fontsize{#1}{#1}\selectfont\color{popink}OnBuch\color{poporange}+}}
\newcommand{\signatureonbuch}{%
  \begin{tcolorbox}[enhanced,colback=popink,colframe=popink,boxrule=2.2pt,arc=10pt,
      drop shadow=poporange,left=14pt,right=14pt,top=10pt,bottom=10pt,before skip=0pt,after skip=0pt]
    \tikz[baseline=-0.7ex]\node[circle,fill=popgreen,draw=popcream,line width=1.3pt,minimum size=22pt,inner sep=0]{\popcheck{white}};%
    \hspace{9pt}\begin{minipage}[c]{0.62\linewidth}
      {\popxbold\fontsize{12}{13}\selectfont\color{popcream}Signé\ \textbullet\ {\poplogo OnBuch\color{poporange}+}}\par\vspace{2pt}
      {\raggedright\popsemi\fontsize{8}{10.5}\selectfont\color{popline}\MakeUppercase{Équipe pédagogique OnBuch+}\\\MakeUppercase{Conforme au programme officiel MINESEC}\par}
    \end{minipage}\hfill
    {\poplogo\fontsize{22}{22}\selectfont\color{popcream}OnBuch\color{poporange}+}
  \end{tcolorbox}%
}

% ============================================================
% Page de garde
% #1 titre · #2 matière · #3 niveau · #4 module · #5 n° leçon · #6 durée ·
% #7 description / objectifs (texte ou itemize)
% ============================================================
\newcommand{\pagegarde}[7]{%
  \thispagestyle{empty}
  \newgeometry{a4paper,margin=1.7cm,top=1.6cm,bottom=1.4cm}%
  \begin{tikzpicture}[remember picture,overlay]
    \fill[poporange!18] ([xshift=-3.2cm,yshift=-3.6cm]current page.north east) circle (4.6cm);
    \fill[poppurple!14] ([xshift=2.4cm,yshift=7.6cm]current page.south west) circle (3.8cm);
  \end{tikzpicture}%
  {\poplogo\fontsize{30}{30}\selectfont\color{popink}OnBuch\color{poporange}+}\hfill
  \raisebox{0.6ex}{\poppill{Cours signé \textbullet\ Premium}{popink}{popcream}}\par
  \vspace{1pt}\popsquiggle{poppurple}\par
  \vspace{8pt}
  \begin{tcolorbox}[enhanced,colback=poporange,colframe=popink,boxrule=2.8pt,arc=13pt,
      drop shadow=popink,left=18pt,right=18pt,top=12pt,bottom=13pt,after skip=10pt]
    \poppill{#2 \textbullet\ #3}{popink}{popcream}\hspace{5pt}\poppill{#4}{white}{popink}\par\vspace{8pt}
    {\popdisplay\fontsize{14}{14}\selectfont\color{popink}LEÇON #5}\par\vspace{4pt}
    {\raggedright\hyphenpenalty=10000\exhyphenpenalty=10000\popdisplay\fontsize{25}{27}\selectfont\color{popcream}\MakeUppercase{#1}\par}
    \vspace{8pt}
    {\popxbold\fontsize{9.5}{9.5}\selectfont\color{popink}\MakeUppercase{Durée conseillée : #6}}
  \end{tcolorbox}
  \begin{tcolorbox}[enhanced,colback=white,colframe=popink,boxrule=2.2pt,arc=10pt,
      drop shadow=popink,left=16pt,right=16pt,top=10pt,bottom=10pt,after skip=8pt]
    \poppill{Dans cette leçon}{poppurple}{white}\par\vspace{5pt}
    \popsemi\fontsize{9.3}{12.6}\selectfont\color{popink}#7
  \end{tcolorbox}
  \noindent\begin{minipage}[t]{0.487\linewidth}
    \begin{tcolorbox}[enhanced,colback=popgreen,colframe=popink,boxrule=2.2pt,arc=10pt,
        drop shadow=popink,left=11pt,right=11pt,top=10pt,bottom=10pt,height=3.1cm,valign=center]
      \tcbox[colback=white,colframe=popink,boxrule=1.3pt,arc=5pt,boxsep=0pt,nobeforeafter,
        left=3pt,right=3pt,top=3pt,bottom=3pt,valign=center]{\qrcode[height=1.9cm]{https://play.google.com/store/apps/details?id=com.luvvix.onbuch&hl=fr}}%
      \hspace{8pt}\begin{minipage}[c]{0.5\linewidth}
        {\popdisplay\fontsize{11}{12.5}\selectfont\color{white}\MakeUppercase{Télécharge OnBuch}}\par\vspace{4pt}
        {\raggedright\popsemi\fontsize{8.4}{11}\selectfont\color{white}Gratuit \textbullet\ Google Play\\Tuteur IA, annales, concours, résultats\par}
      \end{minipage}
    \end{tcolorbox}
  \end{minipage}\hfill
  \begin{minipage}[t]{0.487\linewidth}
    \begin{tcolorbox}[enhanced,colback=poppurple,colframe=popink,boxrule=2.2pt,arc=10pt,
        drop shadow=popink,left=11pt,right=11pt,top=10pt,bottom=10pt,height=3.1cm,valign=center]
      \tikz[baseline=-0.7ex]\node[circle,fill=white,draw=popink,line width=1.3pt,minimum size=1.6cm,inner sep=0]{\popdisplay\fontsize{24}{24}\selectfont\color{poppurple}+};%
      \hspace{8pt}\begin{minipage}[c]{0.6\linewidth}
        {\popdisplay\fontsize{11}{12.5}\selectfont\color{white}\MakeUppercase{Passe à OnBuch+}}\par\vspace{4pt}
        {\raggedright\popsemi\fontsize{8.4}{11}\selectfont\color{white}Tous les cours signés, Tuteur IA illimité, fiches PDF — onglet OnBuch+ de l'app\par}
      \end{minipage}
    \end{tcolorbox}
  \end{minipage}
  \vspace{8pt}
  \popcontenu
  \vfill
  \signatureonbuch
  \restoregeometry
  \newpage
}

% Rangée « Ce cours contient » (page de garde)
\newcommand{\poptile}[3]{% #1 couleur #2 gros texte #3 légende
  \begin{tcolorbox}[enhanced,colback=#1,colframe=popink,boxrule=2pt,arc=9pt,shadow={2.5pt}{-2.5pt}{0pt}{popink},
      left=6pt,right=6pt,top=9pt,bottom=9pt,halign=center,valign=center,height=1.75cm,width=\linewidth,nobeforeafter]
    {\popdisplay\fontsize{12.5}{13}\selectfont\color{white}\MakeUppercase{#2}}\par\vspace{3pt}
    {\centering\popsemi\fontsize{7.8}{9.5}\selectfont\color{white}#3\par}
  \end{tcolorbox}}
\newcommand{\popcontenu}{%
  \par\noindent{\popxboldls\fontsize{7.5}{7.5}\selectfont\color{popmuted}CE COURS SIGNÉ CONTIENT}\par\vspace{7pt}
  \noindent\begin{minipage}[t]{0.235\linewidth}\poptile{popblue}{Cours}{complet, illustré, pas à pas}\end{minipage}\hfill
  \begin{minipage}[t]{0.235\linewidth}\poptile{poporange}{Méthodes}{et exercices résolus}\end{minipage}\hfill
  \begin{minipage}[t]{0.235\linewidth}\poptile{poppink}{Exercices}{type examen + corrigés}\end{minipage}\hfill
  \begin{minipage}[t]{0.235\linewidth}\poptile{popgreen}{Fiche bilan}{l'essentiel à retenir}\end{minipage}\par}

% Sommaire
\newtcolorbox{sommaire}{enhanced,colback=white,colframe=popink,boxrule=2.2pt,arc=10pt,
  drop shadow=popink,left=18pt,right=18pt,top=13pt,bottom=13pt,before skip=6pt,after skip=20pt,
  before upper={\poppill{Sommaire}{poppurple}{white}\par\vspace{9pt}\popsemi\fontsize{10.6}{14.5}\selectfont\color{popink}}}

% Signature de fin de cours — poussée tout en bas de la dernière page
\newcommand{\findecours}{%
  \par\vfill\nopagebreak
  \begin{center}\popsquiggle{poporange}\end{center}\vspace{4pt}
  \signatureonbuch
}
\AtBeginDocument{\def\llbracket{\mathopen{[\mkern-3.5mu[}}\def\rrbracket{\mathclose{]\mkern-3.5mu]}}} % intervalles d'entiers (glyphe ⟦ absent de la police)
% Glyphes absents de Fira Math : redéfinis à partir de glyphes disponibles
\AtBeginDocument{%
  \def\setminus{\mathbin{\backslash}}\let\smallsetminus\setminus
  \def\vdots{\mathord{\vbox{\baselineskip3.2pt\lineskiplimit0pt\kern4pt\hbox{.}\hbox{.}\hbox{.}}}}%
  \def\ddots{\mathinner{\mkern1mu\raise6.5pt\hbox{.}\mkern2mu\raise3.6pt\hbox{.}\mkern2mu\raise0.7pt\hbox{.}\mkern1mu}}%
  \def\triangle{\mathord{\scalebox{0.9}{$\symup{\Delta}$}}}%
  \def\nabla{\mathord{\raisebox{\depth}{\scalebox{1}[-1]{$\symup{\Delta}$}}}}%
  \def\checkmark{\popcheck{popdark}}%
  \def\star{\mathbin{\ast}}\def\bigstar{\mathord{\ast}}%
  \def\diamond{\mathbin{\scalebox{0.6}{$\Diamond$}}}%
  \def\bigcup{\mathop{\mathchoice{\raisebox{-0.3em}{\scalebox{1.7}{$\cup$}}}{\scalebox{1.25}{$\cup$}}{\cup}{\cup}}\displaylimits}%
  \def\bigcap{\mathop{\mathchoice{\raisebox{-0.3em}{\scalebox{1.7}{$\cap$}}}{\scalebox{1.25}{$\cap$}}{\cap}{\cap}}\displaylimits}%
  \def\lesseqqgtr{\mathrel{\lessgtr}}\def\bigoplus{\mathop{\oplus}\displaylimits}%
}
\providecommand{\ssi}{si et seulement si} % abréviation fréquente
\AtBeginDocument{\providecommand{\wideparen}{\overparen}} % arc de cercle

%% FIGFIX-BEGIN v5 — correctifs globaux des figures (docs/onbuch-generation/tools/figfix)
\usepackage{adjustbox}\usepackage{etoolbox}\tcbuselibrary{raster}
% toute figure TikZ/pgfplots plus large que la ligne est réduite à la largeur disponible
\BeforeBeginEnvironment{tikzpicture}{\begin{adjustbox}{max width=\linewidth}}
\AfterEndEnvironment{tikzpicture}{\end{adjustbox}}
% légendes pgfplots lisibles par-dessus les courbes
\pgfplotsset{every axis/.append style={legend style={fill=white,fill opacity=0.92,text opacity=1,draw=black!15,inner sep=3pt}}}
% étiquettes d'arêtes (syntaxe edge["..."]) avec fond blanc
\tikzset{every edge quotes/.append style={fill=white,inner sep=1.2pt}}
%% FIGFIX-END
\begin{document}

\pagegarde{Barycentres et lignes de niveau}{Mathématiques}{1ère TI}{Configurations et transformations élémentaires du plan}{N12}{12 h}{Cette leçon introduit le barycentre de points pondérés, outil fondamental pour traiter les problèmes d'équilibre, d'alignement et de concours de droites. Elle se poursuit par l'étude des lignes de niveau, qui permettent de caractériser et construire des ensembles de points définis par des relations métriques.
\vspace{4pt}
\begin{multicols}{2}\small
\begin{itemize}
\item À la fin de cette leçon, je sais définir et déterminer le barycentre de deux points pondérés
\item À la fin de cette leçon, je sais construire géométriquement le barycentre de deux, trois ou quatre points pondérés
\item À la fin de cette leçon, je sais utiliser les barycentres partiels pour construire un barycentre de manière performante
\item À la fin de cette leçon, je sais démontrer que trois points sont alignés à l'aide des barycentres
\item À la fin de cette leçon, je sais démontrer que des droites sont concourantes à l'aide des barycentres
\item À la fin de cette leçon, je sais réduire la somme vectorielle Σ αᵢ MAᵢ selon que la somme des coefficients est nulle ou non
\end{itemize}\end{multicols}}

\begin{sommaire}
\begin{multicols}{2}
\begin{enumerate}
\item Barycentre de deux points pondérés
\item Barycentre de trois ou quatre points pondérés
\item Réduction de la somme vectorielle Σ αᵢ MAᵢ⃗
\item Lignes de niveau MA² + MB² = k et MA² − MB² = k
\item Lignes de niveau MA⃗ · MB⃗ = k et MA/MB = k
\item Activité d'intégration
\item Méthodes et exercices résolus
\item Exercices
\item Corrigés des exercices
\item Fiche bilan
\end{enumerate}
\end{multicols}
\end{sommaire}

\begin{objectifs}
\begin{itemize}
\item À la fin de cette leçon, je sais définir et déterminer le barycentre de deux points pondérés
\item À la fin de cette leçon, je sais construire géométriquement le barycentre de deux, trois ou quatre points pondérés
\item À la fin de cette leçon, je sais utiliser les barycentres partiels pour construire un barycentre de manière performante
\item À la fin de cette leçon, je sais démontrer que trois points sont alignés à l'aide des barycentres
\item À la fin de cette leçon, je sais démontrer que des droites sont concourantes à l'aide des barycentres
\item À la fin de cette leçon, je sais réduire la somme vectorielle Σ αᵢ MAᵢ selon que la somme des coefficients est nulle ou non
\item À la fin de cette leçon, je sais déterminer et caractériser les lignes de niveau MA² + MB² = k et MA² − MB² = k
\item À la fin de cette leçon, je sais déterminer et caractériser les lignes de niveau MA⃗ · MB⃗ = k et MA/MB = k
\item À la fin de cette leçon, je sais construire ces lignes de niveau et résoudre des problèmes concrets d'ensemble de points
\end{itemize}
\end{objectifs}

\begin{prerequis}
\begin{itemize}
\item Vecteurs du plan : colinéarité, relation de Chasles, opérations sur les vecteurs (Seconde)
\item Produit scalaire et ses propriétés (bilinearité, identités remarquables vectorielles)
\item Coordonnées dans un repère, milieu d'un segment, distance de deux points
\item Ensembles de points usuels : cercle, médiatrice, droite
\item Équations de cercles et de droites dans un repère orthonormé
\end{itemize}
\end{prerequis}

\newpage

\coursec{Barycentre de deux points pondérés}

Au marché Mokolo de Yaoundé, un commerçant équilibre sa balance romaine en déplaçant un contrepoids le long du fléau : où placer exactement le poids pour que le plateau reste horizontal ? De même, pour installer une antenne parabolique entre deux poteaux de masses différentes, le technicien cherche le point d'équilibre idéal. Ces problèmes d'équilibre et de « centre de gravité » se traduisent en mathématiques par la notion de \cle{barycentre}. Et lorsqu'un opérateur télécom cherche tous les points où la somme des carrés des distances à deux relais est constante, il trace sans le savoir une \cle{ligne de niveau}. Cette leçon donne les outils pour résoudre rigoureusement ces situations.

Dans cette première section, nous partons de l'intuition physique — une tige chargée à ses extrémités — pour construire la définition mathématique du barycentre de deux points, démontrer son existence et son unicité, établir la relation fondamentale $\overrightarrow{AG} = \dfrac{\beta}{\alpha+\beta}\overrightarrow{AB}$, étudier la position de $G$ selon les signes des coefficients, et donner les méthodes de construction.

\courssub{Situation d'introduction : la balance romaine}

Considérons une tige rigide $[AB]$, de masse négligeable, aux extrémités de laquelle sont suspendues deux masses : une masse $\alpha$ en $A$ et une masse $\beta$ en $B$. On cherche le point $G$ où placer le point d'appui (le crochet de suspension) pour que la tige reste horizontale, en équilibre.

\begin{popfigure}
\begin{tikzpicture}[scale=1.0,>=Stealth]
% Fléau gradué
\draw[line width=2.2pt,popdark] (0,0) -- (8,0);
\foreach \x in {0,0.5,...,8} \draw[popmuted] (\x,0.06) -- (\x,-0.06);
% Masses
\draw[popink,line width=1pt] (0.8,0) -- (0.8,-0.9);
\draw[fill=popblueL,draw=popblue] (0.8,-0.9) circle (0.38);
\node at (0.8,-0.9) {\small $\alpha$};
\draw[popink,line width=1pt] (7.2,0) -- (7.2,-0.9);
\draw[fill=poporangeL,draw=poporange] (7.2,-0.9) circle (0.38);
\node at (7.2,-0.9) {\small $\beta$};
% Point d'équilibre
\draw[fill=popink] (4.27,0) circle (0.09);
\draw[popink,line width=1pt] (4.27,0) -- (4.27,1.1);
\draw[popink,line width=1pt] (4.0,1.1) -- (4.54,1.1);
\node[popink,above] at (4.27,1.15) {$G$ : point d'équilibre};
% Labels
\node[below,popblue] at (0.8,0.15) {$A$};
\node[below,poporange] at (7.2,0.15) {$B$};
% Flèches de poids
\draw[->,popblue,line width=1.2pt] (0.8,-1.35) -- (0.8,-2.1);
\draw[->,poporange,line width=1.2pt] (7.2,-1.35) -- (7.2,-2.1);
\node[popmuted] at (4,-2.35) {\small Les poids tirent la tige vers le bas ; l'équilibre exige que $G$ soit plus proche de la masse la plus lourde.};
\end{tikzpicture}
\legende{Balance romaine : le point d'équilibre $G$ se rapproche de l'extrémité la plus chargée.}
\end{popfigure}

L'expérience quotidienne nous enseigne trois choses :
\begin{itemize}
\item si les deux masses sont égales, l'équilibre est atteint au \emph{milieu} de la tige ;
\item si la masse en $A$ est plus lourde, il faut déplacer le point d'appui \emph{vers $A$} ;
\item le point d'équilibre existe toujours et il est \emph{unique}.
\end{itemize}

La physique dit même quelque chose de très précis : la condition d'équilibre s'écrit $\alpha \cdot GA = \beta \cdot GB$ (égalité des moments). En tenant compte du sens des vecteurs, cette égalité devient une relation vectorielle remarquablement simple : $\alpha\overrightarrow{GA} + \beta\overrightarrow{GB} = \vec{0}$. C'est exactement cette relation que les mathématiques prennent comme définition.

\begin{savaistu}[D'où vient le mot « barycentre » ?]
Le mot \emph{barycentre} vient du grec \emph{barus} (lourd, pesant) et \emph{kentron} (centre) : c'est littéralement le « centre des pesanteurs ». Archimède (III\textsuperscript{e} siècle av. J.-C.) l'utilisait déjà, sans le nommer, dans ses lois du levier : « Donnez-moi un point d'appui et je soulèverai le monde. » Aujourd'hui, le barycentre est partout : centre de gravité d'un avion, positionnement d'antennes, moyennes pondérées de notes au Probatoire, et même calcul des orbites en astronomie.
\end{savaistu}

\courssub{Définition du barycentre de deux points pondérés}

\begin{definition}[Barycentre de deux points pondérés]
Soient $A$ et $B$ deux points du plan, et $\alpha$, $\beta$ deux réels tels que $\alpha + \beta \neq 0$. On appelle \cle{barycentre} des points pondérés $(A,\alpha)$ et $(B,\beta)$ l'unique point $G$ vérifiant :
\[ \alpha\overrightarrow{GA} + \beta\overrightarrow{GB} = \vec{0}. \]
On note $G = \mathrm{bar}\{(A,\alpha)\,;\,(B,\beta)\}$. Les réels $\alpha$ et $\beta$ sont appelés \cle{coefficients} (ou masses, ou poids) affectés aux points.
\end{definition}

\begin{attention}[La condition $\alpha+\beta \neq 0$ est indispensable]
Si $\alpha + \beta = 0$ avec $\alpha \neq 0$, la relation devient $\alpha(\overrightarrow{GA} - \overrightarrow{GB}) = \alpha\overrightarrow{BA} = \vec{0}$, ce qui est impossible car $A \neq B$. Il n'existe alors \textbf{aucun} point $G$ solution. Retenir : on ne peut parler de barycentre que lorsque la \cle{somme des coefficients est non nulle}.
\end{attention}

\courssub{Existence, unicité et relation fondamentale}

\begin{propriete}[Existence et unicité du barycentre — démonstration exigible]
Soient $(A,\alpha)$ et $(B,\beta)$ deux points pondérés avec $\alpha + \beta \neq 0$. Il existe un unique point $G$ tel que $\alpha\overrightarrow{GA} + \beta\overrightarrow{GB} = \vec{0}$, et ce point vérifie la \cle{relation fondamentale} :
\[ \overrightarrow{AG} = \frac{\beta}{\alpha+\beta}\,\overrightarrow{AB}. \]
\end{propriete}

\begin{experience}[Démonstration guidée, pas à pas]
Partons de la relation de définition et transformons-la grâce à la relation de Chasles, en faisant apparaître le point $A$ :
\begin{align*}
\alpha\overrightarrow{GA} + \beta\overrightarrow{GB} &= \vec{0} \\
\alpha\overrightarrow{GA} + \beta\left(\overrightarrow{GA} + \overrightarrow{AB}\right) &= \vec{0} & \text{(Chasles : } \overrightarrow{GB} = \overrightarrow{GA} + \overrightarrow{AB}\text{)} \\
(\alpha + \beta)\overrightarrow{GA} + \beta\overrightarrow{AB} &= \vec{0} & \text{(regroupement)} \\
(\alpha + \beta)\overrightarrow{GA} &= -\beta\overrightarrow{AB} \\
\overrightarrow{AG} &= \frac{\beta}{\alpha+\beta}\,\overrightarrow{AB}. & \text{(on multiplie par } -1 \text{, puis on divise par } \alpha+\beta \neq 0\text{)}
\end{align*}
La dernière ligne montre deux choses à la fois :
\begin{itemize}
\item \textbf{Existence} : le point $G$ défini par $\overrightarrow{AG} = \dfrac{\beta}{\alpha+\beta}\overrightarrow{AB}$ existe (on sait construire une fraction d'un vecteur) ;
\item \textbf{Unicité} : tout point vérifiant la définition vérifie nécessairement cette relation, qui ne définit qu'un seul point. Il n'y a donc qu'un seul barycentre. \hfill $\blacksquare$
\end{itemize}
\end{experience}

\begin{aretenir}[Conséquence immédiate : alignement]
La relation $\overrightarrow{AG} = \dfrac{\beta}{\alpha+\beta}\overrightarrow{AB}$ montre que les vecteurs $\overrightarrow{AG}$ et $\overrightarrow{AB}$ sont \emph{colinéaires}. Par conséquent :
\[ \text{Le barycentre } G \text{ de } (A,\alpha) \text{ et } (B,\beta) \text{ est toujours aligné avec } A \text{ et } B. \]
Autrement dit, $G$ appartient à la droite $(AB)$. C'est la première chose à vérifier dans toute construction.
\end{aretenir}

\begin{attention}[Erreur fréquente : inverser les coefficients]
Dans la formule $\overrightarrow{AG} = \dfrac{\beta}{\alpha+\beta}\overrightarrow{AB}$, c'est le coefficient de $B$ (c'est-à-dire $\beta$) qui apparaît au numérateur devant $\overrightarrow{AB}$, et non celui de $A$ ! Moyen mnémotechnique : \emph{« le point $G$ est attiré par la plus grosse masse »}. Si $\beta$ est grand, la fraction $\dfrac{\beta}{\alpha+\beta}$ est proche de $1$, donc $G$ est proche de $B$ : c'est cohérent, car $B$ est le point lourd. Vérifie toujours ta construction avec cette intuition physique.
\end{attention}

\courssub{Exemples chiffrés et position de $G$ selon les signes des coefficients}

\begin{exemplebox}[Trois configurations types]
Sur une droite graduée, plaçons $A$ d'abscisse $0$ et $B$ d'abscisse $10$.
\begin{enumerate}
\item $(A,1)\,;\,(B,1)$ : $\overrightarrow{AG} = \dfrac{1}{2}\overrightarrow{AB}$, donc $G$ a pour abscisse $5$ : c'est le milieu de $[AB]$.
\item $(A,2)\,;\,(B,3)$ : $\overrightarrow{AG} = \dfrac{3}{5}\overrightarrow{AB}$, donc $G$ a pour abscisse $6$. Remarque : $G$ est plus proche de $B$ (coefficient $3$) que de $A$ (coefficient $2$).
\item $(A,5)\,;\,(B,-2)$ : $\alpha+\beta = 3 \neq 0$ et $\overrightarrow{AG} = \dfrac{-2}{3}\overrightarrow{AB}$, donc $G$ a pour abscisse $-\dfrac{20}{3} \approx -6{,}67$ : $G$ est \emph{à l'extérieur} du segment, du côté de $A$.
\end{enumerate}
\end{exemplebox}

\begin{popfigure}
\begin{tikzpicture}[scale=0.85,>=Stealth]
% Cas 1 : alpha = beta
\draw[line width=1.6pt,popdark] (0,4) -- (10,4);
\draw[fill=popblue] (0,4) circle (0.09); \node[below] at (0,3.9) {$A$};
\draw[fill=poporange] (10,4) circle (0.09); \node[below] at (10,3.9) {$B$};
\draw[fill=popink] (5,4) circle (0.11); \node[above,popink] at (5,4.1) {$G$};
\node[left,align=right] at (-0.3,4) {\small $(A,1);(B,1)$\\ \small $G$ milieu};
% Cas 2 : alpha=2, beta=3
\draw[line width=1.6pt,popdark] (0,2) -- (10,2);
\draw[fill=popblue] (0,2) circle (0.09); \node[below] at (0,1.9) {$A$};
\draw[fill=poporange] (10,2) circle (0.09); \node[below] at (10,1.9) {$B$};
\draw[fill=popink] (6,2) circle (0.11); \node[above,popink] at (6,2.1) {$G$};
\node[left,align=right] at (-0.3,2) {\small $(A,2);(B,3)$\\ \small $G$ plus près de $B$};
% Cas 3 : signes contraires (corrigé : G en -6.67)
\draw[line width=1.6pt,popdark] (-7,0) -- (10,0);
\draw[fill=popblue] (0,0) circle (0.09); \node[below] at (0,-0.1) {$A$};
\draw[fill=poporange] (10,0) circle (0.09); \node[below] at (10,-0.1) {$B$};
\draw[fill=popink] (-6.67,0) circle (0.11); \node[above,popink] at (-6.67,0.1) {$G$};
\node[left,align=right] at (-7.3,0) {\small $(A,5);(B,-2)$\\ \small $G$ à l'extérieur};
\end{tikzpicture}
\legende{Position du barycentre $G$ selon les coefficients : milieu si $\alpha=\beta$, décalé vers la plus grosse masse, à l'extérieur si les signes sont contraires.}
\end{popfigure}

\begin{propriete}[Position de $G$ sur la droite $(AB)$]
Soit $G = \mathrm{bar}\{(A,\alpha)\,;\,(B,\beta)\}$ avec $\alpha+\beta \neq 0$.
\begin{itemize}
\item Si $\alpha\beta > 0$ (coefficients de \textbf{même signe}), alors $G$ appartient au segment $[AB]$.
\item Si $\alpha\beta < 0$ (coefficients de \textbf{signes contraires}), alors $G$ est à l'extérieur de $[AB]$, du côté du point dont le coefficient a la plus grande valeur absolue.
\item Dans tous les cas, $G$ est \cle{plus proche du point dont le coefficient a la plus grande valeur absolue}.
\end{itemize}
\end{propriete}

Justifions le premier point : si $\alpha$ et $\beta$ sont de même signe, le nombre $t = \dfrac{\beta}{\alpha+\beta}$ vérifie $0 \le t \le 1$ (par exemple si $\alpha,\beta > 0$, alors $0 < \beta < \alpha+\beta$). Or $\overrightarrow{AG} = t\overrightarrow{AB}$ avec $t \in [0,1]$ signifie exactement que $G \in [AB]$.

\courssub{Cas particulier : l'isobarycentre}

\begin{definition}[Isobarycentre]
Lorsque $\alpha = \beta$, le barycentre est appelé \cle{isobarycentre} de $A$ et $B$. La relation fondamentale donne $\overrightarrow{AG} = \dfrac{1}{2}\overrightarrow{AB}$ : l'isobarycentre de $A$ et $B$ est \textbf{le milieu du segment $[AB]$}.
\end{definition}

\courssub{Homogénéité des coefficients}

\begin{propriete}[Homogénéité — démonstration guidée]
Le barycentre d'un système de points pondérés ne change pas si l'on multiplie tous les coefficients par un même réel non nul $k$ :
\[ \mathrm{bar}\{(A,\alpha)\,;\,(B,\beta)\} = \mathrm{bar}\{(A,k\alpha)\,;\,(B,k\beta)\} \quad (k \neq 0). \]
\end{propriete}

\begin{experience}[Démonstration guidée]
Soit $G = \mathrm{bar}\{(A,\alpha)\,;\,(B,\beta)\}$. Par définition : $\alpha\overrightarrow{GA} + \beta\overrightarrow{GB} = \vec{0}$.
\begin{enumerate}
\item Multiplions les deux membres par $k \neq 0$ : $k\left(\alpha\overrightarrow{GA} + \beta\overrightarrow{GB}\right) = k\vec{0} = \vec{0}$.
\item Distribuons : $(k\alpha)\overrightarrow{GA} + (k\beta)\overrightarrow{GB} = \vec{0}$.
\item De plus, $k\alpha + k\beta = k(\alpha+\beta) \neq 0$ car $k \neq 0$ et $\alpha+\beta \neq 0$.
\item La relation obtenue dit exactement que $G = \mathrm{bar}\{(A,k\alpha)\,;\,(B,k\beta)\}$. \hfill $\blacksquare$
\end{enumerate}
\end{experience}

\begin{exemplebox}[Application pratique de l'homogénéité]
$\mathrm{bar}\{(A,2)\,;\,(B,3)\} = \mathrm{bar}\{(A,4)\,;\,(B,6)\} = \mathrm{bar}\{(A,-2)\,;\,(B,-3)\} = \mathrm{bar}\left\{\left(A,\tfrac{2}{5}\right)\,;\,\left(B,\tfrac{3}{5}\right)\right\}$. On choisit toujours les coefficients les plus simples pour les calculs : diviser par la somme pour obtenir des coefficients de somme $1$ est souvent très commode.
\end{exemplebox}

\courssub{Coordonnées du barycentre dans un repère}

\begin{propriete}[Coordonnées du barycentre]
Dans un repère $(O\,;\vec{i},\vec{j})$, si $A(x_A\,;y_A)$ et $B(x_B\,;y_B)$, alors $G = \mathrm{bar}\{(A,\alpha)\,;\,(B,\beta)\}$ a pour coordonnées :
\[ x_G = \frac{\alpha x_A + \beta x_B}{\alpha+\beta} \qquad \text{et} \qquad y_G = \frac{\alpha y_A + \beta y_B}{\alpha+\beta}. \]
\end{propriete}

Ces formules se déduisent immédiatement de la relation $\overrightarrow{AG} = \dfrac{\beta}{\alpha+\beta}\overrightarrow{AB}$ traduite en coordonnées. Remarque culturelle : une moyenne pondérée n'est rien d'autre qu'une abscisse de barycentre ! Ta moyenne générale, avec ses coefficients par matière, est le barycentre de tes notes.

\begin{exoresolu}[Calcul de coordonnées]
Énoncé. Dans un repère orthonormé, on donne $A(1\,;2)$ et $B(7\,;-4)$. Déterminer les coordonnées de $G = \mathrm{bar}\{(A,2)\,;\,(B,3)\}$.
\tcblower
Solution détaillée. La somme des coefficients est $\alpha+\beta = 2+3 = 5 \neq 0$ : le barycentre existe. Appliquons les formules :
\begin{align*}
x_G &= \frac{2 \times 1 + 3 \times 7}{5} = \frac{2+21}{5} = \frac{23}{5} = 4{,}6, \\
y_G &= \frac{2 \times 2 + 3 \times (-4)}{5} = \frac{4-12}{5} = -\frac{8}{5} = -1{,}6.
\end{align*}
Donc $G\left(\dfrac{23}{5}\,;\,-\dfrac{8}{5}\right)$. Vérification avec la relation fondamentale : $\overrightarrow{AB}(6\,;-6)$ et $\dfrac{3}{5}\overrightarrow{AB} = \left(\dfrac{18}{5}\,;-\dfrac{18}{5}\right)$, d'où $G = A + \dfrac{3}{5}\overrightarrow{AB} = \left(1+\dfrac{18}{5}\,;\,2-\dfrac{18}{5}\right) = \left(\dfrac{23}{5}\,;\,-\dfrac{8}{5}\right)$. Les deux méthodes concordent.
\end{exoresolu}

\courssub{Méthodes de construction du barycentre}

\begin{methode}[Construire $G = \mathrm{bar}\{(A,2)\,;\,(B,3)\}$]
\begin{enumerate}
\item \textbf{Écrire la relation fondamentale} : $\overrightarrow{AG} = \dfrac{\beta}{\alpha+\beta}\overrightarrow{AB} = \dfrac{3}{5}\overrightarrow{AB}$. (Attention : c'est le coefficient de $B$ au numérateur.)
\item \textbf{Graduer le segment} $[AB]$ en $5$ parts égales (à la règle, ou par le théorème de Thalès avec une demi-droite auxiliaire graduée en $5$).
\item \textbf{Reporter} : $G$ est le point tel que $AG = \dfrac{3}{5}AB$, c'est-à-dire la 3\textsuperscript{e} graduation en partant de $A$.
\item \textbf{Contrôler} : $G$ doit être entre $A$ et $B$ (coefficients de même signe) et plus proche de $B$ (car $3 > 2$). Si ce n'est pas le cas, il y a une erreur.
\end{enumerate}
\end{methode}

\begin{popfigure}
\begin{tikzpicture}[scale=1.0,>=Stealth]
\draw[line width=1.8pt,popdark] (0,0) -- (10,0);
\foreach \x in {0,2,4,6,8,10} \draw[popmuted] (\x,0.12) -- (\x,-0.12);
\draw[fill=popblue] (0,0) circle (0.1); \node[below=3pt,popblue] at (0,0) {$A\ (2)$};
\draw[fill=poporange] (10,0) circle (0.1); \node[below=3pt,poporange] at (10,0) {$B\ (3)$};
\draw[fill=popink] (6,0) circle (0.12); \node[above=3pt,popink] at (6,0) {$G$};
\draw[<->,popgreen,line width=1pt] (0,0.55) -- (6,0.55);
\node[popgreen,above] at (3,0.6) {\small $\frac{3}{5}AB$};
\draw[<->,poppurple,line width=1pt] (6,-0.55) -- (10,-0.55);
\node[poppurple,below] at (8,-0.6) {\small $\frac{2}{5}AB$};
\node[popmuted] at (5,-1.3) {\small Segment $[AB]$ partagé en $5$ parts égales ; $G$ est à $\frac{3}{5}$ de $A$.};
\end{tikzpicture}
\legende{Construction de $G = \mathrm{bar}\{(A,2)\,;\,(B,3)\}$ par report de la fraction $\frac{3}{5}$ de $\overrightarrow{AB}$.}
\end{popfigure}

\begin{methode}[Construire un barycentre à coefficients de signes contraires]
Pour $G = \mathrm{bar}\{(A,5)\,;\,(B,-2)\}$ :
\begin{enumerate}
\item Somme : $5 + (-2) = 3 \neq 0$. Relation : $\overrightarrow{AG} = \dfrac{-2}{3}\overrightarrow{AB}$.
\item Le coefficient de proportionnalité est \emph{négatif} : $G$ est donc de l'autre côté de $A$ par rapport à $B$.
\item On prolonge la droite $(AB)$ du côté de $A$ et on reporte $AG = \dfrac{2}{3}AB$ dans le sens opposé à $\overrightarrow{AB}$.
\item Contrôle : $G$ est du côté de $A$, le point de plus grand coefficient en valeur absolue ($|5| > |-2|$). Cohérent.
\end{enumerate}
\end{methode}

\begin{exoresolu}[Retrouver les coefficients à partir de la position]
Énoncé. Sur une droite, $G$ est le point du segment $[AB]$ tel que $AG = \dfrac{1}{4}AB$. Montrer que $G$ est le barycentre de $(A,\alpha)$ et $(B,\beta)$ pour des coefficients que l'on déterminera.
\tcblower
Solution détaillée. On sait que $\overrightarrow{AG} = \dfrac{1}{4}\overrightarrow{AB}$. D'après la relation fondamentale, $\overrightarrow{AG} = \dfrac{\beta}{\alpha+\beta}\overrightarrow{AB}$. Il faut donc avoir :
\[ \frac{\beta}{\alpha+\beta} = \frac{1}{4} \quad \Longleftrightarrow \quad 4\beta = \alpha + \beta \quad \Longleftrightarrow \quad \alpha = 3\beta. \]
On choisit par exemple $\beta = 1$, ce qui donne $\alpha = 3$. La somme des coefficients est $\alpha+\beta = 4 \neq 0$, le barycentre existe bien.
Donc $G = \mathrm{bar}\{(A,3)\,;\,(B,1)\}$.
Vérification : les coefficients sont de même signe ($3>0, 1>0$), donc $G \in [AB]$, et $\dfrac{\beta}{\alpha+\beta} = \dfrac{1}{4}$, ce qui correspond bien à $AG = \frac{1}{4}AB$.
\end{exoresolu}

\begin{attention}[Pièges de rédaction au Probatoire]
\begin{itemize}
\item Toujours \textbf{vérifier et écrire} que $\alpha + \beta \neq 0$ avant de conclure à l'existence du barycentre : c'est la première ligne attendue.
\item Ne pas confondre $\alpha\overrightarrow{GA} + \beta\overrightarrow{GB} = \vec{0}$ (relation de définition, centrée en $G$) avec $\overrightarrow{AG} = \dfrac{\beta}{\alpha+\beta}\overrightarrow{AB}$ (relation de construction, centrée en $A$).
\item Une égalité de \emph{longueurs} $3GA = GB$ ne suffit pas à caractériser un barycentre : il faut une égalité \emph{vectorielle}, car le sens compte.
\item Dans les coordonnées, le coefficient $\alpha$ multiplie les coordonnées de $A$ (pas d'inversion ici : l'inversion n'apparaît que dans la relation avec $\overrightarrow{AB}$).
\end{itemize}
\end{attention}

\begin{exercice}[Pour s'entraîner]
\begin{enumerate}
\item Construire $G = \mathrm{bar}\{(A,1)\,;\,(B,4)\}$ puis $H = \mathrm{bar}\{(A,3)\,;\,(B,-1)\}$ sur une même figure où $AB = 6$ cm. Que constate-t-on pour $H$ ?
\item Dans un repère, $A(-2\,;5)$ et $B(4\,;-1)$. Calculer les coordonnées de $K = \mathrm{bar}\{(A,2)\,;\,(B,1)\}$.
\item Déterminer deux réels $\alpha$ et $\beta$ tels que le milieu $I$ de $[AB]$ soit le barycentre de $(A,\alpha)$ et $(B,\beta)$ avec $\alpha + \beta = 10$.
\end{enumerate}
\end{exercice}

\begin{corrige}[Exercice]
\begin{enumerate}
\item $\overrightarrow{AG} = \dfrac{4}{5}\overrightarrow{AB}$ : $G$ est à $4{,}8$ cm de $A$. $\overrightarrow{AH} = \dfrac{-1}{2}\overrightarrow{AB}$ : $H$ est à l'extérieur de $[AB]$, à $3$ cm de $A$ du côté opposé à $B$ (coefficients de signes contraires, $|3| > |-1|$ donc $H$ du côté de $A$).
\item $x_K = \dfrac{2\times(-2) + 1\times 4}{3} = \dfrac{0}{3} = 0$ ; $y_K = \dfrac{2\times 5 + 1\times(-1)}{3} = \dfrac{9}{3} = 3$. Donc $K(0\,;3)$.
\item $I$ milieu de $[AB]$ est l'isobarycentre : il faut $\alpha = \beta$. Avec $\alpha + \beta = 10$, on obtient $\alpha = \beta = 5$.
\end{enumerate}
\end{corrige}

\begin{aretenir}[L'essentiel de la section]
\begin{itemize}
\item $G = \mathrm{bar}\{(A,\alpha)\,;\,(B,\beta)\}$ existe si et seulement si $\alpha + \beta \neq 0$, et vérifie $\alpha\overrightarrow{GA} + \beta\overrightarrow{GB} = \vec{0}$.
\item Relation fondamentale : $\overrightarrow{AG} = \dfrac{\beta}{\alpha+\beta}\overrightarrow{AB}$ ; en particulier $G \in (AB)$ toujours.
\item $G \in [AB]$ si $\alpha\beta > 0$ ; $G$ est plus proche du point de plus grand coefficient en valeur absolue.
\item Isobarycentre ($\alpha = \beta$) = milieu de $[AB]$.
\item Homogénéité : multiplier tous les coefficients par $k \neq 0$ ne change pas le barycentre.
\item Coordonnées : $x_G = \dfrac{\alpha x_A + \beta x_B}{\alpha+\beta}$, $y_G = \dfrac{\alpha y_A + \beta y_B}{\alpha+\beta}$.
\end{itemize}
\end{aretenir}

\coursec{Barycentre de trois ou quatre points pondérés}

Dans la section précédente, tu as appris à équilibrer une tige portant deux masses : le barycentre de deux points pondérés. Mais la réalité est rarement aussi simple. Un maçon de Yaoundé qui veut soulever une plaque triangulaire sans qu'elle bascule cherche le point d'équilibre de \emph{trois} points ; un ingénieur qui conçoit une table à quatre pieds raisonne sur \emph{quatre} points. Toute cette section répond à une question naturelle : ce que nous avons fait avec deux points peut-il se généraliser à trois, quatre, voire $n$ points ? La réponse est oui, et presque sans effort : la définition, l'existence, l'unicité et les formules de coordonnées s'étendent mot pour mot. La vraie nouveauté, c'est un outil redoutable : la propriété d'\cle{associativité} (ou des \cle{barycentres partiels}), qui permet de construire un barycentre de trois points en ne manipulant que des barycentres de deux points — que tu sais déjà construire.

\courssub{Définition, existence et unicité}

\begin{definition}[Barycentre de $n$ points pondérés]
Soient $A_1, A_2, \ldots, A_n$ des points du plan et $\alpha_1, \alpha_2, \ldots, \alpha_n$ des réels tels que
\[
\alpha_1 + \alpha_2 + \cdots + \alpha_n \neq 0.
\]
On appelle \cle{barycentre} du système $\{(A_1, \alpha_1), (A_2, \alpha_2), \ldots, (A_n, \alpha_n)\}$ l'unique point $G$ tel que
\[
\alpha_1 \overrightarrow{GA_1} + \alpha_2 \overrightarrow{GA_2} + \cdots + \alpha_n \overrightarrow{GA_n} = \vec{0}.
\]
On note $G = \mathrm{bar}\{(A_1, \alpha_1), \ldots, (A_n, \alpha_n)\}$.
\end{definition}

\begin{propriete}[Existence et unicité — admise]
Si $\alpha_1 + \alpha_2 + \cdots + \alpha_n \neq 0$, le barycentre $G$ existe et est unique. De plus, pour tout point $O$ du plan :
\[
\overrightarrow{OG} = \frac{\alpha_1 \overrightarrow{OA_1} + \alpha_2 \overrightarrow{OA_2} + \cdots + \alpha_n \overrightarrow{OA_n}}{\alpha_1 + \alpha_2 + \cdots + \alpha_n}.
\]
Ce résultat est la généralisation directe du cas de deux points démontré dans la section précédente ; on l'admet.
\end{propriete}

\begin{attention}[La condition sur la somme des coefficients]
La condition $\sum \alpha_i \neq 0$ est \textbf{indispensable}. Si la somme des coefficients est nulle, le barycentre n'existe pas (on verra en section 3 que la somme vectorielle $\sum \alpha_i \overrightarrow{MA_i}$ devient alors un vecteur constant non nul). Avant tout calcul, vérifie toujours cette condition : c'est la première ligne de rédaction attendue au Probatoire.
\end{attention}

Deux propriétés héritées du cas de deux points restent vraies, et tu dois les citer sans hésiter :

\begin{propriete}[Homogénéité et isobarycentre]
\begin{itemize}
\item \textbf{Homogénéité :} le barycentre d'un système ne change pas si l'on multiplie \emph{tous} les coefficients par un même réel non nul :
\[
\mathrm{bar}\{(A_i, \alpha_i)\} = \mathrm{bar}\{(A_i, \lambda \alpha_i)\} \quad (\lambda \neq 0).
\]
\item \textbf{Isobarycentre :} lorsque tous les coefficients sont égaux, le barycentre est appelé \cle{isobarycentre}. L'isobarycentre de deux points $A$ et $B$ est le milieu de $[AB]$ ; l'isobarycentre de trois points $A$, $B$, $C$ est le \cle{centre de gravité} du triangle $ABC$.
\end{itemize}
\end{propriete}

\courssub{Coordonnées du barycentre dans un repère}

\begin{propriete}[Coordonnées du barycentre]
Dans un repère $(O ; \vec{\imath}, \vec{\jmath})$, si $A_i(x_i \,; y_i)$ pour $i = 1, \ldots, n$ et si $\sum \alpha_i \neq 0$, alors le barycentre $G$ a pour coordonnées :
\[
x_G = \frac{\alpha_1 x_1 + \alpha_2 x_2 + \cdots + \alpha_n x_n}{\alpha_1 + \alpha_2 + \cdots + \alpha_n}
\qquad\text{et}\qquad
y_G = \frac{\alpha_1 y_1 + \alpha_2 y_2 + \cdots + \alpha_n y_n}{\alpha_1 + \alpha_2 + \cdots + \alpha_n}.
\]
\end{propriete}

Autrement dit : chaque coordonnée de $G$ est la \emph{moyenne pondérée} des coordonnées des points. C'est exactement le même mécanisme que ta moyenne trimestrielle au lycée : chaque note compte avec son coefficient.

\begin{exemplebox}[Calcul de coordonnées]
Dans un repère, on donne $A(1\,;2)$, $B(4\,;0)$, $C(-1\,;5)$. Déterminons les coordonnées de $G = \mathrm{bar}\{(A,1),(B,2),(C,3)\}$.

La somme des coefficients est $1+2+3 = 6 \neq 0$ : $G$ existe.
\begin{align*}
x_G &= \frac{1 \times 1 + 2 \times 4 + 3 \times (-1)}{6} = \frac{1 + 8 - 3}{6} = \frac{6}{6} = 1,\\[2pt]
y_G &= \frac{1 \times 2 + 2 \times 0 + 3 \times 5}{6} = \frac{2 + 0 + 15}{6} = \frac{17}{6}.
\end{align*}
Donc $G\left(1 \,; \dfrac{17}{6}\right)$.
\end{exemplebox}

\courssub{Associativité du barycentre : les barycentres partiels}

Voici le théorème central de cette section. L'intuition est physique : pour trouver le point d'équilibre de trois masses, on peut d'abord remplacer deux d'entre elles par une masse unique — la somme des deux — placée à leur point d'équilibre commun, puis chercher l'équilibre entre cette masse fictive et la troisième. Rien ne change.

\begin{propriete}[Associativité du barycentre — démonstration exigible]
Soit un système de points pondérés $(A_i, \alpha_i)$ de somme totale non nulle. On partage le système en deux sous-systèmes de sommes de coefficients non nulles. Si $G_1$ est le barycentre du premier sous-système, affecté de la somme $s_1$ de ses coefficients, et $G_2$ celui du second, affecté de la somme $s_2$, alors :
\[
G = \mathrm{bar}\{(G_1, s_1), (G_2, s_2)\}.
\]
On peut donc remplacer un sous-système par son \cle{barycentre partiel} affecté de la \textbf{somme} de ses coefficients.
\end{propriete}

\begin{experience}[Démonstration guidée de l'associativité (cas de trois points)]
Montrons que si $G' = \mathrm{bar}\{(B,\beta),(C,\gamma)\}$ avec $\beta + \gamma \neq 0$, et si $\alpha + \beta + \gamma \neq 0$, alors
\[
\mathrm{bar}\{(A,\alpha),(B,\beta),(C,\gamma)\} = \mathrm{bar}\{(A,\alpha),(G',\beta+\gamma)\}.
\]
\begin{enumerate}
\item \textbf{Traduis l'hypothèse sur $G'$.} Par définition du barycentre partiel :
\[
\beta \overrightarrow{G'B} + \gamma \overrightarrow{G'C} = \vec{0}.
\]
\item \textbf{Écris la fonction vectorielle de Leibniz.} Pour tout point $M$ du plan, posons
\[
\vec{f}(M) = \alpha \overrightarrow{MA} + \beta \overrightarrow{MB} + \gamma \overrightarrow{MC}.
\]
En introduisant le point $G'$ dans les deux derniers vecteurs (relation de Chasles) :
\[
\vec{f}(M) = \alpha \overrightarrow{MA} + \beta\left(\overrightarrow{MG'} + \overrightarrow{G'B}\right) + \gamma\left(\overrightarrow{MG'} + \overrightarrow{G'C}\right).
\]
\item \textbf{Regroupe et utilise l'étape 1.}
\[
\vec{f}(M) = \alpha \overrightarrow{MA} + (\beta + \gamma)\overrightarrow{MG'} + \underbrace{\left(\beta \overrightarrow{G'B} + \gamma \overrightarrow{G'C}\right)}_{=\,\vec{0}} = \alpha \overrightarrow{MA} + (\beta+\gamma)\overrightarrow{MG'}.
\]
\item \textbf{Conclus.} Le point $G$ cherché vérifie $\vec{f}(G) = \vec{0}$, c'est-à-dire
\[
\alpha \overrightarrow{GA} + (\beta+\gamma)\overrightarrow{GG'} = \vec{0},
\]
ce qui signifie exactement $G = \mathrm{bar}\{(A,\alpha),(G',\beta+\gamma)\}$. $\square$
\end{enumerate}
\end{experience}

\begin{attention}[L'erreur classique : le coefficient du barycentre partiel]
Quand tu remplaces le sous-système $\{(B,2),(C,3)\}$ par son barycentre partiel $G'$, tu dois affecter $G'$ de la \textbf{somme} $2+3 = 5$, et non d'un coefficient quelconque (ni $1$, ni le produit $6$ !). Écrire $G = \mathrm{bar}\{(A,1),(G',1)\}$ est l'erreur la plus fréquente au Probatoire : elle place $G$ au milieu de $[AG']$ alors qu'il doit être bien plus proche de $G'$. Retiens : \emph{le barycentre partiel hérite de la somme des coefficients de son sous-système.}
\end{attention}

\courssub{Construction performante d'un barycentre de trois points}

\begin{methode}[Construire $G = \mathrm{bar}\{(A,\alpha),(B,\beta),(C,\gamma)\}$]
\begin{enumerate}
\item Vérifier que $\alpha + \beta + \gamma \neq 0$.
\item \textbf{Choisir un regroupement} de deux points dont la somme des coefficients est non nulle (idéalement de même signe pour que le barycentre partiel soit \emph{entre} eux). Construire leur barycentre partiel, par exemple $G' = \mathrm{bar}\{(B,\beta),(C,\gamma)\}$, en utilisant $\overrightarrow{BG'} = \dfrac{\gamma}{\beta+\gamma}\overrightarrow{BC}$.
\item Affecter $G'$ du coefficient $\beta + \gamma$.
\item Construire $G = \mathrm{bar}\{(A,\alpha),(G',\beta+\gamma)\}$ sur la droite $(AG')$ grâce à $\overrightarrow{AG} = \dfrac{\beta+\gamma}{\alpha+\beta+\gamma}\overrightarrow{AG'}$.
\end{enumerate}
\end{methode}

\begin{exemplebox}[Construction de $G = \mathrm{bar}\{(A,1),(B,2),(C,3)\}$]
Somme totale : $1+2+3 = 6 \neq 0$, donc $G$ existe.

\textbf{Étape 1 : barycentre partiel.} Soit $G' = \mathrm{bar}\{(B,2),(C,3)\}$. La somme partielle $2+3=5 \neq 0$. Alors
\[
\overrightarrow{BG'} = \frac{3}{5}\overrightarrow{BC}.
\]
On place $G'$ sur $[BC]$, aux trois cinquièmes en partant de $B$.

\textbf{Étape 2 : regroupement.} Par associativité, $G = \mathrm{bar}\{(A,1),(G',5)\}$, donc
\[
\overrightarrow{AG} = \frac{5}{6}\overrightarrow{AG'}.
\]
On place $G$ sur $[AG']$, aux cinq sixièmes en partant de $A$. La construction est terminée — avec uniquement des barycentres de deux points.
\end{exemplebox}

\begin{popfigure}
\begin{tikzpicture}[scale=1.6,>=Stealth]
\coordinate (A) at (0,2.6);
\coordinate (B) at (-2,0);
\coordinate (C) at (2.4,0);
\coordinate (Gp) at ($(B)!0.6!(C)$);
\coordinate (G) at ($(A)!0.8333!(Gp)$);
\draw[popline,thick] (A)--(B)--(C)--cycle;
\draw[popmuted,dashed] (A)--(Gp);
\fill[popblue] (A) circle (1.4pt) node[above left,popdark]{$A\ (1)$};
\fill[popblue] (B) circle (1.4pt) node[below left,popdark]{$B\ (2)$};
\fill[popblue] (C) circle (1.4pt) node[below right,popdark]{$C\ (3)$};
\fill[poporange] (Gp) circle (1.6pt) node[below,popdark]{$G'\ (5)$};
\fill[popink] (G) circle (1.8pt) node[right,popdark]{$G$};
\draw[->,poporange,thick] (B) -- ($(B)!0.6!(C)$) node[midway,below left,popdark,font=\small]{$\frac{3}{5}$};
\draw[->,popgreen,thick] (A) -- ($(A)!0.8333!(Gp)$) node[midway,above right,popdark,font=\small]{$\frac{5}{6}$};
\end{tikzpicture}
\legende{Construction de $G = \mathrm{bar}\{(A,1),(B,2),(C,3)\}$ : on construit d'abord le barycentre partiel $G' = \mathrm{bar}\{(B,2),(C,3)\}$ avec $\overrightarrow{BG'} = \frac{3}{5}\overrightarrow{BC}$, puis $G$ sur $[AG']$ avec $\overrightarrow{AG} = \frac{5}{6}\overrightarrow{AG'}$.}
\end{popfigure}

\courssub{Isobarycentre de trois points : le centre de gravité du triangle}

\begin{definition}[Centre de gravité]
L'isobarycentre de trois points $A$, $B$, $C$ (coefficients tous égaux à $1$) est appelé \cle{centre de gravité} du triangle $ABC$. C'est le point $G$ tel que
\[
\overrightarrow{GA} + \overrightarrow{GB} + \overrightarrow{GC} = \vec{0}.
\]
\end{definition}

\begin{propriete}[Le centre de gravité et les médianes]
Soit $ABC$ un triangle, $A'$ le milieu de $[BC]$, $B'$ le milieu de $[CA]$, $C'$ le milieu de $[AB]$. Le centre de gravité $G$ appartient à chacune des trois médianes et se situe aux deux tiers de chacune en partant du sommet :
\[
\overrightarrow{AG} = \frac{2}{3}\overrightarrow{AA'}, \qquad
\overrightarrow{BG} = \frac{2}{3}\overrightarrow{BB'}, \qquad
\overrightarrow{CG} = \frac{2}{3}\overrightarrow{CC'}.
\]
En particulier, \textbf{les trois médianes d'un triangle sont concourantes en $G$}.
\end{propriete}

\begin{exemplebox}[Démonstration rédigée : concours des médianes]
\textbf{Énoncé.} Démontrer que les médianes d'un triangle $ABC$ sont concourantes.

\textbf{Rédaction.} Soit $G$ l'isobarycentre de $A$, $B$, $C$ : $G = \mathrm{bar}\{(A,1),(B,1),(C,1)\}$ (somme $3 \neq 0$).

Notons $A'$ le milieu de $[BC]$. Alors $A' = \mathrm{bar}\{(B,1),(C,1)\}$. Par associativité du barycentre :
\[
G = \mathrm{bar}\{(A,1),(A',2)\}.
\]
Donc $G$ appartient à la droite $(AA')$, c'est-à-dire à la médiane issue de $A$, et $\overrightarrow{AG} = \dfrac{2}{3}\overrightarrow{AA'}$.

Le même raisonnement avec $B'$ milieu de $[CA]$ donne $G = \mathrm{bar}\{(B,1),(B',2)\}$, donc $G \in (BB')$ ; et avec $C'$ milieu de $[AB]$, $G = \mathrm{bar}\{(C,1),(C',2)\}$, donc $G \in (CC')$.

Le point $G$ appartient aux trois médianes : elles sont donc \textbf{concourantes en $G$}, situé aux deux tiers de chacune à partir du sommet. $\square$
\end{exemplebox}

\begin{popfigure}
\begin{tikzpicture}[scale=1.5,>=Stealth]
\coordinate (A) at (0,2.8);
\coordinate (B) at (-2.2,0);
\coordinate (C) at (2.4,0);
\coordinate (Ap) at ($(B)!0.5!(C)$);
\coordinate (Bp) at ($(A)!0.5!(C)$);
\coordinate (Cp) at ($(A)!0.5!(B)$);
\coordinate (G) at ($(A)!0.6667!(Ap)$);
\draw[popline,thick] (A)--(B)--(C)--cycle;
\draw[popblue,thick] (A)--(Ap);
\draw[poporange,thick] (B)--(Bp);
\draw[popgreen,thick] (C)--(Cp);
\fill[popdark] (A) circle (1.3pt) node[above,popdark]{$A$};
\fill[popdark] (B) circle (1.3pt) node[below left,popdark]{$B$};
\fill[popdark] (C) circle (1.3pt) node[below right,popdark]{$C$};
\fill[popblue] (Ap) circle (1.2pt) node[below,popdark]{$A'$};
\fill[poporange] (Bp) circle (1.2pt) node[above right,popdark]{$B'$};
\fill[popgreen] (Cp) circle (1.2pt) node[above left,popdark]{$C'$};
\fill[popink] (G) circle (1.8pt) node[below right,popdark]{$G$};
\end{tikzpicture}
\legende{Les trois médianes $(AA')$, $(BB')$, $(CC')$ du triangle $ABC$ sont concourantes au centre de gravité $G$, isobarycentre de $A$, $B$, $C$, situé aux deux tiers de chaque médiane à partir du sommet.}
\end{popfigure}

\begin{savaistu}[Le centre de gravité dans la vie réelle]
Le centre de gravité d'une plaque triangulaire homogène (en bois, en tôle) est exactement l'isobarycentre de ses trois sommets. C'est le point où l'on peut poser la plaque sur la pointe d'un doigt sans qu'elle bascule. Les menuisiers de Douala l'utilisent intuitivement pour équilibrer les plateaux de table, et les ingénieurs du Génie civil en tiennent compte pour stabiliser les structures triangulaires (fermes de toit, pylônes). Archimède, il y a plus de deux mille ans, avait déjà fondé toute sa théorie du levier sur cette idée de barycentre.
\end{savaistu}

\courssub{Applications : alignement de points et concours de droites}

L'associativité fournit une stratégie élégante pour deux types de problèmes classiques au Probatoire.

\begin{propriete}[Barycentre et alignement]
Si un même point $G$ est barycentre de $(A,\alpha),(B,\beta)$ et aussi de $(C,\gamma),(D,\delta)$, alors $G$ appartient à la fois à $(AB)$ et à $(CD)$ : les droites $(AB)$ et $(CD)$ passent par $G$. Réciproquement, tout barycentre de deux points est \textbf{aligné} avec ces deux points.
\end{propriete}

\begin{methode}[Montrer que trois points sont alignés]
Pour montrer que $G$, $M$, $N$ sont alignés, on cherche à écrire $G$ comme barycentre de $M$ et $N$ (avec des coefficients de somme non nulle), ou on montre que $M$ est barycentre de deux points situés sur une même droite passant par $G$ et $N$. L'outil clé : un barycentre de deux points est toujours sur la droite qu'ils définissent.
\end{methode}

\begin{methode}[Montrer que des droites sont concourantes]
Pour montrer que trois droites $(d_1)$, $(d_2)$, $(d_3)$ sont concourantes :
\begin{enumerate}
\item Introduire un point $G$ défini comme barycentre d'un système bien choisi (souvent l'isobarycentre ou un barycentre lié à l'énoncé).
\item Par associativité, écrire $G$ de \emph{plusieurs façons} comme barycentre de deux points, chaque écriture le plaçant sur l'une des droites.
\item Conclure : $G$ appartient à toutes les droites, donc elles sont concourantes en $G$.
\end{enumerate}
C'est exactement la méthode utilisée ci-dessus pour les médianes.
\end{methode}

\begin{exoresolu}[Alignement, concours et barycentres partiels]
\textbf{Énoncé.} Soit $ABC$ un triangle. On définit $I = \mathrm{bar}\{(A,1),(B,2)\}$ et $J = \mathrm{bar}\{(A,1),(C,2)\}$. Soit $G = \mathrm{bar}\{(A,1),(B,2),(C,2)\}$.
Montrer que les droites $(AA')$, $(IC)$ et $(JB)$ sont concourantes en $G$, où $A'$ est le milieu de $[BC]$. Préciser la position de $G$ sur la médiane $[AA']$.

\tcblower
\textbf{Solution détaillée.}

La somme des coefficients de $G$ est $1+2+2 = 5 \neq 0$ : $G$ existe.

\textbf{Premier regroupement (sur la base $BC$).} Posons $A' = \mathrm{bar}\{(B,2),(C,2)\}$. Par homogénéité (division par $2$), $A' = \mathrm{bar}\{(B,1),(C,1)\}$, donc $A'$ est le \textbf{milieu de $[BC]$}. La somme partielle est $4$. Par associativité :
\[
G = \mathrm{bar}\{(A,1),(A',4)\}.
\]
Donc $G \in (AA')$, la médiane issue de $A$, et $\overrightarrow{AG} = \dfrac{4}{5}\overrightarrow{AA'}$.

\textbf{Deuxième regroupement (en isolant $C$).} $I = \mathrm{bar}\{(A,1),(B,2)\}$, affecté de la somme $3$. Par associativité :
\[
G = \mathrm{bar}\{(I,3),(C,2)\}.
\]
Donc $G \in (IC)$.

\textbf{Troisième regroupement (en isolant $B$).} De même, $J = \mathrm{bar}\{(A,1),(C,2)\}$ (somme $3$) donne $G = \mathrm{bar}\{(J,3),(B,2)\}$, donc $G \in (JB)$.

\textbf{Conclusion.} Le point $G$ appartient simultanément aux trois droites $(AA')$, $(IC)$ et $(JB)$ : elles sont \textbf{concourantes en $G$}. De plus, $G$ se situe aux $\dfrac{4}{5}$ de la médiane $[AA']$ en partant du sommet $A$.
\end{exoresolu}

\begin{exercice}[Pour t'entraîner]
Soit $ABCD$ un quadrilatère. On pose $G = \mathrm{bar}\{(A,1),(B,1),(C,1),(D,1)\}$ (isobarycentre des quatre sommets).
\begin{enumerate}
\item En regroupant $(A,1),(B,1)$ d'une part et $(C,1),(D,1)$ d'autre part, montrer que $G$ est le milieu du segment joignant les milieux de $[AB]$ et $[CD]$.
\item En déduire que les trois segments joignant les milieux des côtés opposés et les milieux des diagonales sont concourants en $G$.
\end{enumerate}
\end{exercice}

\begin{corrige}[Exercice 1]
\begin{enumerate}
\item Soit $M$ le milieu de $[AB]$ et $N$ le milieu de $[CD]$. Alors $M = \mathrm{bar}\{(A,1),(B,1)\}$ (somme $2$) et $N = \mathrm{bar}\{(C,1),(D,1)\}$ (somme $2$). Par associativité :
\[
G = \mathrm{bar}\{(M,2),(N,2)\} = \mathrm{bar}\{(M,1),(N,1)\}
\]
par homogénéité. Donc $G$ est le milieu de $[MN]$.
\item Soit $P$ le milieu de $[AD]$ et $Q$ le milieu de $[BC]$ : le même regroupement donne $G$ milieu de $[PQ]$. Soit $R$ le milieu de $[AC]$ et $S$ le milieu de $[BD]$ (milieux des diagonales) : $G$ est aussi milieu de $[RS]$. Les trois segments $[MN]$, $[PQ]$, $[RS]$ ont le même milieu $G$ : ils sont donc concourants en $G$.
\end{enumerate}
\end{corrige}

\begin{aretenir}[L'essentiel de la section]
\begin{itemize}
\item $G = \mathrm{bar}\{(A_i, \alpha_i)\}$ (avec $\sum \alpha_i \neq 0$) signifie $\sum \alpha_i \overrightarrow{GA_i} = \vec{0}$ ; existence et unicité garanties.
\item Coordonnées : moyennes pondérées des coordonnées des points.
\item \textbf{Associativité} : on peut remplacer un sous-système par son barycentre partiel affecté de la \textbf{somme} de ses coefficients — c'est l'outil de construction et de démonstration numéro un.
\item L'isobarycentre de trois points est le centre de gravité du triangle, situé aux $\dfrac{2}{3}$ de chaque médiane ; les médianes sont concourantes en ce point.
\item Double expression d'un même barycentre $\Rightarrow$ alignements et droites concourantes.
\end{itemize}
\end{aretenir}

\coursec{Réduction de la somme vectorielle $\sum \alpha_i \overrightarrow{MA_i}$}

Dans la section précédente, tu as appris à déterminer et à construire le barycentre de trois ou quatre points pondérés, notamment grâce à la technique des barycentres partiels. Nous allons maintenant exploiter le barycentre comme un véritable \emph{outil de calcul vectoriel} : il va nous permettre de \cle{réduire} des sommes de vecteurs qui, à première vue, semblent dépendre de façon compliquée d'un point $M$ variable du plan. C'est cette réduction qui ouvrira la porte, dans les deux prochaines sections, à l'étude des lignes de niveau.

\courssub{Le problème posé}

Considérons $n$ points pondérés $(A_1, \alpha_1), (A_2, \alpha_2), \ldots, (A_n, \alpha_n)$ du plan. Pour tout point $M$ du plan, on peut former le vecteur
\[
\vec{f}(M) = \sum_{i=1}^{n} \alpha_i \overrightarrow{MA_i} = \alpha_1 \overrightarrow{MA_1} + \alpha_2 \overrightarrow{MA_2} + \cdots + \alpha_n \overrightarrow{MA_n}.
\]

Lorsque $M$ se déplace dans le plan, chacun des vecteurs $\overrightarrow{MA_i}$ change : la somme $\vec{f}(M)$ semble donc être un objet très compliqué, qui dépend de $M$ de manière inextricable.

\begin{savaistu}[Une question naturelle]
En physique, lorsqu'on étudie un système de masses placées en $A_1, \ldots, A_n$ (par exemple les piliers d'un pont sur le Wouri), on a constamment besoin de sommes du type $\sum \alpha_i \overrightarrow{MA_i}$ pour calculer des moments de forces. Les physiciens ont découvert que ces sommes se simplifient miraculeusement lorsqu'on introduit le centre de masse $G$. C'est exactement le théorème que nous allons démontrer.
\end{savaistu}

La question centrale de cette section est donc : \textbf{peut-on écrire $\vec{f}(M)$ sous une forme simple, qui fasse apparaître explicitement la dépendance en $M$ ?} La réponse est oui, et elle dépend d'un seul nombre : la somme $S = \alpha_1 + \alpha_2 + \cdots + \alpha_n$ des coefficients.

\courssub{Le théorème de réduction}

\begin{propriete}[Théorème de réduction de $\sum \alpha_i \overrightarrow{MA_i}$ — démonstration exigible]
Soient $(A_1, \alpha_1), \ldots, (A_n, \alpha_n)$ des points pondérés et $S = \sum_{i=1}^{n} \alpha_i$.

\textbf{Cas 1 : $S \neq 0$.} Le système admet un barycentre $G$, et pour tout point $M$ du plan :
\[
\boxed{\;\sum_{i=1}^{n} \alpha_i \overrightarrow{MA_i} = S \, \overrightarrow{MG} = \left(\sum_{i=1}^{n} \alpha_i\right) \overrightarrow{MG}\;}
\]

\textbf{Cas 2 : $S = 0$.} Le vecteur $\sum_{i=1}^{n} \alpha_i \overrightarrow{MA_i}$ est \cle{constant} : il ne dépend pas du point $M$. Plus précisément, pour tous points $M$ et $N$ du plan :
\[
\sum_{i=1}^{n} \alpha_i \overrightarrow{MA_i} = \sum_{i=1}^{n} \alpha_i \overrightarrow{NA_i} = \sum_{i=1}^{n} \alpha_i \overrightarrow{\Omega A_i}
\]
où $\Omega$ est un point quelconque fixé une fois pour toutes.
\end{propriete}

\begin{experience}[Démonstration complète du théorème]
\textbf{Cas 1 : supposons $S \neq 0$.} Le barycentre $G$ du système existe. L'idée est d'introduire $G$ dans chaque vecteur $\overrightarrow{MA_i}$ grâce à la \cle{relation de Chasles} : pour tout indice $i$,
\[
\overrightarrow{MA_i} = \overrightarrow{MG} + \overrightarrow{GA_i}.
\]
En multipliant par $\alpha_i$ puis en sommant :
\begin{align*}
\sum_{i=1}^{n} \alpha_i \overrightarrow{MA_i}
&= \sum_{i=1}^{n} \alpha_i \left(\overrightarrow{MG} + \overrightarrow{GA_i}\right) \\
&= \sum_{i=1}^{n} \alpha_i \overrightarrow{MG} + \sum_{i=1}^{n} \alpha_i \overrightarrow{GA_i} \\
&= \left(\sum_{i=1}^{n} \alpha_i\right) \overrightarrow{MG} + \underbrace{\sum_{i=1}^{n} \alpha_i \overrightarrow{GA_i}}_{=\,\vec{0}} \\
&= S \, \overrightarrow{MG}.
\end{align*}
La somme $\sum \alpha_i \overrightarrow{GA_i}$ est nulle par la \textbf{caractérisation vectorielle du barycentre} : $G = \mathrm{bar}\{(A_i, \alpha_i)\}$ signifie précisément que $\sum \alpha_i \overrightarrow{GA_i} = \vec{0}$. $\blacksquare$

\medskip
\textbf{Cas 2 : supposons $S = 0$.} Soient $M$ et $N$ deux points quelconques du plan. Introduisons $N$ par la relation de Chasles dans chaque vecteur :
\begin{align*}
\sum_{i=1}^{n} \alpha_i \overrightarrow{MA_i}
&= \sum_{i=1}^{n} \alpha_i \left(\overrightarrow{MN} + \overrightarrow{NA_i}\right) \\
&= \left(\sum_{i=1}^{n} \alpha_i\right) \overrightarrow{MN} + \sum_{i=1}^{n} \alpha_i \overrightarrow{NA_i} \\
&= 0 \cdot \overrightarrow{MN} + \sum_{i=1}^{n} \alpha_i \overrightarrow{NA_i} \\
&= \sum_{i=1}^{n} \alpha_i \overrightarrow{NA_i}.
\end{align*}
Le vecteur obtenu est le même pour $M$ et pour $N$ : il est donc \textbf{indépendant de $M$}, c'est-à-dire constant. Pour le calculer, on peut remplacer $M$ par n'importe quel point pratique, par exemple $A_1$ : le vecteur constant vaut $\sum_{i=1}^{n} \alpha_i \overrightarrow{A_1 A_i}$. $\blacksquare$
\end{experience}

\begin{methode}[Comment réduire une somme $\sum \alpha_i \overrightarrow{MA_i}$ en pratique]
\begin{enumerate}
\item \textbf{Calculer la somme des coefficients} $S = \sum \alpha_i$.
\item \textbf{Si $S \neq 0$} : nommer $G$ le barycentre du système (le définir proprement par une égalité vectorielle), puis écrire directement $\sum \alpha_i \overrightarrow{MA_i} = S\,\overrightarrow{MG}$.
\item \textbf{Si $S = 0$} : le vecteur est constant. Pour le calculer, remplacer $M$ par l'un des points du système (souvent $A_1$), ce qui fait disparaître un terme, puis réduire avec Chasles.
\item \textbf{Conclure} en interprétant géométriquement le résultat (colinéarité, norme, lieu de points...).
\end{enumerate}
\end{methode}

\courssub{Conséquence : la caractérisation vectorielle du barycentre}

Appliquons le Cas 1 avec $M = G$. On obtient $S\,\overrightarrow{GG} = \vec{0}$, c'est-à-dire :

\begin{aretenir}[Caractérisation vectorielle du barycentre]
Si $S = \sum \alpha_i \neq 0$, alors
\[
G = \mathrm{bar}\{(A_1, \alpha_1), \ldots, (A_n, \alpha_n)\} \iff \sum_{i=1}^{n} \alpha_i \overrightarrow{GA_i} = \vec{0}.
\]
C'est cette égalité qui a servi de pivot dans la démonstration du théorème de réduction, et c'est elle qu'il faut écrire en premier dans toute copie du Probatoire lorsqu'un barycentre est introduit.
\end{aretenir}

\courssub{Exemples rédigés pas à pas}

\begin{exemplebox}[Cas d'une somme de coefficients non nulle]
Soient $A$, $B$, $C$ trois points du plan. Réduisons, pour tout point $M$, le vecteur
\[
\vec{u}(M) = \overrightarrow{MA} + 2\overrightarrow{MB} + \overrightarrow{MC}.
\]
\textbf{Étape 1.} Somme des coefficients : $S = 1 + 2 + 1 = 4 \neq 0$.

\textbf{Étape 2.} Soit $G = \mathrm{bar}\{(A,1), (B,2), (C,1)\}$, c'est-à-dire le point défini par $\overrightarrow{GA} + 2\overrightarrow{GB} + \overrightarrow{GC} = \vec{0}$.

\textbf{Étape 3.} Par le théorème de réduction :
\[
\vec{u}(M) = 4\,\overrightarrow{MG} \quad \text{pour tout point } M.
\]
Les trois vecteurs $\overrightarrow{MA}$, $2\overrightarrow{MB}$, $\overrightarrow{MC}$ se comportent donc, collectivement, comme l'unique vecteur $4\overrightarrow{MG}$.
\end{exemplebox}

\begin{exemplebox}[Cas d'une somme de coefficients nulle : vecteur constant]
Soient $A$, $B$, $C$ trois points. Réduisons
\[
\vec{v}(M) = \overrightarrow{MA} + 2\overrightarrow{MB} - 3\overrightarrow{MC}.
\]
\textbf{Étape 1.} Somme des coefficients : $S = 1 + 2 - 3 = 0$. Le vecteur $\vec{v}(M)$ est donc \textbf{constant}.

\textbf{Étape 2.} Pour le calculer, remplaçons $M$ par $A$ (le terme $\overrightarrow{MA}$ s'annule) :
\begin{align*}
\vec{v}(M) &= \vec{v}(A) = \vec{0} + 2\overrightarrow{AB} - 3\overrightarrow{AC} \\
&= 2\overrightarrow{AB} - 3\overrightarrow{AC}.
\end{align*}
\textbf{Conclusion :} pour tout point $M$ du plan, $\overrightarrow{MA} + 2\overrightarrow{MB} - 3\overrightarrow{MC} = 2\overrightarrow{AB} - 3\overrightarrow{AC}$, un vecteur \emph{fixe} qui ne dépend que de la configuration $(A, B, C)$.
\end{exemplebox}

\begin{attention}[Le piège classique du Probatoire]
Quand la somme des coefficients est nulle, \textbf{ne conclus jamais} que le vecteur $\sum \alpha_i \overrightarrow{MA_i}$ est \emph{nul} ! Le théorème dit qu'il est \emph{constant}, c'est-à-dire indépendant de $M$. Il n'est nul que dans un cas très particulier : lorsque le calcul en un point donne $\vec{0}$ (par exemple si $\sum \alpha_i \overrightarrow{A_1 A_i} = \vec{0}$). Rédiger « $S = 0$ donc la somme est nulle » est une erreur qui coûte des points à chaque session.
\end{attention}

\courssub{Illustration : trois vecteurs remplacés par un seul}

\begin{popfigure}
\begin{center}
\begin{tikzpicture}[scale=1.1, >=Stealth]
\coordinate (A) at (0,0);
\coordinate (B) at (5,0.6);
\coordinate (C) at (1.6,3.2);
\coordinate (G) at (2.2,1.3);
\coordinate (M) at (4.6,3.4);
% points
\fill[popblue] (A) circle (2pt) node[below left, popdark] {$A$};
\fill[popblue] (B) circle (2pt) node[below right, popdark] {$B$};
\fill[popblue] (C) circle (2pt) node[above, popdark] {$C$};
\fill[poporange] (G) circle (2.5pt) node[below, popdark] {$G$};
\fill[poppurple] (M) circle (2pt) node[above right, popdark] {$M$};
% vecteurs MA, MB, MC (pointillés)
\draw[->, popmuted, dashed, thick] (M) -- (A) node[midway, above left, popmuted] {$\overrightarrow{MA}$};
\draw[->, popmuted, dashed, thick] (M) -- (B) node[midway, below right, popmuted] {$\overrightarrow{MB}$};
\draw[->, popmuted, dashed, thick] (M) -- (C) node[midway, above, popmuted] {$\overrightarrow{MC}$};
% vecteur MG et 3MG
\draw[->, poporange, very thick] (M) -- (G) node[midway, right, popdark] {$\overrightarrow{MG}$};
\coordinate (G3) at ($(M)!3!(G)$);
\draw[->, popgreen, very thick] (M) -- (G3) node[midway, below right, popdark] {$3\overrightarrow{MG}$};
\fill[popgreen] (G3) circle (2pt) node[below left, popdark] {$G'$};
% triangle de référence
\draw[popline, thin] (A) -- (B) -- (C) -- cycle;
\end{tikzpicture}
\end{center}
\legende{Pour le système $\{(A,1),(B,1),(C,1)\}$ de barycentre $G$ (isobarycentre) : la somme $\overrightarrow{MA} + \overrightarrow{MB} + \overrightarrow{MC}$ (vecteurs en pointillés) se réduit à l'unique vecteur $3\overrightarrow{MG}$, représenté en vert. Le point $G'$ est tel que $\overrightarrow{MG'} = 3\overrightarrow{MG}$.}
\end{popfigure}

\courssub{Applications : ensembles de points et lien avec les lignes de niveau}

La réduction transforme une condition compliquée portant sur $M$ en une condition simple portant sur la distance ou la direction de $\overrightarrow{MG}$. C'est exactement la stratégie que nous utiliserons dans les sections 4 et 5.

\begin{exoresolu}[Ensemble des points tels qu'une somme vectorielle a une norme fixée]
\textbf{Énoncé.} Soient $A$ et $B$ deux points distincts tels que $AB = 6$ cm. Déterminer l'ensemble $(\mathcal{E})$ des points $M$ du plan tels que
\[
\left\| 2\overrightarrow{MA} + \overrightarrow{MB} \right\| = 9 \text{ cm}.
\]
\tcblower
\textbf{Solution.} \textbf{Étape 1 : réduction.} La somme des coefficients est $2 + 1 = 3 \neq 0$. Soit $G = \mathrm{bar}\{(A,2), (B,1)\}$, défini par $2\overrightarrow{GA} + \overrightarrow{GB} = \vec{0}$. Par le théorème de réduction, pour tout point $M$ :
\[
2\overrightarrow{MA} + \overrightarrow{MB} = 3\overrightarrow{MG}.
\]
\textbf{Étape 2 : traduction de la condition.}
\[
\left\| 2\overrightarrow{MA} + \overrightarrow{MB} \right\| = 9 \iff \left\| 3\overrightarrow{MG} \right\| = 9 \iff 3\,MG = 9 \iff MG = 3.
\]
\textbf{Étape 3 : conclusion.} L'ensemble $(\mathcal{E})$ est le \textbf{cercle de centre $G$ et de rayon $3$ cm}. Au passage, plaçons $G$ : de $2\overrightarrow{GA} + \overrightarrow{GB} = \vec{0}$ on tire $\overrightarrow{AG} = \dfrac{1}{3}\overrightarrow{AB}$, donc $G$ est le point de $[AB]$ situé à $2$ cm de $A$.
\end{exoresolu}

\begin{exoresolu}[Ensemble des points tels qu'une somme vectorielle est colinéaire à un vecteur donné]
\textbf{Énoncé.} Soient $A$, $B$, $C$ trois points non alignés. Déterminer l'ensemble $(\mathcal{D})$ des points $M$ tels que le vecteur $\overrightarrow{MA} + \overrightarrow{MB} - 2\overrightarrow{MC}$ soit colinéaire à $\overrightarrow{AB}$.
\tcblower
\textbf{Solution.} \textbf{Étape 1 : réduction.} Somme des coefficients : $1 + 1 - 2 = 0$. Le vecteur est donc constant. Calculons-le en remplaçant $M$ par $C$ :
\[
\overrightarrow{MA} + \overrightarrow{MB} - 2\overrightarrow{MC} = \overrightarrow{CA} + \overrightarrow{CB} - \vec{0} = \overrightarrow{CA} + \overrightarrow{CB}.
\]
\textbf{Étape 2 : discussion.} Le vecteur constant $\vec{w} = \overrightarrow{CA} + \overrightarrow{CB}$ est non nul (sinon $C$ serait le milieu de $[AB]$, et $A$, $B$, $C$ seraient alignés, ce qui est exclu). Deux cas se présentent :
\begin{itemize}
\item si $\vec{w}$ est colinéaire à $\overrightarrow{AB}$, alors \emph{tout} point $M$ convient : $(\mathcal{D})$ est le plan entier ;
\item sinon, \emph{aucun} point $M$ ne convient : $(\mathcal{D}) = \varnothing$.
\end{itemize}
\textbf{Remarque.} Ce type de question est un piège classique : quand la somme des coefficients est nulle, l'ensemble cherché est soit le plan entier, soit l'ensemble vide — jamais une droite ou un cercle.
\end{exoresolu}

\begin{exoresolu}[Une somme de trois vecteurs colinéaire à une direction]
\textbf{Énoncé.} Soient $A$, $B$, $C$ trois points non alignés et $G = \mathrm{bar}\{(A,1),(B,1),(C,1)\}$. Déterminer l'ensemble des points $M$ tels que $\overrightarrow{MA} + \overrightarrow{MB} + \overrightarrow{MC}$ soit colinéaire à $\overrightarrow{BC}$.
\tcblower
\textbf{Solution.} La somme des coefficients est $3 \neq 0$, donc pour tout $M$ :
\[
\overrightarrow{MA} + \overrightarrow{MB} + \overrightarrow{MC} = 3\overrightarrow{MG}.
\]
La condition « $3\overrightarrow{MG}$ colinéaire à $\overrightarrow{BC}$ » équivaut à « $\overrightarrow{MG}$ colinéaire à $\overrightarrow{BC}$ », c'est-à-dire : $M = G$, ou $M$ appartient à la droite passant par $G$ et parallèle à $(BC)$. Comme $G$ lui-même appartient à cette droite, l'ensemble cherché est exactement \textbf{la parallèle à $(BC)$ passant par $G$}.
\end{exoresolu}

\begin{methode}[Stratégie générale face à un ensemble de points défini vectoriellement]
\begin{enumerate}
\item Calculer $S = \sum \alpha_i$.
\item Si $S \neq 0$ : introduire $G$, réduire en $S\,\overrightarrow{MG}$, puis traduire la condition sur $M$ en une condition sur $MG$ (norme $\Rightarrow$ cercle ; colinéarité $\Rightarrow$ droite ; orthogonalité $\Rightarrow$ droite ou cercle selon le cas).
\item Si $S = 0$ : calculer le vecteur constant $\vec{w}$. L'ensemble est alors le plan entier ou l'ensemble vide, selon que $\vec{w}$ vérifie ou non la condition imposée.
\end{enumerate}
\end{methode}

\begin{savaistu}[L'outil-clé des lignes de niveau au Probatoire]
La réduction de $\sum \alpha_i \overrightarrow{MA_i}$ est l'outil-clé de \emph{presque tous} les problèmes de lignes de niveau posés au Probatoire. Dans les sections suivantes, nous l'utiliserons pour montrer que l'ensemble des points $M$ tels que $MA^2 + MB^2 = k$ est un cercle, que celui des points tels que $MA^2 - MB^2 = k$ est une droite, et que celui des points tels que $\overrightarrow{MA} \cdot \overrightarrow{MB} = k$ est encore un cercle. À chaque fois, le même schéma : \textbf{réduire, traduire, conclure}. Les élèves qui maîtrisent cette réduction traitent ces exercices en quelques lignes ; les autres s'y perdent en calculs interminables. Entraîne-toi : c'est un investissement rentable !
\end{savaistu}

\begin{exercice}[Pour s'entraîner]
Soient $A$, $B$, $C$ trois points du plan tels que $ABC$ soit un triangle équilatéral de côté $6$ cm.
\begin{enumerate}
\item Réduire le vecteur $\vec{u}(M) = 2\overrightarrow{MA} - \overrightarrow{MB} - \overrightarrow{MC}$. Que peut-on dire de sa direction et de sa norme ?
\item Déterminer l'ensemble des points $M$ tels que $\left\| \overrightarrow{MA} + \overrightarrow{MB} + \overrightarrow{MC} \right\| = 12$ cm.
\end{enumerate}
\end{exercice}

\begin{corrige}[Exercice 1]
\textbf{1.} Somme des coefficients : $2 - 1 - 1 = 0$, donc $\vec{u}(M)$ est constant. En remplaçant $M$ par $A$ :
\[
\vec{u}(M) = -\overrightarrow{AB} - \overrightarrow{AC}.
\]
Soit $I$ le milieu de $[BC]$ : $-\overrightarrow{AB} - \overrightarrow{AC} = -2\overrightarrow{AI}$. Or dans un triangle équilatéral de côté $6$, la hauteur vaut $AI = \dfrac{6\sqrt{3}}{2} = 3\sqrt{3}$. Donc $\vec{u}$ est un vecteur constant, de direction $(AI)$ (la médiane issue de $A$), de sens opposé à $\overrightarrow{AI}$, et de norme $2 \times 3\sqrt{3} = 6\sqrt{3}$ cm.

\textbf{2.} Somme des coefficients : $3 \neq 0$. Soit $G$ l'isobarycentre de $A$, $B$, $C$ (le centre de gravité du triangle). Alors $\overrightarrow{MA} + \overrightarrow{MB} + \overrightarrow{MC} = 3\overrightarrow{MG}$, et la condition devient $3\,MG = 12$, soit $MG = 4$. L'ensemble cherché est le cercle de centre $G$ et de rayon $4$ cm.
\end{corrige}

\coursec{Lignes de niveau $MA^2 + MB^2 = k$ et $MA^2 - MB^2 = k$}

Dans la section précédente, nous avons appris à \emph{réduire} une somme vectorielle $\sum_{i=1}^{n} \alpha_i \overrightarrow{MA_i}$ : selon que la somme des coefficients est nulle ou non, cette expression se transforme en un vecteur constant ou en un multiple de $\overrightarrow{MG}$, où $G$ est le barycentre du système. Cette technique de réduction va maintenant nous servir à résoudre un nouveau type de problème, très classique au Probatoire : \textbf{déterminer et construire l'ensemble des points $M$ du plan vérifiant une condition métrique donnée}, par exemple $MA^2 + MB^2 = 50$.

Imagine la situation suivante : deux émetteurs MTN sont installés en deux points $A$ et $B$ distants de \SI{6}{km}. Un téléphone reçoit des signaux dont les intensités dépendent des carrés des distances $MA$ et $MB$. On cherche les positions $M$ où la somme de ces intensités est constante. Mathématiquement, c'est exactement une \cle{ligne de niveau}.

\courssub{Notion de ligne de niveau}

\begin{definition}[Ligne de niveau d'une fonction du plan]
Soit $f$ une fonction qui à tout point $M$ du plan associe un réel $f(M)$, et soit $k$ un réel fixé. On appelle \cle{ligne de niveau} $k$ de $f$ l'ensemble des points $M$ du plan tels que :
\[ f(M) = k. \]
On le note souvent $(\mathcal{L}_k)$ ou $\Gamma_k$.
\end{definition}

\begin{savaistu}[D'où vient ce vocabulaire ?]
Sur une carte topographique (par exemple les cartes de l'Institut National de Cartographie du Cameroun), les \emph{courbes de niveau} relient les points de même altitude : la ligne de niveau $800$ est l'ensemble des points situés à \SI{800}{m} d'altitude. Autour du Mont Cameroun, ces courbes sont des lignes fermées emboîtées. En météorologie, les \emph{isobares} relient les points de même pression. En mathématiques, le principe est identique : on fixe une valeur $k$ et on cherche tous les points $M$ où la fonction vaut exactement $k$.
\end{savaistu}

Dans cette section, $A$ et $B$ sont deux points distincts fixés, et nous étudions les lignes de niveau des fonctions :
\[ f(M) = MA^2 + MB^2 \qquad \text{et} \qquad g(M) = MA^2 - MB^2. \]

\begin{methode}[La démarche générale — à retenir absolument]
Pour déterminer une ligne de niveau faisant intervenir $MA$ et $MB$ :
\begin{enumerate}
\item \textbf{Introduire le milieu $I$ de $[AB]$} (l'isobarycentre de $A$ et $B$).
\item \textbf{Décomposer avec la relation de Chasles} : écrire $\overrightarrow{MA} = \overrightarrow{MI} + \overrightarrow{IA}$ et $\overrightarrow{MB} = \overrightarrow{MI} + \overrightarrow{IB}$.
\item \textbf{Développer} en utilisant le produit scalaire et le fait que $\overrightarrow{IA} + \overrightarrow{IB} = \vec{0}$ (car $I$ est le milieu de $[AB]$).
\item \textbf{Réduire} l'équation à une condition simple sur $MI$ ou sur $\overrightarrow{IM} \cdot \overrightarrow{AB}$.
\item \textbf{Conclure sur la nature de l'ensemble} (cercle, point, ensemble vide, droite) et \textbf{construire}.
\end{enumerate}
\end{methode}

\courssub{Ligne de niveau $MA^2 + MB^2 = k$ : le théorème de la médiane}

\textbf{Intuition.} Quand le point $M$ s'éloigne du segment $[AB]$, les distances $MA$ et $MB$ augmentent toutes les deux, donc la somme $MA^2 + MB^2$ augmente. Par symétrie de la configuration, les points donnant la même somme devraient se répartir « en rond » autour du milieu $I$ de $[AB]$. Nous allons le prouver : l'ensemble cherché sera un cercle de centre $I$.

\begin{propriete}[Théorème de la médiane]
Soit $A$ et $B$ deux points distincts du plan et $I$ le milieu de $[AB]$. Pour tout point $M$ du plan :
\[ MA^2 + MB^2 = 2MI^2 + \frac{AB^2}{2}. \]
\emph{Cette démonstration est exigible au Probatoire.}
\end{propriete}

\begin{experience}[Démonstration du théorème de la médiane]
Partons du membre de gauche et introduisons le point $I$ par la relation de Chasles :
\begin{align*}
MA^2 + MB^2 &= \overrightarrow{MA}^2 + \overrightarrow{MB}^2 \\
&= \left(\overrightarrow{MI} + \overrightarrow{IA}\right)^2 + \left(\overrightarrow{MI} + \overrightarrow{IB}\right)^2 \\
&= MI^2 + 2\,\overrightarrow{MI}\cdot\overrightarrow{IA} + IA^2 + MI^2 + 2\,\overrightarrow{MI}\cdot\overrightarrow{IB} + IB^2 \\
&= 2MI^2 + 2\,\overrightarrow{MI}\cdot\left(\overrightarrow{IA} + \overrightarrow{IB}\right) + IA^2 + IB^2.
\end{align*}
Or $I$ est le milieu de $[AB]$, donc $\overrightarrow{IA} + \overrightarrow{IB} = \vec{0}$, ce qui annule le terme du milieu. De plus $IA = IB = \dfrac{AB}{2}$, donc :
\[ IA^2 + IB^2 = \frac{AB^2}{4} + \frac{AB^2}{4} = \frac{AB^2}{2}. \]
En regroupant : $MA^2 + MB^2 = 2MI^2 + \dfrac{AB^2}{2}$. \hfill$\blacksquare$
\end{experience}

\begin{propriete}[Ensemble des points $M$ tels que $MA^2 + MB^2 = k$]
Soit $A$ et $B$ deux points distincts, $I$ le milieu de $[AB]$ et $k$ un réel. L'ensemble $(\mathcal{C}_k)$ des points $M$ tels que $MA^2 + MB^2 = k$ est donné par la discussion suivante :
\begin{itemize}
\item si $k > \dfrac{AB^2}{2}$ : $(\mathcal{C}_k)$ est le \textbf{cercle de centre $I$} et de rayon $r = \sqrt{\dfrac{2k - AB^2}{4}} = \dfrac{1}{2}\sqrt{2k - AB^2}$ ;
\item si $k = \dfrac{AB^2}{2}$ : $(\mathcal{C}_k)$ est réduit au \textbf{point $I$} ;
\item si $k < \dfrac{AB^2}{2}$ : $(\mathcal{C}_k)$ est \textbf{l'ensemble vide}.
\end{itemize}
\end{propriete}

\textbf{Justification.} D'après le théorème de la médiane, l'équation $MA^2 + MB^2 = k$ équivaut à $2MI^2 + \dfrac{AB^2}{2} = k$, c'est-à-dire :
\[ MI^2 = \frac{k}{2} - \frac{AB^2}{4} = \frac{2k - AB^2}{4}. \]
Le signe du membre de droite décide de tout : s'il est strictement positif, $MI$ est une constante strictement positive, donc $M$ décrit un cercle de centre $I$ ; s'il est nul, $MI = 0$ donc $M = I$ ; s'il est strictement négatif, aucun point ne convient car un carré est toujours positif ou nul.

\begin{attention}[L'erreur classique : oublier la discussion]
Beaucoup d'élèves écrivent automatiquement « l'ensemble est le cercle de centre $I$ et de rayon $\sqrt{\ldots}$ » sans vérifier le signe de la quantité sous la racine. Or si $k < \dfrac{AB^2}{2}$, la quantité $2k - AB^2$ est \textbf{négative} : la racine n'existe pas et l'ensemble est \textbf{vide}. Au Probatoire, la discussion selon les valeurs de $k$ rapporte des points : rédige-la toujours, même rapidement. Retiens aussi la valeur seuil : $k_0 = \dfrac{AB^2}{2}$.
\end{attention}

\begin{exoresolu}[AB = 6 cm : déterminer et construire l'ensemble $MA^2 + MB^2 = 50$]
Soit $A$ et $B$ deux points tels que $AB = 6$ cm. Déterminer et construire l'ensemble $(\Gamma)$ des points $M$ du plan tels que $MA^2 + MB^2 = 50$.
\tcblower
\textbf{Solution détaillée.}

Soit $I$ le milieu de $[AB]$. D'après le théorème de la médiane, pour tout point $M$ :
\[ MA^2 + MB^2 = 2MI^2 + \frac{AB^2}{2} = 2MI^2 + \frac{36}{2} = 2MI^2 + 18. \]
L'équation $MA^2 + MB^2 = 50$ équivaut donc à :
\[ 2MI^2 + 18 = 50 \iff 2MI^2 = 32 \iff MI^2 = 16 \iff MI = 4. \]
\textbf{Discussion :} la valeur seuil est $\dfrac{AB^2}{2} = 18$. Comme $k = 50 > 18$, l'ensemble est non vide.

\textbf{Conclusion :} $(\Gamma)$ est le cercle de centre $I$ (milieu de $[AB]$) et de rayon $4$ cm.

\textbf{Construction :} tracer le segment $[AB]$ de longueur $6$ cm, placer son milieu $I$ (à $3$ cm de $A$), puis tracer le cercle de centre $I$ et de rayon $4$ cm au compas. Vérification : pour le point $M$ du cercle situé sur la droite $(AB)$ du côté de $B$, on a $MA = 7$ et $MB = 1$, donc $MA^2 + MB^2 = 49 + 1 = 50$. \checkmark
\end{exoresolu}

\begin{exemplebox}[Les cas limites avec AB = 6 cm]
Toujours avec $AB = 6$ cm, donc $\dfrac{AB^2}{2} = 18$ :
\begin{itemize}
\item Pour $k = 18$ : l'équation donne $2MI^2 = 0$, soit $MI = 0$. L'ensemble est réduit au \textbf{point $I$}.
\item Pour $k = 12$ : l'équation donne $2MI^2 = 12 - 18 = -6$, impossible car $MI^2 \geq 0$. L'ensemble est \textbf{vide} : aucun point $M$ du plan ne vérifie $MA^2 + MB^2 = 12$.
\end{itemize}
Remarque : la somme $MA^2 + MB^2$ est donc toujours supérieure ou égale à $18$ cm², et sa valeur minimale $18$ est atteinte exactement au milieu $I$.
\end{exemplebox}

\begin{popfigure}
\begin{tikzpicture}[scale=0.85]
% Points A et B
\coordinate (A) at (-3,0);
\coordinate (B) at (3,0);
\coordinate (I) at (0,0);
% Cercles de niveau : k=50 -> r=4 ; k=34 -> r=sqrt(8)=2.83 ; k=26 -> r=2
\draw[popblue, thick] (I) circle (4);
\draw[poporange, thick] (I) circle (2.83);
\draw[popgreen, thick] (I) circle (2);
\fill[poppurple] (I) circle (2.2pt);
% Segment AB
\draw[popdark, thick] (A) -- (B);
\foreach \p/\n/\pos in {A/A/below, B/B/below, I/I/below}
  \fill[popdark] (\p) circle (1.8pt) node[\pos, popdark] {$\n$};
% Labels des cercles
\node[popblue] at (0,4.35) {$k=50$};
\node[poporange] at (0,3.15) {$k=34$};
\node[popgreen] at (0,2.3) {$k=26$};
\node[poppurple] at (1.15,0.35) {$k=18$};
% Rayon
\draw[popblue, ->] (I) -- (45:4) node[midway, above left, popblue] {$r=4$};
\end{tikzpicture}
\legende{Avec $AB = 6$ cm : lignes de niveau de $M \mapsto MA^2 + MB^2$. Ce sont des cercles concentriques de centre $I$ ; pour $k = 18$, l'ensemble se réduit au point $I$ ; en dessous de $18$, il n'y a aucun point.}
\end{popfigure}

\courssub{Ligne de niveau $MA^2 - MB^2 = k$}

\textbf{Intuition.} Cette fois on fait la \emph{différence} des carrés des distances. Si $M$ est sur la médiatrice de $[AB]$, alors $MA = MB$ donc $MA^2 - MB^2 = 0$. Si $M$ se rapproche de $B$, $MB$ diminue et la différence devient positive. On devine que les lignes de niveau seront des lignes « d'équilibre » perpendiculaires à $(AB)$, se décalant le long de la droite $(AB)$ quand $k$ varie.

\begin{propriete}[Réduction de $MA^2 - MB^2$]
Soit $A$ et $B$ deux points distincts et $I$ le milieu de $[AB]$. Pour tout point $M$ du plan :
\[ MA^2 - MB^2 = 2\,\overrightarrow{AB} \cdot \overrightarrow{IM}. \]
\emph{Démonstration exigible.}
\end{propriete}

\begin{experience}[Démonstration]
On utilise l'identité remarquable vectorielle $u^2 - v^2 = (\vec{u} + \vec{v})\cdot(\vec{u} - \vec{v})$ :
\begin{align*}
MA^2 - MB^2 &= \overrightarrow{MA}^2 - \overrightarrow{MB}^2 \\
&= \left(\overrightarrow{MA} + \overrightarrow{MB}\right)\cdot\left(\overrightarrow{MA} - \overrightarrow{MB}\right).
\end{align*}
D'une part, $I$ étant le milieu de $[AB]$, on a $\overrightarrow{MA} + \overrightarrow{MB} = 2\overrightarrow{MI}$ (propriété fondamentale de l'isobarycentre). D'autre part, par Chasles :
\[ \overrightarrow{MA} - \overrightarrow{MB} = \overrightarrow{MA} + \overrightarrow{BM} = \overrightarrow{BM} + \overrightarrow{MA} = \overrightarrow{BA}. \]
Donc :
\[ MA^2 - MB^2 = 2\overrightarrow{MI}\cdot\overrightarrow{BA} = 2\,\overrightarrow{AB}\cdot\overrightarrow{IM}. \]
(en utilisant $\overrightarrow{MI} = -\overrightarrow{IM}$ et $\overrightarrow{BA} = -\overrightarrow{AB}$). \hfill$\blacksquare$
\end{experience}

\begin{propriete}[Ensemble des points $M$ tels que $MA^2 - MB^2 = k$]
Soit $A$ et $B$ deux points distincts, $I$ le milieu de $[AB]$ et $k$ un réel. L'ensemble $(\mathcal{D}_k)$ des points $M$ tels que $MA^2 - MB^2 = k$ est la \textbf{droite perpendiculaire à $(AB)$} passant par le point $H$ de la droite $(AB)$ défini par :
\[ \overrightarrow{IH} = \frac{k}{2\,AB^2}\,\overrightarrow{AB}. \]
En particulier, pour $k = 0$, on retrouve la \textbf{médiatrice de $[AB]$}.
\end{propriete}

\textbf{Justification.} Soit $H$ le point de $(AB)$ défini ci-dessus (il existe et est unique). Pour tout point $M$, décomposons $\overrightarrow{IM} = \overrightarrow{IH} + \overrightarrow{HM}$. Alors :
\[ MA^2 - MB^2 = 2\,\overrightarrow{AB}\cdot\overrightarrow{IM} = 2\,\overrightarrow{AB}\cdot\overrightarrow{IH} + 2\,\overrightarrow{AB}\cdot\overrightarrow{HM} = k + 2\,\overrightarrow{AB}\cdot\overrightarrow{HM}, \]
car $2\,\overrightarrow{AB}\cdot\overrightarrow{IH} = 2 \times \dfrac{k}{2AB^2} \times AB^2 = k$. L'équation $MA^2 - MB^2 = k$ équivaut donc à $\overrightarrow{AB}\cdot\overrightarrow{HM} = 0$, c'est-à-dire $(HM) \perp (AB)$ : l'ensemble est la perpendiculaire à $(AB)$ passant par $H$.

\begin{methode}[Construire la ligne de niveau $MA^2 - MB^2 = k$]
\begin{enumerate}
\item Calculer la valeur seuil : il n'y en a pas ici, une droite existe \textbf{pour tout réel $k$}.
\item Placer le milieu $I$ de $[AB]$, puis déterminer le point $H$ de $(AB)$ grâce à la relation $\overrightarrow{IH} = \dfrac{k}{2AB^2}\,\overrightarrow{AB}$ (attention au sens : si $k > 0$, $H$ est du côté de $B$ ; si $k < 0$, du côté de $A$).
\item Tracer la perpendiculaire à $(AB)$ en $H$ : c'est l'ensemble cherché.
\end{enumerate}
\end{methode}

\begin{exoresolu}[AB = 6 cm : déterminer et construire l'ensemble $MA^2 - MB^2 = 12$]
Soit $A$ et $B$ deux points tels que $AB = 6$ cm. Déterminer et construire l'ensemble $(\Delta)$ des points $M$ tels que $MA^2 - MB^2 = 12$.
\tcblower
\textbf{Solution détaillée.}

Soit $I$ le milieu de $[AB]$. Pour tout point $M$ :
\[ MA^2 - MB^2 = 2\,\overrightarrow{AB}\cdot\overrightarrow{IM}. \]
Cherchons le point $H$ de $(AB)$ tel que $HA^2 - HB^2 = 12$. Posons $H$ sur $(AB)$ avec $\overrightarrow{IH} = x\,\overrightarrow{AB}$. La condition $2\,\overrightarrow{AB}\cdot\overrightarrow{IH} = 12$ donne :
\[ 2x\,AB^2 = 12 \iff 2x \times 36 = 12 \iff x = \frac{1}{6}. \]
Donc $\overrightarrow{IH} = \dfrac{1}{6}\overrightarrow{AB}$, c'est-à-dire $IH = \dfrac{AB}{6} = 1$ cm, $H$ étant placé \textbf{du côté de $B$} (car $k = 12 > 0$).

\textbf{Conclusion :} $(\Delta)$ est la droite perpendiculaire à $(AB)$ passant par le point $H$ de $[AB]$ tel que $IH = 1$ cm (donc $AH = 4$ cm et $BH = 2$ cm).

\textbf{Vérification :} pour $H$ lui-même : $HA^2 - HB^2 = 16 - 4 = 12$. \checkmark
\end{exoresolu}

\begin{popfigure}
\begin{tikzpicture}[scale=0.8]
% Axe AB
\draw[popdark, thick] (-3,0) -- (3,0);
\coordinate (A) at (-3,0);
\coordinate (B) at (3,0);
\coordinate (I) at (0,0);
\foreach \p/\n in {A/A, B/B, I/I}
  \fill[popdark] (\p) circle (1.8pt) node[below, popdark] {$\n$};
% Droites de niveau : k = -24 -> x=-2 ; k=-12 -> x=-1 ; k=0 -> x=0 ; k=12 -> x=1 ; k=24 -> x=2
\draw[popgreen, thick] (-2,-2.2) -- (-2,2.2) node[above, popgreen] {$k=-24$};
\draw[poporange, thick] (-1,-2.2) -- (-1,2.2) node[above, poporange] {$k=-12$};
\draw[poppurple, thick] (0,-2.2) -- (0,2.2) node[above, poppurple] {$k=0$};
\draw[popblue, thick] (1,-2.2) -- (1,2.2) node[above, popblue] {$k=12$};
\draw[popink, thick] (2,-2.2) -- (2,2.2) node[above, popink] {$k=24$};
% Point H pour k=12
\fill[popblue] (1,0) circle (2pt) node[below, popblue] {$H$};
% Angle droit
\draw[popblue] (1,0) ++(0,0.25) -- ++(-0.25,0.25) -- ++(-0.25,0) -- cycle;
\node[poppurple] at (0,-2.6) {médiatrice};
\end{tikzpicture}
\legende{Avec $AB = 6$ cm : lignes de niveau de $M \mapsto MA^2 - MB^2$. Ce sont des droites parallèles, toutes perpendiculaires à $(AB)$ ; la ligne $k = 0$ est la médiatrice de $[AB]$.}
\end{popfigure}

\begin{attention}[Pièges fréquents sur $MA^2 - MB^2 = k$]
\begin{itemize}
\item \textbf{Ne pas confondre} avec $MA^2 + MB^2 = k$ : la somme donne un cercle (ou un point, ou le vide), la différence donne toujours une droite.
\item \textbf{Le sens de $H$ dépend du signe de $k$} : si $k > 0$, $H$ est du côté de $B$ (les points sont plus proches de $B$ que de $A$) ; si $k < 0$, $H$ est du côté de $A$. Une vérification avec $HA^2 - HB^2$ permet de contrôler.
\item \textbf{Ne pas écrire} $MA^2 - MB^2 = (MA - MB)^2$ : c'est faux ! On a $MA^2 - MB^2 = (MA - MB)(MA + MB)$, mais ces longueurs ne se manipulent pas aussi simplement ; passe toujours par les vecteurs et le milieu $I$.
\end{itemize}
\end{attention}

\begin{aretenir}[Bilan de la section]
Soit $A$ et $B$ deux points distincts, $I$ le milieu de $[AB]$ et $k$ un réel.
\begin{itemize}
\item $MA^2 + MB^2 = 2MI^2 + \dfrac{AB^2}{2}$ (théorème de la médiane). L'ensemble $MA^2 + MB^2 = k$ est : un cercle de centre $I$ si $k > \dfrac{AB^2}{2}$, le point $I$ si $k = \dfrac{AB^2}{2}$, vide si $k < \dfrac{AB^2}{2}$.
\item $MA^2 - MB^2 = 2\,\overrightarrow{AB}\cdot\overrightarrow{IM}$. L'ensemble $MA^2 - MB^2 = k$ est toujours une droite perpendiculaire à $(AB)$, passant par le point $H$ de $(AB)$ tel que $\overrightarrow{IH} = \dfrac{k}{2AB^2}\overrightarrow{AB}$.
\item La méthode reine : introduire le milieu $I$, décomposer avec Chasles, développer, réduire, discuter, construire.
\end{itemize}
\end{aretenir}

\begin{exercice}[Pour s'entraîner]
Soit $A$ et $B$ deux points tels que $AB = 8$ cm.
\begin{enumerate}
\item Déterminer et construire l'ensemble des points $M$ tels que $MA^2 + MB^2 = 72$.
\item Que dire de l'ensemble $MA^2 + MB^2 = 30$ ?
\item Déterminer et construire l'ensemble des points $M$ tels que $MA^2 - MB^2 = -16$.
\end{enumerate}
\end{exercice}

\begin{corrige}[Exercice]
Ici $\dfrac{AB^2}{2} = \dfrac{64}{2} = 32$.
\begin{enumerate}
\item $2MI^2 + 32 = 72 \iff MI^2 = 20 \iff MI = \sqrt{20} = 2\sqrt{5}$. Comme $72 > 32$, l'ensemble est le cercle de centre $I$ et de rayon $2\sqrt{5} \approx 4{,}47$ cm.
\item $30 < 32$ : l'ensemble est vide.
\item $2\,\overrightarrow{AB}\cdot\overrightarrow{IH} = -16$ avec $AB^2 = 64$ donne $\overrightarrow{IH} = -\dfrac{16}{128}\overrightarrow{AB} = -\dfrac{1}{8}\overrightarrow{AB}$, soit $IH = 1$ cm du côté de $A$. L'ensemble est la perpendiculaire à $(AB)$ en ce point $H$.
\end{enumerate}
\end{corrige}

\coursec{Lignes de niveau $\overrightarrow{MA}\cdot\overrightarrow{MB}=k$ et $\dfrac{MA}{MB}=k$}

Dans la section précédente, nous avons caractérisé les lignes de niveau $MA^2+MB^2=k$ (cercles centrés au milieu de $[AB]$) et $MA^2-MB^2=k$ (droites perpendiculaires à $(AB)$). Nous poursuivons notre étude avec deux nouvelles conditions, plus subtiles : une condition portant sur le \cle{produit scalaire} $\overrightarrow{MA}\cdot\overrightarrow{MB}$, et une condition portant sur le \cle{rapport des distances} $\dfrac{MA}{MB}$. Tu vas découvrir que ces deux familles de lignes de niveau cachent deux des plus beaux cercles de la géométrie : le cercle de diamètre $[AB]$ et le cercle d'Apollonius.

\courssub{Ligne de niveau $\overrightarrow{MA}\cdot\overrightarrow{MB}=k$}

\textbf{Intuition.} Le produit scalaire $\overrightarrow{MA}\cdot\overrightarrow{MB}$ est lié au cosinus de l'angle $\widehat{AMB}$ : $\overrightarrow{MA}\cdot\overrightarrow{MB}=MA\times MB\times\cos(\widehat{AMB})$. Dire que ce produit est nul, c'est dire que $\cos(\widehat{AMB})=0$, donc que l'angle $\widehat{AMB}$ est droit. Or tu connais depuis la classe de Seconde un théorème célèbre : un triangle rectangle en $M$ est inscrit dans le cercle de diamètre son hypoténuse. L'ensemble cherché pour $k=0$ devrait donc être le cercle de diamètre $[AB]$. C'est exactement ce que nous allons démontrer, en toute généralité pour tout $k$, grâce au milieu $I$ de $[AB]$.

\begin{propriete}[Identité de polarisation au milieu]
Soient $A$ et $B$ deux points du plan et $I$ le milieu de $[AB]$. Pour tout point $M$ du plan :
\[
\overrightarrow{MA}\cdot\overrightarrow{MB}=MI^2-\frac{AB^2}{4}.
\]
\emph{Démonstration exigible au Probatoire.} On décompose par la relation de Chasles en passant par $I$ :
\begin{align*}
\overrightarrow{MA}\cdot\overrightarrow{MB}
&=\left(\overrightarrow{MI}+\overrightarrow{IA}\right)\cdot\left(\overrightarrow{MI}+\overrightarrow{IB}\right)\\
&=MI^2+\overrightarrow{MI}\cdot\overrightarrow{IB}+\overrightarrow{IA}\cdot\overrightarrow{MI}+\overrightarrow{IA}\cdot\overrightarrow{IB}\\
&=MI^2+\overrightarrow{MI}\cdot\left(\overrightarrow{IA}+\overrightarrow{IB}\right)+\overrightarrow{IA}\cdot\overrightarrow{IB}.
\end{align*}
Or $I$ est le milieu de $[AB]$, donc $\overrightarrow{IA}+\overrightarrow{IB}=\vec{0}$, et $\overrightarrow{IB}=-\overrightarrow{IA}$ avec $IA=\dfrac{AB}{2}$. Il reste :
\[
\overrightarrow{MA}\cdot\overrightarrow{MB}=MI^2-IA^2=MI^2-\frac{AB^2}{4}. \qquad \blacksquare
\]
\end{propriete}

\begin{propriete}[Ligne de niveau $\overrightarrow{MA}\cdot\overrightarrow{MB}=k$]
Soient $A$ et $B$ deux points distincts, $I$ le milieu de $[AB]$ et $k$ un réel. L'ensemble $(\mathcal{E})$ des points $M$ tels que $\overrightarrow{MA}\cdot\overrightarrow{MB}=k$ est caractérisé par $MI^2=k+\dfrac{AB^2}{4}$. On distingue trois cas :
\begin{itemize}
\item si $k>-\dfrac{AB^2}{4}$ : $(\mathcal{E})$ est le \cle{cercle} de centre $I$ et de rayon $\sqrt{k+\dfrac{AB^2}{4}}$ ;
\item si $k=-\dfrac{AB^2}{4}$ : $(\mathcal{E})$ est réduit au point $I$ ;
\item si $k<-\dfrac{AB^2}{4}$ : $(\mathcal{E})$ est \cle{vide}.
\end{itemize}
\textbf{Cas remarquable $k=0$ :} $\overrightarrow{MA}\cdot\overrightarrow{MB}=0$ équivaut à $MI=\dfrac{AB}{2}$ : l'ensemble est le \cle{cercle de diamètre $[AB]$}. C'est le théorème de l'angle droit : $\widehat{AMB}=90^\circ$ si et seulement si $M$ appartient au cercle de diamètre $[AB]$ (points $A$ et $B$ exclus si l'on exige un vrai triangle).
\end{propriete}

\begin{popfigure}
\begin{tikzpicture}[scale=1.1]
\coordinate (A) at (-2,0);
\coordinate (B) at (2,0);
\coordinate (I) at (0,0);
\coordinate (M) at (1.2,1.6);
\draw[popblue, thick] (I) circle (2);
\draw[popdark, thick] (A) -- (B);
\draw[poporange, thick] (A) -- (M) -- (B);
\draw[popgreen] (M) ++(-0.35,-0.467) -- ++(0.35,-0.262) -- ++(0.35,0.467);
\foreach \p/\n/\pos in {A/A/left, B/B/right, I/I/below, M/M/above right}
  \fill[popink] (\p) circle (1.5pt) node[\pos, popdark]{$\n$};
\draw[dashed, popmuted] (I) -- (M) node[midway, above, popdark, font=\small]{$\frac{AB}{2}$};
\end{tikzpicture}
\legende{Cas $k=0$ : $\overrightarrow{MA}\cdot\overrightarrow{MB}=0$ si et seulement si $M$ est sur le cercle de diamètre $[AB]$. L'angle $\widehat{AMB}$ est droit.}
\end{popfigure}

\begin{exemplebox}[Avec $AB=6$ cm]
Soit $AB=6$ cm. Déterminons l'ensemble des points $M$ tels que $\overrightarrow{MA}\cdot\overrightarrow{MB}=7$. Ici $\dfrac{AB^2}{4}=9$, donc $MI^2=7+9=16$, c'est-à-dire $MI=4$ : l'ensemble est le cercle de centre $I$ (milieu de $[AB]$) et de rayon $4$ cm. Pour $k=-9$, on obtient le seul point $I$ ; pour $k=-10$, l'ensemble est vide.
\end{exemplebox}

\courssub{Ligne de niveau $\dfrac{MA}{MB}=k$ ($k>0$)}

\textbf{Intuition.} Cherchons les points $M$ « deux fois plus loin de $A$ que de $B$ ». Si $k=1$, on cherche les points équidistants de $A$ et $B$ : c'est la médiatrice de $[AB]$, une droite. Mais si $k\neq 1$ ? Les points ne peuvent plus se répartir symétriquement : l'ensemble va « s'enrouler » autour du point le plus proche. Nous allons montrer que c'est un cercle.

\begin{propriete}[Cas $k=1$]
L'ensemble des points $M$ tels que $MA=MB$ est la \cle{médiatrice} du segment $[AB]$.
\end{propriete}

Supposons désormais $k>0$ et $k\neq 1$. Comme $MA$ et $MB$ sont des distances positives, $\dfrac{MA}{MB}=k$ équivaut à $MA^2=k^2MB^2$, c'est-à-dire :
\[
MA^2-k^2MB^2=0.
\]
Factorisons cette différence de deux carrés « vectorielle » en introduisant deux barycentres. La somme des coefficients $1-k^2$ est non nulle (car $k\neq 1$), ce qui autorise l'introduction des points $I$ (division interne) et $E$ (division externe) définis par :
\[
I=\mathrm{bar}\{(A,1)\,;\,(B,k)\} \qquad \text{et} \qquad E=\mathrm{bar}\{(A,1)\,;\,(B,-k)\}.
\]

\begin{experience}[Démonstration guidée : le cercle d'Apollonius]
Montrons pas à pas que $\dfrac{MA}{MB}=k$ (avec $k>0$, $k\neq 1$) équivaut à dire que $M$ appartient au cercle de diamètre $[IE]$.
\begin{enumerate}
\item \textbf{Réduction par les barycentres.} Par le théorème de réduction de la somme vectorielle (somme des coefficients non nulle), pour tout point $M$ :
\[
\overrightarrow{MA}+k\,\overrightarrow{MB}=(1+k)\,\overrightarrow{MI}
\qquad\text{et}\qquad
\overrightarrow{MA}-k\,\overrightarrow{MB}=(1-k)\,\overrightarrow{ME}.
\]
\item \textbf{Factorisation.} En développant le produit scalaire :
\begin{align*}
\left(\overrightarrow{MA}-k\,\overrightarrow{MB}\right)\cdot\left(\overrightarrow{MA}+k\,\overrightarrow{MB}\right)
&=MA^2-k^2MB^2.
\end{align*}
\item \textbf{Conclusion.} Ainsi $MA^2-k^2MB^2=(1-k)(1+k)\,\overrightarrow{ME}\cdot\overrightarrow{MI}=(1-k^2)\,\overrightarrow{ME}\cdot\overrightarrow{MI}$. Comme $1-k^2\neq 0$ :
\[
MA^2-k^2MB^2=0 \iff \overrightarrow{ME}\cdot\overrightarrow{MI}=0.
\]
\item D'après le cas remarquable précédent ($k=0$ pour le produit scalaire), cela signifie que $M$ appartient au \cle{cercle de diamètre $[IE]$}. $\blacksquare$
\end{enumerate}
\end{experience}

\begin{propriete}[Cercle d'Apollonius]
Soient $A$ et $B$ deux points distincts et $k>0$, $k\neq 1$. L'ensemble des points $M$ tels que $\dfrac{MA}{MB}=k$ est le cercle de diamètre $[IE]$, où
\[
I=\mathrm{bar}\{(A,1)\,;\,(B,k)\}, \qquad E=\mathrm{bar}\{(A,1)\,;\,(B,-k)\}.
\]
Ce cercle est appelé \cle{cercle d'Apollonius} associé au rapport $k$.
\begin{itemize}
\item Le point $I$ (coefficients de même signe) est le point de la droite $(AB)$ \textbf{entre $A$ et $B$} (division interne) tel que $IA=k\,IB$.
\item Le point $E$ (coefficients de signes opposés) est le point de la droite $(AB)$ \textbf{à l'extérieur de $[AB]$} (division externe) tel que $EA=k\,EB$.
\end{itemize}
\end{propriete}

\begin{methode}[Construire l'ensemble $\dfrac{MA}{MB}=k$]
\begin{enumerate}
\item Chercher d'abord les \textbf{deux points de la droite $(AB)$} solutions de $\dfrac{MA}{MB}=k$ :
\begin{itemize}
\item $I$ \textbf{entre $A$ et $B$} tel que $IA=k\,IB$ (division interne) ;
\item $E$ \textbf{hors de $[AB]$} tel que $EA=k\,EB$ (division externe, du côté du point le plus proche, c'est-à-dire $B$ si $k>1$, $A$ si $k<1$).
\end{itemize}
\item Ces deux points sont les extrémités d'un \textbf{diamètre} du cercle cherché.
\item Le centre est le milieu de $[IE]$ et le rayon est $\dfrac{IE}{2}$ : tracer le cercle.
\end{enumerate}
\end{methode}

\begin{exoresolu}[Construction avec $AB=8$ cm et $k=2$]
Soient $A$ et $B$ tels que $AB=8$ cm. Déterminer et construire l'ensemble $(\mathcal{C})$ des points $M$ tels que $\dfrac{MA}{MB}=2$.
\tcblower
\textbf{Recherche des points de $(AB)$.} Soit $M$ sur $(AB)$ avec $MA=2\,MB$.
\begin{itemize}
\item \emph{Point intérieur $I$ (division interne)} : $I$ est entre $A$ et $B$, donc $IA+IB=8$ et $IA=2\,IB$, d'où $3\,IB=8$, soit $IB=\dfrac{8}{3}$ cm et $IA=\dfrac{16}{3}$ cm.
\item \emph{Point extérieur $E$ (division externe)} : $E$ est du côté de $B$ (le point le plus proche car $k>1$). Alors $EA=EB+8$ et $EA=2\,EB$, donc $EB=8$ cm et $EA=16$ cm.
\end{itemize}
\textbf{Conclusion.} $(\mathcal{C})$ est le cercle de diamètre $[IE]$, avec $IE=IB+BE=\dfrac{8}{3}+8=\dfrac{32}{3}$ cm. Le rayon est donc $\dfrac{16}{3}\approx 5{,}33$ cm, et le centre $\Omega$ est le milieu de $[IE]$, situé à $\dfrac{16}{3}$ cm de $I$ et de $E$. On a $\Omega B = \Omega E - BE = \dfrac{16}{3} - 8 = -\dfrac{8}{3}$ (sens $B \to E$) : $\Omega$ est à $\dfrac{8}{3}$ cm de $B$, à l'extérieur de $[AB]$, du côté de $B$ (et de $E$).
\end{exoresolu}

\begin{popfigure}
\begin{tikzpicture}[scale=0.55]
\coordinate (A) at (0,0);
\coordinate (B) at (8,0);
% k=2 : I (interne) à 16/3 de A, E (externe) à 16 de A ; centre à 32/3, rayon 16/3
\coordinate (I) at (5.333,0);
\coordinate (E) at (16,0);
\coordinate (O1) at (10.667,0);
% k=1/2 : symétrique ; I' à 8/3, E' à -8 ; centre à -8/3, rayon 16/3
\coordinate (I2) at (2.667,0);
\coordinate (E2) at (-8,0);
\coordinate (O2) at (-2.667,0);
\draw[popblue, thick] (O1) circle (5.333);
\draw[poporange, thick] (O2) circle (5.333);
\draw[popdark, thick] (-9.5,0) -- (17.5,0);
\foreach \p/\n/\pos in {A/A/above, B/B/above, I/I/below, E/E/below, I2/I'/below, E2/E'/below}
  \fill[popink] (\p) circle (2.5pt) node[\pos, popdark]{$\n$};
\fill[popblue] (O1) circle (2.5pt) node[above, popblue]{$\Omega$};
\fill[poporange] (O2) circle (2.5pt) node[above, poporange]{$\Omega'$};
\node[popblue] at (10.667,4.2) {$\frac{MA}{MB}=2$};
\node[poporange] at (-2.667,4.2) {$\frac{MA}{MB}=\frac12$};
\end{tikzpicture}
\legende{Cercles d'Apollonius pour $AB=8$ cm : en bleu $\frac{MA}{MB}=2$ (diamètre $[IE]$), en orange $\frac{MA}{MB}=\frac12$ (diamètre $[I'E']$). Les deux cercles sont symétriques par rapport à la médiatrice de $[AB]$.}
\end{popfigure}

\begin{attention}[Une droite seulement si $k=1$]
Erreur fréquente : affirmer que l'ensemble $\dfrac{MA}{MB}=k$ est toujours une droite. C'est \textbf{faux} ! C'est une droite (la médiatrice de $[AB]$) \textbf{uniquement si $k=1$}. Pour tout $k>0$ différent de $1$, c'est un cercle. Autre piège : oublier que $k$ est un rapport de distances, donc $k>0$ ; et ne pas oublier d'élever au carré avant de factoriser : on utilise $MA^2-k^2MB^2=0$, jamais $MA-kMB=0$ avec des longueurs.
\end{attention}

\begin{savaistu}[Apollonius de Perga]
Apollonius de Perga (vers 262 -- 190 av. J.-C.), mathématicien grec surnommé « le Grand Géomètre », a étudié ces cercles plus de deux siècles avant Jésus-Christ, bien avant l'invention des coordonnées et des vecteurs ! Son traité \emph{Des coniques} est resté une référence pendant plus de quinze siècles. Aujourd'hui, les cercles d'Apollonius interviennent en géolocalisation : si un téléphone à Douala reçoit un signal deux fois plus fort d'une antenne MTN que d'une antenne voisine, sa position possible est un cercle d'Apollonius.
\end{savaistu}

\courssub{Synthèse : les quatre familles de lignes de niveau}

\begin{aretenir}[Tableau récapitulatif]
Dans tout le tableau, $A$ et $B$ sont deux points distincts, $I$ est le milieu de $[AB]$.
\begin{center}
\begin{tabularx}{\linewidth}{|l|X|}
\hline
\rowcolor{popblueL} \textbf{Condition} & \textbf{Ensemble des points $M$} \\
\hline
$MA^2+MB^2=k$ & Cercle de centre $I$ et de rayon $\sqrt{\dfrac{k}{2}-\dfrac{AB^2}{4}}$ si $k>\dfrac{AB^2}{2}$ (point $I$ si égalité, vide sinon) \\
\hline
$MA^2-MB^2=k$ & Droite perpendiculaire à $(AB)$, d'équation vectorielle $2\,\overrightarrow{AB}\cdot\overrightarrow{IM}=k$ (équivalent : $\overrightarrow{AB}\cdot\overrightarrow{IH}=\frac{k}{2}$ pour $H$ pied de la perpendiculaire) \\
\hline
$\overrightarrow{MA}\cdot\overrightarrow{MB}=k$ & Cercle de centre $I$ et de rayon $\sqrt{k+\dfrac{AB^2}{4}}$ si $k>-\dfrac{AB^2}{4}$ ; pour $k=0$ : cercle de diamètre $[AB]$ \\
\hline
$\dfrac{MA}{MB}=k$ ($k>0$) & Médiatrice de $[AB]$ si $k=1$ ; cercle d'Apollonius de diamètre $[IE]$ si $k\neq 1$ ($I$ division interne, $E$ division externe) \\
\hline
\end{tabularx}
\end{center}
\end{aretenir}

\begin{exercice}[Pour t'entraîner]
Soient $A$ et $B$ deux points tels que $AB=10$ cm.
\begin{enumerate}
\item Déterminer et construire l'ensemble des points $M$ tels que $\overrightarrow{MA}\cdot\overrightarrow{MB}=11$.
\item Déterminer et construire l'ensemble des points $M$ tels que $\dfrac{MA}{MB}=3$.
\end{enumerate}
\end{exercice}

\begin{corrige}[Exercice]
\textbf{1.} Avec $I$ milieu de $[AB]$ : $MI^2=11+\dfrac{100}{4}=36$, donc $MI=6$ cm. C'est le cercle de centre $I$ et de rayon $6$ cm.
\textbf{2.} Points de $(AB)$ : $I$ entre $A$ et $B$ avec $IA=3\,IB$ et $IA+IB=10$, d'où $IB=2{,}5$ cm, $IA=7{,}5$ cm ; $E$ au-delà de $B$ avec $EA=3\,EB$ et $EA-EB=10$, d'où $EB=5$ cm, $EA=15$ cm. L'ensemble est le cercle de diamètre $[IE]$ avec $IE=IB+BE=2{,}5+5=7{,}5$ cm, soit un rayon de $3{,}75$ cm.
\end{corrige}

\begin{savaistu}[Vers la Terminale (hors strict programme de Première)]
Ces techniques se généralisent : en Terminale, tu étudieras des lignes de niveau faisant intervenir trois points ou plus, grâce à la réduction de $\sum \alpha_i\overrightarrow{MA_i}$. Le cercle d'Apollonius se réinterprétera avec les \textbf{homothéties} ($I$ et $E$ sont les centres des deux homothéties transformant un cercle en un autre) et avec les \textbf{nombres complexes}, où la condition $\dfrac{|z-a|}{|z-b|}=k$ décrira directement ces cercles dans le plan complexe. Tout le travail accompli cette année sur les barycentres portera alors ses fruits !
\end{savaistu}

\newpage

\coursec{Activité d'intégration}

\courssub{Objectifs de l'activité}

À l'issue de cette activité d'intégration, tu seras capable de :
\begin{itemize}
\item modéliser un problème concret d'implantation à l'aide du \cle{barycentre} de points pondérés ;
\item déterminer et construire le barycentre de deux, puis de trois points pondérés, en utilisant la technique des \cle{barycentres partiels} ;
\item réduire une somme vectorielle du type $\sum \alpha_i \overrightarrow{MA_i}$ selon que la somme des coefficients est nulle ou non ;
\item déterminer, caractériser et construire une \cle{ligne de niveau} du type $MA^2 + MB^2 = k$ ;
\item caractériser les ensembles de points définis par $\dfrac{MA}{MB} = k$ (médiatrice et cercle d'Apollonius) ;
\item travailler en groupe, argumenter, rédiger un rapport et présenter une figure à l'échelle.
\end{itemize}

\courssub{Situation}

Le quartier de Mvog-Mbi, à Yaoundé, connaît une croissance rapide de sa population. L'opérateur de téléphonie mobile qui dessert la zone décide d'installer une nouvelle \cle{antenne relais} pour améliorer la qualité du réseau. Trois localités sont concernées par le projet : le village $A$ (Mvog-Mbi centre), le village $B$ (Nkol-Bikok) et le village $C$ (Etoug-Ebe).

Sur le plan de l'ingénieur, les trois villages sont modélisés par trois points du plan :
\begin{itemize}
\item $A$ et $B$ sont distants de $AB = \SI{10}{km}$ ;
\item le village $C$ est situé de sorte que $AC = \SI{8}{km}$ et $BC = \SI{6}{km}$ (on remarquera que $6^2 + 8^2 = 10^2$ : le triangle $ABC$ est rectangle en $C$).
\end{itemize}

L'ingénieur répartit le coût du raccordement en tenant compte du \textbf{trafic téléphonique} de chaque village : plus un village génère de trafic, plus il est important de rapprocher l'antenne de lui. Le coût total est modélisé par la quantité $3\,MA^2 + 2\,MB^2$ (les distances interviennent au carré car la perte de signal et le coût en câble croissent quadratiquement), et l'on cherche le point $M$ qui la minimise. Les trafics estimés donnent les pondérations suivantes : $A$ compte pour $3$, $B$ compte pour $2$, puis $C$, qui rejoint le projet dans un second temps, compte pour $1$.

Par ailleurs, l'émetteur est réglé de telle sorte qu'un point $M$ du territoire reçoit un \textbf{signal de bonne qualité} lorsque
\[ MA^2 + MB^2 \leq 68 \quad (\text{en } \mathrm{km}^2), \]
l'égalité $MA^2 + MB^2 = 68$ définissant la limite de la zone de couverture principale.

L'entreprise te missionne, toi et ton groupe, pour répondre aux questions d'implantation et de couverture, et pour produire une \textbf{figure à l'échelle $1 \text{ cm} \leftrightarrow 1 \text{ km}$} qui servira de document officiel au conseil municipal.

\courssub{Consignes}

\textbf{Partie 1 — Implantation avec deux villages}

\begin{enumerate}
\item On admet que le point qui minimise la quantité $\alpha\,MA^2 + \beta\,MB^2$ (avec $\alpha, \beta > 0$) est le barycentre de $\{(A,\alpha);(B,\beta)\}$. Justifier alors que le point d'implantation optimale $G$ est le barycentre du système $\{(A, 3) ; (B, 2)\}$.
\item Exprimer $\overrightarrow{AG}$ en fonction de $\overrightarrow{AB}$, puis calculer les distances $AG$ et $GB$.
\item Construire $G$ sur la figure à l'échelle. Vérifier que $G$ est bien aligné avec $A$ et $B$.
\end{enumerate}

\textbf{Partie 2 — Le village $C$ rejoint le projet}

\begin{enumerate}
\setcounter{enumi}{3}
\item En utilisant le barycentre partiel $G$ de la Partie 1, montrer que le nouveau point optimal $G'$ est le barycentre de $\{(G, 5) ; (C, 1)\}$.
\item Exprimer $\overrightarrow{GG'}$ en fonction de $\overrightarrow{GC}$ et calculer la distance $GG'$.
\item Construire $G'$ sur la figure. Le nouveau site est-il très éloigné de l'ancien ? Commenter.
\end{enumerate}

\textbf{Partie 3 — Zone limite de couverture}

\begin{enumerate}
\setcounter{enumi}{6}
\item Soit $I$ le milieu de $[AB]$. En utilisant la réduction de $MA^2 + MB^2$ (formule de la médiane), montrer que l'équation $MA^2 + MB^2 = 68$ équivaut à $MI^2 = 9$.
\item En déduire la nature exacte de la zone limite de couverture, son centre et son rayon.
\item Construire cette ligne de niveau sur la figure. Le village $C$ est-il bien couvert ? Justifier par un calcul.
\end{enumerate}

\textbf{Partie 4 — Bonus : points « équidistants en signal »}

\begin{enumerate}
\setcounter{enumi}{9}
\item Déterminer et caractériser l'ensemble des points $M$ tels que $\dfrac{MA}{MB} = 1$.
\item On admet que l'ensemble des points $M$ tels que $\dfrac{MA}{MB} = \dfrac{3}{2}$ est un cercle (cercle d'Apollonius). Déterminer les deux points de la droite $(AB)$ appartenant à cet ensemble, en utilisant des barycentres, puis préciser le centre et le rayon du cercle.
\end{enumerate}

\textbf{Organisation du travail :} le travail se fait en groupes de quatre. Chaque groupe désigne un \textbf{rapporteur} qui présentera la figure et les réponses en dix minutes. La figure, réalisée sur papier millimétré à l'échelle $1 \text{ cm} \leftrightarrow 1 \text{ km}$, comportera : le triangle $ABC$, les points $G$, $G'$, $I$, la zone de couverture et la médiatrice de $[AB]$.

\courssub{Ressources à mobiliser}

\begin{aretenir}[Boîte à outils de l'activité]
\begin{itemize}
\item \textbf{Barycentre de deux points :} $G = \mathrm{bar}\{(A, \alpha) ; (B, \beta)\}$ avec $\alpha + \beta \neq 0$ est l'unique point tel que $\alpha \overrightarrow{GA} + \beta \overrightarrow{GB} = \vec{0}$, et l'on a $\overrightarrow{AG} = \dfrac{\beta}{\alpha + \beta}\overrightarrow{AB}$. Le barycentre est toujours \cle{aligné} avec $A$ et $B$. Si $\alpha = \beta$, on parle d'\cle{isobarycentre} : c'est le milieu de $[AB]$.
\item \textbf{Homogénéité :} le barycentre ne change pas si tous les coefficients sont multipliés par un même réel non nul.
\item \textbf{Barycentre partiel :} pour chercher le barycentre de $\{(A, \alpha) ; (B, \beta) ; (C, \gamma)\}$ avec $\alpha + \beta \neq 0$, on regroupe : $G' = \mathrm{bar}\{(G, \alpha + \beta) ; (C, \gamma)\}$ où $G = \mathrm{bar}\{(A, \alpha) ; (B, \beta)\}$.
\item \textbf{Réduction vectorielle :} si $\sum \alpha_i \neq 0$, alors pour tout point $M$ :
\[ \sum_{i=1}^{n} \alpha_i \overrightarrow{MA_i} = \left(\sum_{i=1}^{n} \alpha_i\right) \overrightarrow{MG}, \quad \text{où } G = \mathrm{bar}\{(A_i, \alpha_i)\}. \]
Si $\sum \alpha_i = 0$, la somme $\sum \alpha_i \overrightarrow{MA_i}$ est un vecteur constant, indépendant de $M$.
\item \textbf{Ligne de niveau $MA^2 + MB^2 = k$ :} si $I$ est le milieu de $[AB]$, alors $MA^2 + MB^2 = 2MI^2 + \dfrac{AB^2}{2}$. L'ensemble cherché est un cercle de centre $I$ (éventuellement réduit à un point ou vide).
\item \textbf{Ligne de niveau $\dfrac{MA}{MB} = k$ ($k > 0$) :} si $k = 1$, c'est la médiatrice de $[AB]$ ; si $k \neq 1$, c'est un cercle de diamètre $[G_1G_2]$, où $G_1 = \mathrm{bar}\{(A,1);(B,-k)\}$ et $G_2 = \mathrm{bar}\{(A,1);(B,k)\}$ (les sommes $1 - k$ et $1 + k$ sont non nulles car $k \neq 1$ et $k > 0$).
\end{itemize}
\end{aretenir}

\courssub{Démarche et corrigé commenté}

\begin{exoresolu}[L'antenne relais de Mvog-Mbi : corrigé complet]
\textbf{Partie 1.} 1) Justifier que $G = \mathrm{bar}\{(A,3);(B,2)\}$. 2) Calculer $\overrightarrow{AG}$, $AG$ et $GB$. 3) Construire $G$. \textbf{Partie 2.} 4) Montrer que $G' = \mathrm{bar}\{(G,5);(C,1)\}$. 5) Calculer $GG'$. 6) Construire $G'$ et commenter. \textbf{Partie 3.} 7) Réduire $MA^2 + MB^2 = 68$. 8) Nature de la ligne. 9) $C$ est-il couvert ? \textbf{Partie 4.} 10) Ensemble $\frac{MA}{MB}=1$. 11) Cercle d'Apollonius pour $\frac{MA}{MB}=\frac{3}{2}$.
\tcblower
\textbf{Partie 1 — Implantation avec deux villages.}

\textbf{1.} Le coût total étant modélisé par $3\,MA^2 + 2\,MB^2$, le point qui minimise cette quantité est, d'après le résultat admis, le point d'équilibre du système pondéré, c'est-à-dire le \cle{barycentre} $G = \mathrm{bar}\{(A, 3) ; (B, 2)\}$. Comme $\alpha + \beta = 3 + 2 = 5 \neq 0$, ce barycentre existe et est unique.

\textbf{2.} Par définition, $3\overrightarrow{GA} + 2\overrightarrow{GB} = \vec{0}$. En introduisant $A$ via la relation de Chasles dans $\overrightarrow{GB} = \overrightarrow{GA} + \overrightarrow{AB}$ :
\begin{align*}
3\overrightarrow{GA} + 2\left(\overrightarrow{GA} + \overrightarrow{AB}\right) &= \vec{0} \\
5\overrightarrow{GA} + 2\overrightarrow{AB} &= \vec{0} \\
\overrightarrow{AG} &= \frac{2}{5}\overrightarrow{AB}.
\end{align*}
Donc $AG = \dfrac{2}{5} \times 10 = \SI{4}{km}$ et $GB = 10 - 4 = \SI{6}{km}$.

\textbf{3.} Sur la figure à l'échelle ($1 \text{ cm} \leftrightarrow 1 \text{ km}$), on trace $[AB]$ de longueur $\SI{10}{cm}$ et l'on place $G$ sur $[AB]$ à $\SI{4}{cm}$ de $A$. Le point $G$ est bien \cle{aligné} avec $A$ et $B$, comme tout barycentre de deux points.

\textbf{Partie 2 — Le village $C$ rejoint le projet.}

\textbf{4.} On cherche $G' = \mathrm{bar}\{(A,3);(B,2);(C,1)\}$, avec $3+2+1 = 6 \neq 0$. Comme $3 + 2 = 5 \neq 0$, on applique la technique du \cle{barycentre partiel} : en regroupant $A$ et $B$ en leur barycentre $G$ affecté de la somme $5$, on obtient
\[ G' = \mathrm{bar}\{(G, 5) ; (C, 1)\}. \]

\textbf{5.} De $5\overrightarrow{G'G} + \overrightarrow{G'C} = \vec{0}$, on tire, en introduisant $G$ dans $\overrightarrow{G'C} = \overrightarrow{G'G} + \overrightarrow{GC}$ :
\[ 5\overrightarrow{G'G} + \overrightarrow{G'G} + \overrightarrow{GC} = \vec{0} \quad \Longrightarrow \quad 6\overrightarrow{G'G} = -\overrightarrow{GC} \quad \Longrightarrow \quad \overrightarrow{GG'} = \frac{1}{6}\overrightarrow{GC}. \]
Calculons $GC$. Le triangle $ABC$ est rectangle en $C$ (car $AC^2 + BC^2 = 64 + 36 = 100 = AB^2$). Plaçons-nous dans le repère orthonormé $(C\,;\vec{i},\vec{j})$ avec $\vec{i}$ unitaire sur $(CA)$ et $\vec{j}$ unitaire sur $(CB)$ : alors $A(8;0)$, $B(0;6)$. Le point $G$ vérifie $\overrightarrow{AG} = \frac{2}{5}\overrightarrow{AB}$, donc
\[ G = A + \frac{2}{5}(B - A) = \left(8 - \frac{16}{5}\,;\, \frac{12}{5}\right) = \left(\frac{24}{5}\,;\,\frac{12}{5}\right). \]
Alors $GC^2 = \left(\frac{24}{5}\right)^2 + \left(\frac{12}{5}\right)^2 = \frac{576 + 144}{25} = \frac{720}{25} = \frac{144}{5}$, d'où $GC = \dfrac{12}{\sqrt{5}} = \dfrac{12\sqrt{5}}{5} \approx \SI{5,37}{km}$, et
\[ GG' = \frac{1}{6}GC \approx \SI{0,89}{km}. \]

\textbf{6.} On construit $G'$ sur le segment $[GC]$, au sixième de $GC$ en partant de $G$ (soit environ $\SI{0,9}{cm}$ de $G$ sur la figure). Le déplacement est faible : le trafic de $C$ (coefficient $1$) est modeste devant ceux de $A$ et $B$ (somme $5$), donc l'implantation optimale bouge peu. C'est un résultat rassurant pour l'ingénieur : pas besoin de déplacer l'antenne déjà installée de manière significative.

\textbf{Partie 3 — Zone limite de couverture.}

\textbf{7.} Soit $I$ le milieu de $[AB]$ (c'est l'\cle{isobarycentre} de $A$ et $B$). Pour tout point $M$, la formule de la médiane (obtenue en écrivant $\overrightarrow{MA} = \overrightarrow{MI} + \overrightarrow{IA}$ et $\overrightarrow{MB} = \overrightarrow{MI} + \overrightarrow{IB}$, en développant les carrés scalaires et en utilisant $\overrightarrow{IA} + \overrightarrow{IB} = \vec{0}$) donne :
\[ MA^2 + MB^2 = 2MI^2 + \frac{AB^2}{2}. \]
Avec $AB = 10$, on a $\dfrac{AB^2}{2} = \dfrac{100}{2} = 50$. L'équation $MA^2 + MB^2 = 68$ équivaut donc à :
\begin{align*}
2MI^2 + 50 &= 68 \\
2MI^2 &= 18 \\
MI^2 &= 9.
\end{align*}

\textbf{8.} $MI^2 = 9$ équivaut à $MI = 3$ : la zone limite de couverture est le \textbf{cercle de centre $I$ (milieu de $[AB]$) et de rayon $\SI{3}{km}$}. Les points bien couverts sont les points du disque fermé ($MI \leq 3$), la limite étant le cercle lui-même.

\textbf{9.} Construction : sur la figure, on place $I$ à $\SI{5}{cm}$ de $A$ sur $[AB]$ et l'on trace le cercle de centre $I$ et de rayon $\SI{3}{cm}$. Le village $C$ est-il couvert ? Calculons $IC$. Dans le triangle $ABC$ rectangle en $C$, la médiane issue de l'angle droit vérifie $IC = \dfrac{AB}{2} = \SI{5}{km}$ (le triangle rectangle est inscrit dans le demi-cercle de diamètre $[AB]$, donc $I$, milieu de l'hypoténuse, est équidistant de $A$, $B$ et $C$). Comme $IC = 5 > 3$, le village $C$ est \textbf{hors de la zone de bonne couverture} : il faudra augmenter la puissance de l'émetteur ou installer un répéteur. C'est une conclusion concrète et utile au conseil municipal.

\textbf{Partie 4 — Bonus.}

\textbf{10.} $\dfrac{MA}{MB} = 1$ équivaut à $MA = MB$ : l'ensemble cherché est la \textbf{médiatrice du segment $[AB]$}, droite perpendiculaire à $(AB)$ passant par $I$. Les points de cette médiatrice reçoivent un signal « équilibré » des deux villages.

\textbf{11.} Pour $\dfrac{MA}{MB} = \dfrac{3}{2}$, on a $2MA = 3MB$. Cherchons les points de la droite $(AB)$ solutions. Un point $M$ de $(AB)$ vérifie $2MA = 3MB$ dans deux configurations :
\begin{itemize}
\item $2\overrightarrow{MA} - 3\overrightarrow{MB} = \vec{0}$, c'est-à-dire $M = G_1 = \mathrm{bar}\{(A, 2) ; (B, -3)\}$ (somme $2 - 3 = -1 \neq 0$) : $\overrightarrow{AG_1} = \dfrac{-3}{-1}\overrightarrow{AB} = 3\overrightarrow{AB}$, donc $G_1$ est à $\SI{30}{km}$ de $A$, au-delà de $B$ ;
\item $2\overrightarrow{MA} + 3\overrightarrow{MB} = \vec{0}$, c'est-à-dire $M = G_2 = \mathrm{bar}\{(A, 2) ; (B, 3)\}$ (somme $5 \neq 0$) : $\overrightarrow{AG_2} = \dfrac{3}{5}\overrightarrow{AB}$, donc $AG_2 = \SI{6}{km}$ et $G_2 \in [AB]$.
\end{itemize}
(On retrouve les barycentres $\mathrm{bar}\{(A,1);(B,-k)\}$ et $\mathrm{bar}\{(A,1);(B,k)\}$ de la boîte à outils avec $k = \frac{3}{2}$, multipliés par $2$ grâce à l'\cle{homogénéité}.) L'ensemble des points tels que $\dfrac{MA}{MB} = \dfrac{3}{2}$ est le \textbf{cercle de diamètre $[G_1G_2]$} (cercle d'Apollonius). Son centre $\Omega$ est le milieu de $[G_1G_2]$ : sur la droite $(AB)$ orientée de $A$ vers $B$, $G_2$ a pour abscisse $6$ et $G_1$ l'abscisse $30$, donc $\Omega$ a pour abscisse $\dfrac{6+30}{2} = 18$ (soit $\SI{18}{km}$ de $A$) et le rayon vaut $\dfrac{G_1G_2}{2} = \dfrac{30 - 6}{2} = \SI{12}{km}$. Ce cercle est très grand : il déborde largement du territoire communal, ce qui montre qu'un rapport de signaux de $3/2$ ne peut être obtenu que loin de $B$.
\end{exoresolu}

\begin{attention}[Le barycentre est plus proche du point le plus pondéré]
Dans la Partie 1, on constate que $G$ est à $\SI{4}{km}$ de $A$ (coefficient $3$) et à $\SI{6}{km}$ de $B$ (coefficient $2$) : le barycentre est \textbf{plus proche du point de plus fort coefficient}. C'est cohérent avec la situation : le village $A$ génère plus de trafic, l'antenne doit donc être plus proche de $A$. Erreur fréquente : écrire $\overrightarrow{AG} = \frac{3}{5}\overrightarrow{AB}$ en inversant les coefficients — retiens que dans $\overrightarrow{AG} = \dfrac{\beta}{\alpha + \beta}\overrightarrow{AB}$, c'est le coefficient de \emph{l'autre point} ($B$) qui apparaît au numérateur. Autre piège : n'oublie jamais de vérifier que la somme des coefficients est non nulle avant de conclure à l'existence du barycentre.
\end{attention}

\begin{popfigure}
\begin{center}
\begin{tikzpicture}[scale=0.62,>=Stealth]
% Triangle ABC rectangle en C, AB = 10
\coordinate (C) at (0,0);
\coordinate (A) at (8,0);
\coordinate (B) at (0,6);
% Points barycentriques
\coordinate (G) at (4.8,2.4);   % G = (24/5, 12/5)
\coordinate (Gp) at (4,2);      % G' = G + (1/6)(C - G) = (4, 2)
\coordinate (I) at (4,3);       % milieu de AB
% Cercle de couverture
\draw[poporange,thick,fill=poporangeL,fill opacity=0.5] (I) circle (3);
% Triangle
\draw[popdark,thick] (A)--(B)--(C)--cycle;
\draw[popdark] (0,0.45)--(0.45,0.45)--(0.45,0);
% Segment GC et point G'
\draw[popblue,dashed,thick] (G)--(C);
% Médiatrice de [AB] (perpendiculaire à AB passant par I)
\draw[poppurple,thick] ($(I)!1.6!90:(A)$)--($(I)!1.6!-90:(A)$);
% Points
\foreach \p/\pos/\col in {A/right/popdark, B/above/popdark, C/below left/popdark, G/above right/popblue, Gp/below right/popblue, I/above left/poporange}
  \fill[\col] (\p) circle (2.2pt) node[\pos,\col,font=\small]{$\p$};
\node[poporange,font=\small] at (7.2,4.6) {Couverture : $MI = 3$};
\node[poppurple,font=\small,rotate=32] at (6.4,1.2) {$\frac{MA}{MB}=1$};
\end{tikzpicture}
\end{center}
\legende{Figure de synthèse à l'échelle $1 \text{ cm} \leftrightarrow 1 \text{ km}$ : le triangle $ABC$ rectangle en $C$, le barycentre $G$ de $\{(A,3);(B,2)\}$, le barycentre $G'$ obtenu par barycentre partiel, la zone de couverture (cercle de centre $I$ et de rayon $\SI{3}{km}$) et la médiatrice de $[AB]$.}
\end{popfigure}

\courssub{Bilan de l'activité}

Cette activité a permis de mettre en évidence les points essentiels de la leçon dans un contexte professionnel réaliste :

\begin{itemize}
\item \textbf{Le barycentre comme point d'équilibre économique.} Pondérer les villages par leurs trafics transforme un problème de coût en un problème de barycentre. On a vérifié concrètement que le barycentre se rapproche du point de plus fort coefficient et qu'il est toujours aligné avec les deux points.
\item \textbf{La puissance des barycentres partiels.} L'arrivée du village $C$ n'a pas obligé à refaire tous les calculs : regrouper $A$ et $B$ en $G$ (affecté de la somme $5$) a ramené le problème à trois points à un problème à deux points. C'est la stratégie à retenir pour le Probatoire.
\item \textbf{La réduction vectorielle au service des lignes de niveau.} L'équation $MA^2 + MB^2 = 68$, qui semblait compliquée, s'est réduite grâce à la formule de la médiane à $MI = 3$ : la ligne de niveau est un cercle dont on connaît le centre (le milieu $I$, isobarycentre de $A$ et $B$) et le rayon.
\item \textbf{Le lien entre barycentres et cercle d'Apollonius.} Les ensembles $\frac{MA}{MB} = k$ se traitent avec les barycentres $\mathrm{bar}\{(A,1);(B,-k)\}$ et $\mathrm{bar}\{(A,1);(B,k)\}$ : médiatrice si $k = 1$, cercle sinon.
\item \textbf{Des compétences transversales.} Travailler à l'échelle, vérifier une hypothèse par le calcul (le village $C$ hors couverture), rédiger un rapport et défendre ses conclusions à l'oral sont des savoir-faire et savoir-être directement mobilisables dans la vie professionnelle — et au Probatoire, où la qualité de la rédaction et de la figure compte autant que le résultat.
\end{itemize}

\begin{methode}[Réflexes à retenir pour tout problème de ce type]
\begin{enumerate}
\item Traduire les pondérations de l'énoncé en coefficients et vérifier que la somme est non nulle.
\item Pour trois points ou plus, regrouper par barycentres partiels en choisissant des regroupements dont la somme des coefficients est non nulle.
\item Pour une ligne de niveau $MA^2 + MB^2 = k$, passer par le milieu $I$ de $[AB]$ et la formule $MA^2 + MB^2 = 2MI^2 + \frac{AB^2}{2}$, puis discuter selon le signe de $k - \frac{AB^2}{2}$ : cercle si $k - \frac{AB^2}{2} > 0$, point $I$ si $k - \frac{AB^2}{2} = 0$, ensemble vide si $k - \frac{AB^2}{2} < 0$.
\item Toujours conclure en langage géométrique clair : « l'ensemble cherché est le cercle de centre ... et de rayon ... », puis construire proprement à l'échelle demandée.
\end{enumerate}
\end{methode}

\newpage

\coursec{Méthodes et exercices résolus}

\coursec{Barycentres et lignes de niveau}

\courssub{Barycentre de deux points pondérés : définition, détermination et construction}

\begin{definition}[Barycentre de deux points pondérés]
Soient $A$ et $B$ deux points distincts du plan, et $\alpha, \beta$ deux réels tels que $\alpha+\beta \neq 0$.
Le \textbf{barycentre} du système de points pondérés $\{(A;\alpha), (B;\beta)\}$ est l'unique point $G$ du plan vérifiant :
\[
\alpha \overrightarrow{GA} + \beta \overrightarrow{GB} = \vec{0}
\]
On note $G = \text{Bar}\{(A;\alpha), (B;\beta)\}$.
\end{definition}

\begin{propriete}[Propriétés fondamentales]
\begin{enumerate}
    \item \textbf{Expression vectorielle :} Pour tout point $M$ du plan,
    \[
    \overrightarrow{MG} = \frac{\alpha \overrightarrow{MA} + \beta \overrightarrow{MB}}{\alpha+\beta}
    \]
    \item \textbf{Coordonnées (repère orthonormé) :} Si $A(x_A,y_A), B(x_B,y_B)$, alors
    \[
    G\left(\frac{\alpha x_A + \beta x_B}{\alpha+\beta}\ ;\ \frac{\alpha y_A + \beta y_B}{\alpha+\beta}\right)
    \]
    \item \textbf{Alignement :} $G$ est \textbf{toujours aligné} avec $A$ et $B$.
    \item \textbf{Rapport signé :} $\dfrac{\overrightarrow{GA}}{\overrightarrow{GB}} = -\dfrac{\beta}{\alpha}$.
    \item \textbf{Rapport des distances :} $\dfrac{GA}{GB} = \left|\dfrac{\beta}{\alpha}\right|$.
    \item \textbf{Invariance par homothétie des coefficients :} Multiplier $\alpha$ et $\beta$ par un même réel non nul ne change pas $G$.
    \item \textbf{Position de $G$ par rapport à $[AB]$ :}
    \begin{itemize}
        \item Si $\alpha$ et $\beta$ ont \textbf{le même signe} ($\alpha\beta > 0$) : $G \in [AB]$ (barycentre intérieur).
        \item Si $\alpha$ et $\beta$ ont \textbf{signes contraires} ($\alpha\beta < 0$) : $G \notin [AB]$ (barycentre extérieur). Il se trouve du côté du point dont le coefficient a la plus grande valeur absolue.
    \end{itemize}
\end{enumerate}
\end{propriete}

\begin{methode}[Déterminer et construire le barycentre de deux points pondérés]
\textbf{Objectif :} Trouver le point $G$ barycentre de $\{(A;\alpha), (B;\beta)\}$ avec $\alpha+\beta \neq 0$.
\begin{enumerate}
    \item \textbf{Calcul vectoriel / Coordonnées (Détermination) :}
    \begin{itemize}
        \item Vérifier que $\alpha+\beta \neq 0$.
        \item Appliquer la formule : $\overrightarrow{MG} = \dfrac{\alpha \overrightarrow{MA} + \beta \overrightarrow{MB}}{\alpha+\beta}$ pour tout $M$ (souvent $M=A$ ou $M=O$ origine).
        \item En coordonnées : $G\left(\dfrac{\alpha x_A + \beta x_B}{\alpha+\beta}\ ;\ \dfrac{\alpha y_A + \beta y_B}{\alpha+\beta}\right)$.
    \end{itemize}
    \item \textbf{Construction géométrique à la règle et au compas (Méthode de Thalès sur une droite auxiliaire) :}
    \begin{enumerate}
        \item Tracer la droite $(AB)$.
        \item Tracer une demi-droite $[Ax)$ quelconque, non confondue avec $(AB)$ (souvent perpendiculaire à $(AB)$ en $A$ pour la propreté).
        \item Sur la droite $(Ax)$ (munie d'une orientation), reporter les longueurs \textbf{signées} :
        \begin{itemize}
            \item $A'$ tel que $\overrightarrow{AA'} = \beta \cdot \vec{u}$ (coefficient du \textbf{deuxième} point $B$).
            \item $B'$ tel que $\overrightarrow{AB'} = (\alpha+\beta) \cdot \vec{u}$ (somme des coefficients).
        \end{itemize}
        \emph{Note :} On utilise une même unité $u$. Le signe de $\beta$ et $\alpha+\beta$ donne le sens sur $(Ax)$.
        \item Tracer la droite $(B'B)$.
        \item Par $A'$, tracer la parallèle à $(B'B)$. Elle coupe $(AB)$ en $G$.
    \end{enumerate}
    \item \textbf{Justification (Thalès) :} Dans le triangle $ABB'$, $A' \in (AB')$, $G \in (AB)$ et $(A'G) \parallel (B'B)$.
    \[
    \frac{AG}{AB} = \frac{AA'}{AB'} = \frac{\beta}{\alpha+\beta} \quad \Rightarrow \quad \overrightarrow{AG} = \frac{\beta}{\alpha+\beta}\overrightarrow{AB}
    \]
    Ce qui caractérise bien le barycentre.
    \item \textbf{Vérifications systématiques :} Alignement $A,G,B$ ; Position de $G$ (intérieur/extérieur) cohérente avec les signes ; Rapport $\frac{GA}{GB} = \left|\frac{\beta}{\alpha}\right|$.
\end{enumerate}
\end{methode}

\begin{exemplebox}[Construction : cas mêmes signes et signes contraires]
\begin{minipage}{0.48\linewidth}
\textbf{Cas $\alpha=3, \beta=2$ (mêmes signes)} \\
$G \in [AB]$. \\
Sur $[Ax)$ : $AA' = 2u$, $AB' = 5u$ (même sens). \\
Parallèle à $(B'B)$ par $A'$ $\Rightarrow G \in [AB]$.
\end{minipage}
\hfill
\begin{minipage}{0.48\linewidth}
\textbf{Cas $\alpha=3, \beta=-2$ (signes contraires)} \\
$\alpha+\beta=1 > 0$. $G$ extérieur, côté $A$ ($|\alpha|>|\beta|$). \\
Sur $(Ax)$ : $AA' = -2u$ (sens opposé), $AB' = 1u$ (sens direct). \\
Parallèle à $(B'B)$ par $A'$ $\Rightarrow G \notin [AB]$, côté $A$.
\end{minipage}
\end{exemplebox}

\begin{exoresolu}[Barycentre de deux points — Coordonnées et construction]
Soient $A(1;2)$ et $B(5;-1)$ dans un repère orthonormé. On considère les points pondérés $(A; 3)$ et $(B; -2)$.
\begin{enumerate}
    \item Déterminer les coordonnées du barycentre $G$ de ce système.
    \item Construire $G$ à la règle et au compas (décrire la construction).
    \item Vérifier que $G$ est aligné avec $A$ et $B$ et donner le rapport $\frac{GA}{GB}$.
\end{enumerate}
\tcblower
\textbf{Solution détaillée}
\begin{enumerate}
    \item \textbf{Coordonnées de $G$ :}
    Somme des coefficients : $S = \alpha+\beta = 3 + (-2) = 1 \neq 0$. Le barycentre existe.
    \[
    x_G = \frac{3 \times 1 + (-2) \times 5}{1} = \frac{3 - 10}{1} = -7
    \]
    \[
    y_G = \frac{3 \times 2 + (-2) \times (-1)}{1} = \frac{6 + 2}{1} = 8
    \]
    Donc \textbf{$G(-7; 8)$}.
    
    \item \textbf{Construction géométrique (Méthode correcte) :}
    $\alpha=3, \beta=-2 \Rightarrow \alpha+\beta=1$.
    \begin{enumerate}
        \item Tracer la droite $(AB)$.
        \item Tracer une demi-droite $[Ax)$ (ex: perpendiculaire à $(AB)$ en $A$), orientée par un vecteur unitaire $\vec{u}$.
        \item Construire $A'$ sur $(Ax)$ tel que $\overrightarrow{AA'} = \beta \vec{u} = -2\vec{u}$ (2 unités en sens inverse de $\vec{u}$).
        \item Construire $B'$ sur $(Ax)$ tel que $\overrightarrow{AB'} = (\alpha+\beta)\vec{u} = 1\vec{u}$ (1 unité dans le sens de $\vec{u}$).
        \item Tracer la droite $(B'B)$.
        \item Par $A'$, tracer la parallèle à $(B'B)$. Elle coupe $(AB)$ en $G$.
    \end{enumerate}
    \emph{Justification :} Thalès dans $\triangle ABB'$ : $\frac{AG}{AB} = \frac{AA'}{AB'} = \frac{-2}{1} = -2 \Rightarrow \overrightarrow{AG} = -2\overrightarrow{AB}$. $G$ est bien sur $(AB)$, du côté de $A$ (extérieur), et $AG = 2 AB$.
    
    \item \textbf{Vérification alignement et rapport :}
    $\overrightarrow{AG} = \begin{pmatrix} -8 \\ 6 \end{pmatrix}$, $\overrightarrow{AB} = \begin{pmatrix} 4 \\ -3 \end{pmatrix}$.
    $\overrightarrow{AG} = -2 \overrightarrow{AB}$ $\Rightarrow$ vecteurs colinéaires $\Rightarrow$ \textbf{$A, G, B$ alignés}.
    $\overrightarrow{BG} = \overrightarrow{BA} + \overrightarrow{AG} = \begin{pmatrix} -4 \\ 3 \end{pmatrix} + \begin{pmatrix} -8 \\ 6 \end{pmatrix} = \begin{pmatrix} -12 \\ 9 \end{pmatrix} = 3 \begin{pmatrix} -4 \\ 3 \end{pmatrix}$.
    $GA = \|\overrightarrow{AG}\| = 2\|\overrightarrow{AB}\|$, $GB = \|\overrightarrow{BG}\| = 3\|\overrightarrow{AB}\|$.
    $\frac{GA}{GB} = \frac{2}{3} = \left|\frac{\beta}{\alpha}\right| = \left|\frac{-2}{3}\right|$. \textbf{Correct.}
    \textbf{Conclusion :} $G(-7;8)$ est le barycentre, aligné à $A$ et $B$, extérieur à $[AB]$ côté $A$, $\frac{GA}{GB} = \frac{2}{3}$.
\end{enumerate}
\end{exoresolu}

\begin{popfigure}
\begin{tikzpicture}[scale=0.8, >=Stealth]
% Points
\coordinate (A) at (0,0);
\coordinate (B) at (4,0);
% Auxiliary line
\coordinate (Ax) at (0,2);
\draw[dashed, gray] (A) -- (Ax) node[left] {$Ax$};
% Construction for alpha=3, beta=-2 (sum=1)
% Unit u = 0.5cm
\coordinate (Aprime) at (0,-1); % AA' = -2u
\coordinate (Bprime) at (0,0.5); % AB' = 1u
\draw[popblue, thick] (A) -- (Aprime) node[left] {$A'$};
\draw[poporange, thick] (A) -- (Bprime) node[left] {$B'$};
% Line B'B
\draw[popgreen] (Bprime) -- (B) node[right] {$B$};
% Parallel through A'
\draw[poppink, dashed] (Aprime) -- ++(4.5, -1.5) coordinate (G) node[below right] {$G$};
% Line AB
\draw[thick] (A) node[left] {$A$} -- (B);
% Mark G intersection
\fill (G) circle (2pt);
% Labels
\node at (2, -0.5) {$\alpha=3,\ \beta=-2$};
\node at (2, -1) {$S=1$};
\end{tikzpicture}
\legende{Construction du barycentre pour $\alpha=3, \beta=-2$ (coefficients de signes contraires).}
\end{popfigure}

\courssub{Barycentre de trois ou quatre points pondérés : barycentres partiels, alignements et concurrence}

\begin{definition}[Barycentre de $n$ points pondérés]
Soient $A_1, \dots, A_n$ des points du plan et $\alpha_1, \dots, \alpha_n$ des réels tels que $\sum_{i=1}^n \alpha_i \neq 0$.
Le \textbf{barycentre} $G$ du système $\{(A_i;\alpha_i)\}_{i=1}^n$ est l'unique point vérifiant :
\[
\sum_{i=1}^n \alpha_i \overrightarrow{GA_i} = \vec{0}
\quad \text{ou} \quad
\forall M,\ \overrightarrow{MG} = \frac{\sum \alpha_i \overrightarrow{MA_i}}{\sum \alpha_i}
\]
En coordonnées : $G\left(\frac{\sum \alpha_i x_i}{\sum \alpha_i}\ ;\ \frac{\sum \alpha_i y_i}{\sum \alpha_i}\right)$.
\end{definition}

\begin{propriete}[Barycentres partiels (Associativité)]
Soit $G$ le barycentre de $\{(A_i;\alpha_i)\}_{i=1}^n$ ($S=\sum\alpha_i\neq 0$).
Pour toute partition de l'ensemble des indices en sous-ensembles $I_1, \dots, I_k$, si l'on note $G_j$ le barycentre de $\{(A_i;\alpha_i)\}_{i\in I_j}$ (somme $S_j\neq 0$), alors $G$ est le barycentre de $\{(G_j; S_j)\}_{j=1}^k$.
En particulier, on peut regrouper les points deux à deux successivement.
\end{propriete}

\begin{methode}[Déterminer et construire le barycentre de trois points pondérés]
\textbf{Objectif :} Trouver $G$ barycentre de $(A;\alpha), (B;\beta), (C;\gamma)$ avec $S=\alpha+\beta+\gamma \neq 0$.
\begin{enumerate}
    \item \textbf{Méthode des barycentres partiels (Calcul et Construction) :}
    \begin{itemize}
        \item Construire/Calculer $G_1$ barycentre de $(A;\alpha)$ et $(B;\beta)$ (somme $\alpha+\beta$).
        \item $G$ est alors le barycentre de $(G_1; \alpha+\beta)$ et $(C;\gamma)$.
        \item \emph{Ordre arbitraire :} On peut commencer par n'importe quelle paire.
    \end{itemize}
    \item \textbf{Coordonnées (repère orthonormé) :}
    $G\left(\dfrac{\alpha x_A + \beta x_B + \gamma x_C}{S}\ ;\ \dfrac{\alpha y_A + \beta y_B + \gamma y_C}{S}\right)$.
    \item \textbf{Montrer que trois points sont alignés (Critère barycentrique) :}
    Si un point $P$ est barycentre de deux points $Q$ et $R$ avec coefficients de somme non nulle, alors $P, Q, R$ sont alignés.
    \item \textbf{Montrer que trois droites sont concourantes (Théorème de Céva vectoriel / Barycentres partiels) :}
    Dans un triangle $ABC$, si l'on définit des points $J\in (BC), K\in (CA), L\in (AB)$ comme barycentres partiels des sommets, alors les droites $(AJ), (BK), (CL)$ sont concourantes en $G$, barycentre total du système.
\end{enumerate}
\end{methode}

\begin{exoresolu}[Barycentre de trois points — Alignement et concurrence (Céva)]
Soit $ABC$ un triangle non isocèle. On note $G$ le barycentre des points pondérés $(A; 2), (B; 3), (C; 4)$.
$J$ est le barycentre de $(B; 3)$ et $(C; 4)$.
$K$ est le barycentre de $(A; 2)$ et $(C; 4)$.
$L$ est le barycentre de $(A; 2)$ et $(B; 3)$.
\begin{enumerate}
    \item Déterminer les coordonnées de $G$ dans le repère $(A; \overrightarrow{AB}, \overrightarrow{AC})$.
    \item Montrer que les points $A, J, G$ sont alignés.
    \item Montrer que les droites $(AJ)$, $(BK)$ et $(CL)$ sont concourantes. Préciser le point de concours.
\end{enumerate}
\tcblower
\textbf{Solution détaillée}
\begin{enumerate}
    \item \textbf{Coordonnées de $G$ dans $(A; \vec{AB}, \vec{AC})$ :}
    Dans ce repère affine : $A(0,0)$, $B(1,0)$, $C(0,1)$.
    Somme des coefficients : $S = 2+3+4 = 9 \neq 0$.
    \[
    x_G = \frac{2(0) + 3(1) + 4(0)}{9} = \frac{3}{9} = \frac{1}{3}
    \]
    \[
    y_G = \frac{2(0) + 3(0) + 4(1)}{9} = \frac{4}{9}
    \]
    \textbf{$G\left(\frac{1}{3}; \frac{4}{9}\right)$}.
    
    \item \textbf{Alignement $A, J, G$ :}
    $J$ est barycentre de $(B;3)$ et $(C;4)$. Somme $S_{BC} = 7 \neq 0$.
    $\overrightarrow{AJ} = \frac{3\overrightarrow{AB} + 4\overrightarrow{AC}}{7}$.
    $G$ est barycentre du système global. Par associativité, $G$ est barycentre de $(A;2)$ et $(J; 7)$ (puisque $2+3+4 = 2+7$).
    Par définition du barycentre de deux points, $G$ appartient à la droite $(AJ)$.
    \textbf{Conclusion : $A, J, G$ sont alignés.}
    
    \item \textbf{Concurrence de $(AJ), (BK), (CL)$ :}
    \begin{itemize}
        \item $J$ barycentre de $(B;3),(C;4) \Rightarrow G$ barycentre de $(A;2)$ et $(J;7) \Rightarrow G \in (AJ)$.
        \item $K$ barycentre de $(A;2),(C;4) \Rightarrow$ somme $6$. $G$ barycentre de $(B;3)$ et $(K;6) \Rightarrow G \in (BK)$.
        \item $L$ barycentre de $(A;2),(B;3) \Rightarrow$ somme $5$. $G$ barycentre de $(C;4)$ et $(L;5) \Rightarrow G \in (CL)$.
    \end{itemize}
    Les trois droites $(AJ), (BK), (CL)$ passent par $G$. Elles sont \textbf{concourantes en $G$}.
    \emph{Remarque :} C'est le théorème de Céva sous forme barycentrique. Les coefficients $(2,3,4)$ définissent un point $G$ général (ni centre de gravité, ni centre de masse pondéré par les longueurs des côtés).
\end{enumerate}
\end{exoresolu}

\begin{popfigure}
\begin{tikzpicture}[scale=1.2, >=Stealth]
% Triangle ABC
\coordinate (A) at (0,0);
\coordinate (B) at (4,0);
\coordinate (C) at (1,3);
\draw[thick] (A) -- (B) -- (C) -- cycle;
\node[left] at (A) {$A$};
\node[right] at (B) {$B$};
\node[above] at (C) {$C$};
% Partial barycenters
% J on BC: (3B+4C)/7
\coordinate (J) at ($(B)!3/7!(C)$); % Wait, barycenter (B,3),(C,4) -> J = (3B+4C)/7. Vector BJ = 4/7 BC. So J is at 4/7 from B to C? No.
% G = (alpha A + beta B)/(alpha+beta). AG/AB = beta/(alpha+beta).
% J barycenter (B,3), (C,4). BJ/JC = 4/3 (signed). BJ/BC = 4/7.
\coordinate (J) at ($(B)!4/7!(C)$);
% K on AC: (A,2),(C,4). AK/KC = 4/2 = 2. AK/AC = 4/6 = 2/3.
\coordinate (K) at ($(A)!2/3!(C)$);
% L on AB: (A,2),(B,3). AL/LB = 3/2. AL/AB = 3/5.
\coordinate (L) at ($(A)!3/5!(B)$);
% Draw cevians
\draw[popblue, thick] (A) -- (J) node[midway, above left] {$AJ$};
\draw[poporange, thick] (B) -- (K) node[midway, above right] {$BK$};
\draw[popgreen, thick] (C) -- (L) node[midway, left] {$CL$};
% G intersection
\coordinate (G) at (intersection of A--J and B--K);
\fill[poppink] (G) circle (2.5pt) node[below right] {$G$};
% Labels on sides
\node[below] at (L) {$L$};
\node[right] at (J) {$J$};
\node[left] at (K) {$K$};
% Coefficients
\node[popmuted, font=\small] at (2, -0.5) {Coeff : $A(2), B(3), C(4)$};
\end{tikzpicture}
\legende{Concurrence des droites joignant les sommets aux barycentres partiels des côtés opposés (Céva barycentrique).}
\end{popfigure}

\courssub{Réduction de la somme vectorielle $\sum \alpha_i \overrightarrow{MA_i}$}

\begin{propriete}[Réduction de $\vec{S}(M) = \sum_{i=1}^n \alpha_i \overrightarrow{MA_i}$]
Soient $A_1, \dots, A_n$ des points du plan et $\alpha_1, \dots, \alpha_n \in \mathbb{R}$. On pose $S = \sum_{i=1}^n \alpha_i$.
\begin{enumerate}
    \item \textbf{Cas $S \neq 0$ (Barycentre $G$ existe) :}
    \[
    \sum_{i=1}^n \alpha_i \overrightarrow{MA_i} = S \cdot \overrightarrow{MG}
    \]
    où $G$ est le barycentre de $\{(A_i;\alpha_i)\}$.
    \emph{Conséquence :} L'ensemble $\{M \mid \sum \alpha_i \overrightarrow{MA_i} = \vec{0}\}$ est le singleton $\{G\}$.
    \item \textbf{Cas $S = 0$ (Pas de barycentre / Vecteur libre) :}
    L'expression $\sum \alpha_i \overrightarrow{MA_i}$ \textbf{ne dépend pas de $M$}. C'est un vecteur constant $\vec{v}_0$.
    Pour le calculer, on choisit un $M$ commode (souvent $M = A_1$).
    \emph{Conséquence :} L'ensemble $\{M \mid \sum \alpha_i \overrightarrow{MA_i} = \vec{v}\}$ est :
    \begin{itemize}
        \item Le plan entier si $\vec{v} = \vec{v}_0$.
        \item L'ensemble vide $\varnothing$ sinon.
    \end{itemize}
\end{enumerate}
\end{propriete}

\begin{methode}[Réduire $\sum \alpha_i \overrightarrow{MA_i}$ et résoudre $\sum \alpha_i \overrightarrow{MA_i} = \vec{v}$]
\begin{enumerate}
    \item Calculer $S = \sum \alpha_i$.
    \item \textbf{Si $S \neq 0$ :} Déterminer $G$ barycentre de $(A_i;\alpha_i)$. Écrire $\sum \alpha_i \overrightarrow{MA_i} = S \overrightarrow{MG}$.
    \begin{itemize}
        \item Équation $\sum \alpha_i \overrightarrow{MA_i} = \vec{0} \iff M=G$.
        \item Équation $\sum \alpha_i \overrightarrow{MA_i} = \vec{v} \iff \overrightarrow{MG} = \frac{\vec{v}}{S}$ (un unique point $M$).
    \end{itemize}
    \item \textbf{Si $S = 0$ :} Calculer le vecteur constant $\vec{v}_0 = \sum \alpha_i \overrightarrow{M_0A_i}$ (avec $M_0$ commode, ex $A_1$).
    \begin{itemize}
        \item Si $\vec{v} = \vec{v}_0$ : Ensemble solution = Plan entier.
        \item Si $\vec{v} \neq \vec{v}_0$ : Ensemble solution = $\varnothing$.
    \end{itemize}
\end{enumerate}
\end{methode}

\begin{exoresolu}[Réduction de somme vectorielle et ensembles solutions]
Soient $A, B, C$ trois points non alignés du plan. On considère :
$\vec{E}(M) = 3\overrightarrow{MA} - 2\overrightarrow{MB} + \overrightarrow{MC}$.
\begin{enumerate}
    \item Réduire $\vec{E}(M)$ en fonction de $\overrightarrow{MG}$ où $G$ est un point à déterminer (préciser ses coefficients).
    \item Déterminer l'ensemble $\mathcal{E}$ des points $M$ tels que $\vec{E}(M) = \vec{0}$.
    \item On pose $\vec{F}(M) = 2\overrightarrow{MA} - 3\overrightarrow{MB} + \overrightarrow{MC}$. Réduire $\vec{F}(M)$. Déterminer l'ensemble $\mathcal{F}$ des points $M$ tels que $\vec{F}(M) = \overrightarrow{AB}$.
\end{enumerate}
\tcblower
\textbf{Solution détaillée}
\begin{enumerate}
    \item \textbf{Réduction de $\vec{E}(M)$ :}
    Coefficients : $\alpha_A = 3, \alpha_B = -2, \alpha_C = 1$.
    Somme $S = 3 - 2 + 1 = 2 \neq 0$.
    Le barycentre $G$ existe. C'est le barycentre de $(A;3), (B;-2), (C;1)$.
    \textbf{Propriété :} $\vec{E}(M) = S \cdot \overrightarrow{MG} = 2 \overrightarrow{MG}$.
    
    \item \textbf{Ensemble $\mathcal{E} = \{ M \mid \vec{E}(M) = \vec{0} \}$ :}
    $2 \overrightarrow{MG} = \vec{0} \iff \overrightarrow{MG} = \vec{0} \iff M = G$.
    $\mathcal{E} = \{ G \}$ (singleton).
    
    \item \textbf{Réduction de $\vec{F}(M)$ :}
    Coefficients : $2, -3, 1$. Somme $S' = 2 - 3 + 1 = 0$.
    \textbf{Propriété :} $\vec{F}(M)$ ne dépend pas de $M$. C'est un vecteur constant $\vec{v}_0$.
    Calculons $\vec{v}_0$ en prenant $M = A$ (commode car $\overrightarrow{AA}=\vec{0}$) :
    \[
    \vec{v}_0 = \vec{F}(A) = 2\vec{0} - 3\overrightarrow{AB} + \overrightarrow{AC} = -3\overrightarrow{AB} + \overrightarrow{AC}
    \]
    \textbf{Ensemble $\mathcal{F} = \{ M \mid \vec{F}(M) = \overrightarrow{AB} \}$ :}
    L'équation devient $\vec{v}_0 = \overrightarrow{AB}$.
    $-3\overrightarrow{AB} + \overrightarrow{AC} = \overrightarrow{AB} \iff \overrightarrow{AC} = 4\overrightarrow{AB}$.
    Or $A, B, C$ sont \textbf{non alignés} (données), donc $\overrightarrow{AC}$ n'est pas colinéaire à $\overrightarrow{AB}$.
    L'égalité est \textbf{fausse}.
    \textbf{Conclusion : $\mathcal{F} = \varnothing$ (ensemble vide).}
\end{enumerate}
\end{exoresolu}

\courssub{Lignes de niveau : $MA^2 + MB^2 = k$ (Cercles de centre le milieu de $AB$)}

\begin{propriete}[Relation de Leibniz / Médiane]
Soit $I$ le milieu de $[AB]$. Pour tout point $M$ du plan :
\[
MA^2 + MB^2 = 2 MI^2 + \frac{AB^2}{2}
\]
\emph{Preuve :} $\overrightarrow{MA} = \overrightarrow{MI} + \overrightarrow{IA}$, $\overrightarrow{MB} = \overrightarrow{MI} + \overrightarrow{IB} = \overrightarrow{MI} - \overrightarrow{IA}$.
$MA^2+MB^2 = (\vec{MI}+\vec{IA})^2 + (\vec{MI}-\vec{IA})^2 = 2MI^2 + 2IA^2 = 2MI^2 + \frac{AB^2}{2}$.
\end{propriete}

\begin{methode}[Ligne de niveau $\mathcal{C} = \{ M \mid MA^2 + MB^2 = k \}$]
\begin{enumerate}
    \item Exprimer la condition : $2 MI^2 + \frac{AB^2}{2} = k \iff MI^2 = \frac{k}{2} - \frac{AB^2}{4}$.
    \item \textbf{Discussion selon $k$ :}
    \begin{itemize}
        \item Si $k < \frac{AB^2}{2}$ : $MI^2 < 0$ impossible $\Rightarrow \mathcal{C} = \varnothing$.
        \item Si $k = \frac{AB^2}{2}$ : $MI^2 = 0 \Rightarrow M=I$. $\mathcal{C} = \{I\}$.
        \item Si $k > \frac{AB^2}{2}$ : $\mathcal{C}$ est le \textbf{cercle de centre $I$ (milieu de $AB$) et de rayon $R = \sqrt{\frac{k}{2} - \frac{AB^2}{4}}$}.
    \end{itemize}
    \item \textbf{Construction :} Tracer $I$ milieu de $[AB]$. Calculer $R$. Tracer le cercle $\mathcal{C}(I, R)$.
\end{enumerate}
\end{methode}

\begin{exoresolu}[Ligne de niveau $MA^2 + MB^2 = k$ — Application]
Dans un repère orthonormé, $A(-3;0)$ et $B(3;0)$. On cherche l'ensemble $\mathcal{E}$ des points $M(x;y)$ tels que $MA^2 + MB^2 = 40$.
\begin{enumerate}
    \item Déterminer la nature de $\mathcal{E}$, son centre et son rayon.
    \item Donner l'équation cartésienne de $\mathcal{E}$.
    \item Vérifier si le point $M_0(0; \sqrt{11})$ appartient à $\mathcal{E}$.
\end{enumerate}
\tcblower
\textbf{Solution détaillée}
\begin{enumerate}
    \item \textbf{Nature, centre, rayon :}
    $AB = 6 \Rightarrow AB^2 = 36$.
    $I$ milieu de $[AB]$ : $I(0;0)$ (origine).
    Seuil : $\frac{AB^2}{2} = 18$.
    Donné $k = 40$. Comme $40 > 18$, $\mathcal{E}$ est un \textbf{cercle}.
    Centre : $I(0;0)$.
    Rayon $R$ : $R^2 = \frac{k}{2} - \frac{AB^2}{4} = \frac{40}{2} - \frac{36}{4} = 20 - 9 = 11$.
    $R = \sqrt{11}$.
    
    \item \textbf{Équation cartésienne :}
    Cercle centre $O(0,0)$ rayon $\sqrt{11}$ : $x^2 + y^2 = 11$.
    \emph{Vérification par le calcul direct :}
    $M(x,y)$. $MA^2 = (x+3)^2 + y^2 = x^2+6x+9+y^2$.
    $MB^2 = (x-3)^2 + y^2 = x^2-6x+9+y^2$.
    Somme $= 2x^2 + 2y^2 + 18 = 40 \iff 2(x^2+y^2) = 22 \iff x^2+y^2 = 11$. \textbf{Ok.}
    
    \item \textbf{Appartenance de $M_0(0; \sqrt{11})$ :}
    $x^2+y^2 = 0^2 + (\sqrt{11})^2 = 11$.
    L'équation est vérifiée. \textbf{$M_0 \in \mathcal{E}$}.
\end{enumerate}
\end{exoresolu}

\begin{popfigure}
\begin{tikzpicture}[scale=0.7, >=Stealth]
\begin{axis}[
    axis equal, axis lines=middle,
    xmin=-5, xmax=5, ymin=-4, ymax=4,
    xtick={-3,0,3}, ytick={-3,0,3},
    xlabel=$x$, ylabel=$y$,
    width=0.9\linewidth, height=6cm,
    grid=major, grid style={popline},
]
% Points A, B
\addplot[only marks, mark=*, mark size=2.5, popdark] coordinates {(-3,0) (3,0)};
\node[below left] at (axis cs:-3,0) {$A(-3,0)$};
\node[below right] at (axis cs:3,0) {$B(3,0)$};
% Circle
\addplot[popblue, thick, domain=0:360, samples=100, variable=t] ({sqrt(max(0,11))*cos(t)}, {sqrt(max(0,11))*sin(t)});
% Point M0
\addplot[only marks, mark=*, mark size=2.5, poppink] coordinates {(0,3.3166)};
\node[above left] at (axis cs:0,3.3166) {$M_0(0,\sqrt{11})$};
% Center I
\addplot[only marks, mark=*, mark size=2, poporange] coordinates {(0,0)};
\node[below right] at (axis cs:0,0) {$I(0,0)$};
\end{axis}
\end{tikzpicture}
\legende{Ligne de niveau $MA^2+MB^2=40$ : cercle centre $I$ (milieu de $AB$) rayon $\sqrt{11}$.}
\end{popfigure}

\courssub{Lignes de niveau : $MA^2 - MB^2 = k$ (Droites perpendiculaires à $AB$)}

\begin{propriete}[Différence des carrés des distances]
Soit $I$ le milieu de $[AB]$. Pour tout point $M$ :
\[
MA^2 - MB^2 = \overrightarrow{MA}^2 - \overrightarrow{MB}^2 = (\overrightarrow{MA} - \overrightarrow{MB})\cdot(\overrightarrow{MA} + \overrightarrow{MB}) = \overrightarrow{BA} \cdot (2\overrightarrow{MI}) = 2 \overrightarrow{AB} \cdot \overrightarrow{IM}
\]
Si $H$ est le projeté orthogonal de $M$ sur la droite $(AB)$, alors $\overrightarrow{AB} \cdot \overrightarrow{IM} = \overrightarrow{AB} \cdot \overrightarrow{IH} = AB \times IH$ (allongement signé).
Donc $MA^2 - MB^2 = 2 \cdot AB \cdot IH$.
\end{propriete}

\begin{methode}[Ligne de niveau $\mathcal{D} = \{ M \mid MA^2 - MB^2 = k \}$]
\begin{enumerate}
    \item L'équation équivaut à $2 \cdot AB \cdot IH = k \iff IH = \frac{k}{2 AB}$ (allongement signé sur $(AB)$).
    \item \textbf{Nature :} C'est la \textbf{droite perpendiculaire à $(AB)$} passant par le point $H$ de la droite $(AB)$ tel que $IH = \frac{k}{2 AB}$.
    \begin{itemize}
        \item Si $k=0$ : $H=I$, c'est la \textbf{médiatrice} de $[AB]$.
        \item Si $k>0$ : $IH > 0$, $H$ est du côté de $B$ (même sens que $\overrightarrow{IB}$).
        \item Si $k<0$ : $IH < 0$, $H$ est du côté de $A$.
    \end{itemize}
    \item \textbf{Équation cartésienne (Repère $A(-c,0), B(c,0)$) :} $MA^2-MB^2 = 4cx = k \Rightarrow x = \frac{k}{4c}$. Droite verticale.
\end{enumerate}
\end{methode}

\begin{exoresolu}[Ligne de niveau $MA^2 - MB^2 = k$ — Droite perpendiculaire]
Soient $A(-2;0)$ et $B(4;0)$. Déterminer l'ensemble $\mathcal{L}$ des points $M$ tels que $MA^2 - MB^2 = 12$.
\begin{enumerate}
    \item Montrer que $\mathcal{L}$ est une droite. En donner une équation cartésienne.
    \item Construire cette droite à la règle et au compas.
    \item Le point $N(1; 3)$ appartient-il à $\mathcal{L}$ ?
\end{enumerate}
\tcblower
\textbf{Solution détaillée}
\begin{enumerate}
    \item \textbf{Nature et équation :}
    $A(-2,0), B(4,0) \Rightarrow AB = 6$. Milieu $I(1,0)$.
    Méthode 1 (Produit scalaire) :
    $MA^2 - MB^2 = 2 \overrightarrow{AB} \cdot \overrightarrow{IM}$.
    $\overrightarrow{AB} = (6, 0)$. $\overrightarrow{IM} = (x-1, y)$.
    $2 \times (6,0) \cdot (x-1, y) = 12(x-1)$.
    Équation : $12(x-1) = 12 \iff x-1 = 1 \iff \textbf{x = 2}$.
    C'est une \textbf{droite verticale} (perpendiculaire à $(AB)$ horizontale), passant par $H(2,0)$.
    Vérification $IH = 1$. $\frac{k}{2AB} = \frac{12}{12} = 1$. $I(1,0) \to H(2,0)$ (vers $B$ car $k>0$). Correct.
    
    Méthode 2 (Calcul direct) :
    $MA^2 = (x+2)^2+y^2$, $MB^2 = (x-4)^2+y^2$.
    Différence $= (x^2+4x+4) - (x^2-8x+16) = 12x - 12 = 12 \Rightarrow 12x = 24 \Rightarrow x=2$.
    
    \item \textbf{Construction :}
    \begin{enumerate}
        \item Tracer le segment $[AB]$ et son milieu $I$.
        \item Calculer $IH = \frac{k}{2AB} = \frac{12}{12} = 1$ (unité du repère).
        \item Sur la droite $(AB)$, à partir de $I$, reporter la longueur $1$ vers $B$ (car $k>0$) pour obtenir $H$.
        \item Tracer la perpendiculaire à $(AB)$ en $H$. C'est la droite $\mathcal{L}$.
    \end{enumerate}
    
    \item \textbf{Appartenance de $N(1;3)$ :}
    $N$ a pour abscisse $x_N = 1$.
    L'équation de $\mathcal{L}$ est $x = 2$.
    $1 \neq 2 \Rightarrow \textbf{$N \notin \mathcal{L}$}$.
    \emph{Vérification numérique :} $NA^2 = (3)^2+3^2=18$, $NB^2 = (-3)^2+3^2=18$. $NA^2-NB^2=0 \neq 12$.
\end{enumerate}
\end{exoresolu}

\begin{popfigure}
\begin{tikzpicture}[scale=0.7, >=Stealth]
\begin{axis}[
    axis equal, axis lines=middle,
    xmin=-4, xmax=6, ymin=-4, ymax=4,
    xtick={-2,1,2,4}, ytick={-3,0,3},
    xlabel=$x$, ylabel=$y$,
    width=0.9\linewidth, height=6cm,
    grid=major, grid style={popline},
]
% Points A, B, I
\addplot[only marks, mark=*, mark size=2.5, popdark] coordinates {(-2,0) (4,0) (1,0)};
\node[below] at (axis cs:-2,0) {$A$};
\node[below] at (axis cs:4,0) {$B$};
\node[below] at (axis cs:1,0) {$I$};
% Line L: x=2
\addplot[popblue, thick, domain=-4:4] ({2}, {x});
\node[above right] at (axis cs:2,3) {$\mathcal{L}: x=2$};
% Point H
\addplot[only marks, mark=*, mark size=2, poporange] coordinates {(2,0)};
\node[below] at (axis cs:2,0) {$H$};
% Point N
\addplot[only marks, mark=*, mark size=2.5, poppink] coordinates {(1,3)};
\node[above left] at (axis cs:1,3) {$N(1,3)$};
% Segment IH
\draw[popgreen, thick, <->] (axis cs:1, -0.3) -- (axis cs:2, -0.3) node[midway, below] {$IH=1$};
\end{axis}
\end{tikzpicture}
\legende{Ligne de niveau $MA^2-MB^2=12$ : droite perpendiculaire à $AB$ en $H$ tel que $IH = k/(2AB)$.}
\end{popfigure}

\courssub{Lignes de niveau : $\overrightarrow{MA} \cdot \overrightarrow{MB} = k$ (Cercles de centre le milieu de $AB$)}

\begin{propriete}[Produit scalaire et milieu]
Soit $I$ le milieu de $[AB]$. Pour tout point $M$ :
\[
\overrightarrow{MA} \cdot \overrightarrow{MB} = (\overrightarrow{MI} + \overrightarrow{IA}) \cdot (\overrightarrow{MI} + \overrightarrow{IB}) = (\overrightarrow{MI} + \overrightarrow{IA}) \cdot (\overrightarrow{MI} - \overrightarrow{IA}) = MI^2 - IA^2 = MI^2 - \frac{AB^2}{4}
\]
\end{propriete}

\begin{methode}[Ligne de niveau $\mathcal{C} = \{ M \mid \overrightarrow{MA} \cdot \overrightarrow{MB} = k \}$]
\begin{enumerate}
    \item Équation réduite : $MI^2 = k + \frac{AB^2}{4}$.
    \item \textbf{Discussion selon $k$ :}
    \begin{itemize}
        \item Si $k < -\frac{AB^2}{4}$ : $\mathcal{C} = \varnothing$.
        \item Si $k = -\frac{AB^2}{4}$ : $\mathcal{C} = \{ I \}$ (point unique).
        \item Si $k > -\frac{AB^2}{4}$ : $\mathcal{C}$ est le \textbf{cercle de centre $I$ (milieu de $AB$) et de rayon $R = \sqrt{k + \frac{AB^2}{4}}$}.
    \end{itemize}
    \item \textbf{Cas particulier $k=0$ :} $R = \frac{AB}{2}$. C'est le \textbf{cercle de diamètre $[AB]$}.
    $\overrightarrow{MA} \cdot \overrightarrow{MB} = 0 \iff (MA) \perp (MB) \iff M$ voit $[AB]$ sous un angle droit.
\end{enumerate}
\end{methode}

\begin{exoresolu}[Produit scalaire $\overrightarrow{MA} \cdot \overrightarrow{MB} = k$ — Cercle de diamètre]
Soient $A(0;0)$ et $B(8;0)$. On considère l'ensemble $\mathcal{C}_k = \{ M \mid \overrightarrow{MA} \cdot \overrightarrow{MB} = k \}$.
\begin{enumerate}
    \item Pour quelles valeurs de $k$ l'ensemble $\mathcal{C}_k$ est-il un cercle ? Un point ? Vide ?
    \item Déterminer l'équation cartésienne de $\mathcal{C}_7$.
    \item Montrer que pour $k=0$, $\mathcal{C}_0$ est le cercle de diamètre $[AB]$. Donner son centre et son rayon.
\end{enumerate}
\tcblower
\textbf{Solution détaillée}
\begin{enumerate}
    \item \textbf{Analyse selon $k$ :}
    $AB = 8 \Rightarrow AB^2 = 64 \Rightarrow \frac{AB^2}{4} = 16$. Milieu $I(4,0)$.
    Relation : $MI^2 = k + 16$.
    \begin{itemize}
        \item Cercle $\iff k + 16 > 0 \iff \textbf{k > -16}$. Rayon $R = \sqrt{k+16}$.
        \item Point unique ($\{I\}$) $\iff k + 16 = 0 \iff \textbf{k = -16}$.
        \item Vide $\iff k + 16 < 0 \iff \textbf{k < -16}$.
    \end{itemize}
    
    \item \textbf{Équation de $\mathcal{C}_7$ :}
    $k=7 > -16$, c'est un cercle. $R^2 = 7+16 = 23$. Centre $I(4,0)$.
    Équation : $(x-4)^2 + y^2 = 23$.
    \emph{Vérification :} $M(x,y)$. $\overrightarrow{MA}=(-x,-y)$, $\overrightarrow{MB}=(8-x,-y)$.
    Produit scalaire $= -x(8-x) + y^2 = -8x + x^2 + y^2 = 7$.
    $\iff x^2+y^2-8x-7=0 \iff (x-4)^2 - 16 + y^2 - 7 = 0 \iff (x-4)^2+y^2=23$. \textbf{Ok.}
    
    \item \textbf{Cas $k=0$ :}
    $R^2 = 0 + 16 = 16 \Rightarrow R = 4$.
    Centre $I(4,0)$.
    Comme $AB=8$, le rayon $R=4 = \frac{AB}{2}$ et le centre est le milieu $I$.
    Donc $\mathcal{C}_0$ est bien le \textbf{cercle de diamètre $[AB]$}.
    \emph{Propriété géométrique :} $\overrightarrow{MA} \cdot \overrightarrow{MB} = 0 \iff (MA) \perp (MB) \iff M$ appartient au cercle de diamètre $[AB]$ (Théorème de l'angle inscrit / Thalès).
\end{enumerate}
\end{exoresolu}

\begin{popfigure}
\begin{tikzpicture}[scale=0.7, >=Stealth]
\begin{axis}[
    axis equal, axis lines=middle,
    xmin=-2, xmax=10, ymin=-6, ymax=6,
    xtick={0,4,8}, ytick={-4,0,4},
    xlabel=$x$, ylabel=$y$,
    width=0.9\linewidth, height=6cm,
    grid=major, grid style={popline},
]
% Points A, B
\addplot[only marks, mark=*, mark size=2.5, popdark] coordinates {(0,0) (8,0)};
\node[below left] at (axis cs:0,0) {$A$};
\node[below right] at (axis cs:8,0) {$B$};
% Circle k=7 (R=sqrt(23)~4.8)
\addplot[popblue, thick, domain=0:360, samples=100, variable=t] ({4+sqrt(max(0,23))*cos(t)}, {sqrt(max(0,23))*sin(t)});
\node[popblue, above] at (axis cs:4,5) {$\mathcal{C}_7: R=\sqrt{23}$};
% Circle k=0 (Diameter, R=4)
\addplot[poporange, thick, dashed, domain=0:360, samples=100, variable=t] ({4+4*cos(t)}, {4*sin(t)});
\node[poporange, below] at (axis cs:4,-4.5) {$\mathcal{C}_0$: Cercle diam. $[AB], R=4$};
% Center I
\addplot[only marks, mark=*, mark size=2, popgreen] coordinates {(4,0)};
\node[below right] at (axis cs:4,0) {$I$};
\end{axis}
\end{tikzpicture}
\legende{Famille des cercles $\mathcal{C}_k$ pour $A(0,0), B(8,0)$.}
\end{popfigure}

\courssub{Lignes de niveau : $\frac{MA}{MB} = k$ (Cercle d'Apollonius)}

\begin{propriete}[Rapport de distances et cercle d'Apollonius]
Soient $A, B$ deux points distincts et $k > 0, k \neq 1$. L'ensemble $\mathcal{A}_k = \{ M \mid \frac{MA}{MB} = k \}$ est un cercle, appelé \textbf{cercle d'Apollonius}.
Son centre $\Omega$ appartient à la droite $(AB)$, en dehors du segment $[AB]$.
\begin{itemize}
    \item Si $k > 1$ : $\Omega$ est du côté de $B$ (le coefficient plus grand "attire" le centre).
    \item Si $k < 1$ : $\Omega$ est du côté de $A$.
\end{itemize}
\textbf{Formules vectorielles (repère quelconque) :}
\[
\overrightarrow{A\Omega} = \frac{k^2}{k^2-1} \overrightarrow{AB}
\qquad
R = \frac{k}{|k^2-1|} AB
\]
\textbf{Propriété caractéristique :} Pour tout point $M$ sur $(AB)$ appartenant au cercle, $\frac{MA}{MB} = k$. Les intersections du cercle avec $(AB)$ sont les points $P, Q$ tels que $\frac{PA}{PB} = \frac{QA}{QB} = k$. $\Omega$ est le milieu de $[PQ]$ et $\frac{\Omega A}{\Omega B} = k^2$.
\end{propriete}

\begin{methode}[Ligne de niveau $\mathcal{A}_k = \{ M \mid \frac{MA}{MB} = k \}$]
\begin{enumerate}
    \item \textbf{Cas $k = 1$ :} Médiatrice de $[AB]$ (droite).
    \item \textbf{Cas $k > 0, k \neq 1$ :}
    \begin{enumerate}
        \item Calculer le centre $\Omega$ : $\overrightarrow{A\Omega} = \frac{k^2}{k^2-1} \overrightarrow{AB}$.
        \item Calculer le rayon $R = \frac{k}{|k^2-1|} AB$.
        \item L'ensemble est le cercle $\mathcal{C}(\Omega, R)$.
    \end{enumerate}
    \item \textbf{Construction à la règle et au compas (Thalès) :}
    \begin{enumerate}
        \item Tracer la droite $(AB)$.
        \item Construire les deux points $P, Q$ sur $(AB)$ tels que $\frac{PA}{PB} = \frac{QA}{QB} = k$ (un intérieur, un extérieur à $[AB]$).
        \begin{itemize}
            \item Sur une perpendiculaire en $A$, reporter $k$ unités pour $A'$.
            \item Sur la perpendiculaire en $B$ (même côté), reporter $1$ unité pour $B'$.
            \item $(A'B')$ coupe $(AB)$ en $P$ (intérieur).
            \item Pour $Q$ (extérieur), reporter $1$ unité en sens inverse sur la perpendiculaire en $B$ (ou utiliser l'harmonique).
        \end{itemize}
        \item $\Omega$ est le milieu de $[PQ]$. $R = \Omega P = \Omega Q$.
        \item Tracer le cercle centre $\Omega$ rayon $R$.
    \end{enumerate}
\end{enumerate}
\end{methode}

\begin{exoresolu}[Rapport de distances $\frac{MA}{MB} = k$ — Cercle d'Apollonius]
Dans un repère orthonormé, $A(-1;0)$ et $B(3;0)$. On cherche l'ensemble $\mathcal{E}$ des points $M$ tels que $\frac{MA}{MB} = 2$.
\begin{enumerate}
    \item Déterminer le centre $\Omega$ et le rayon $R$ de $\mathcal{E}$.
    \item Donner l'équation cartésienne de $\mathcal{E}$.
    \item Construire $\mathcal{E}$ à la règle et au compas (indiquer les points remarquables sur $(AB)$).
\end{enumerate}
\tcblower
\textbf{Solution détaillée}
\begin{enumerate}
    \item \textbf{Centre et Rayon :}
    $A(-1,0), B(3,0) \Rightarrow AB = 4$. Milieu $I(1,0)$.
    $k = 2 > 1 \Rightarrow$ Centre $\Omega$ du côté de $B$ (au-delà de $B$).
    \textbf{Formule vectorielle :} $\overrightarrow{A\Omega} = \frac{k^2}{k^2-1} \overrightarrow{AB} = \frac{4}{3} \overrightarrow{AB}$.
    $\overrightarrow{AB} = (4, 0) \Rightarrow \overrightarrow{A\Omega} = \frac{4}{3}(4, 0) = (\frac{16}{3}, 0)$.
    $\Omega = A + \overrightarrow{A\Omega} = (-1 + \frac{16}{3}, 0) = (\frac{13}{3}, 0)$.
    \textbf{Rayon :} $R = \frac{k}{|k^2-1|} AB = \frac{2}{3} \times 4 = \frac{8}{3}$.
    \emph{Vérification :} Distance $\Omega B = \frac{13}{3} - 3 = \frac{4}{3}$. $\Omega A = \frac{13}{3} + 1 = \frac{16}{3}$. Rapport $\frac{\Omega A}{\Omega B} = 4 = k^2$. Correct.
    
    \item \textbf{Équation cartésienne :}
    Cercle centre $\Omega(\frac{13}{3}, 0)$, rayon $R = \frac{8}{3}$.
    $(x - \frac{13}{3})^2 + y^2 = (\frac{8}{3})^2 = \frac{64}{9}$.
    Forme développée : $x^2 - \frac{26}{3}x + \frac{169}{9} + y^2 = \frac{64}{9}$
    $\iff x^2 + y^2 - \frac{26}{3}x + \frac{105}{9} = 0 \iff 9x^2 + 9y^2 - 78x + 105 = 0$
    $\iff \boldsymbol{3x^2 + 3y^2 - 26x + 35 = 0}$.
    
    \item \textbf{Construction :}
    \begin{enumerate}
        \item Tracer $[AB]$ et la droite $(AB)$.
        \item Construire les points $P, Q$ sur $(AB)$ tels que $\frac{PA}{PB} = \frac{QA}{QB} = 2$.
        \begin{itemize}
            \item \textbf{Point $P$ (intérieur) :} Sur une perpendiculaire en $A$, reporter $2$ unités $\to A'$. Sur la perpendiculaire en $B$ (même côté), reporter $1$ unité $\to B'$. $(A'B') \cap (AB) = P$.
            \item \textbf{Point $Q$ (extérieur, côté $B$) :} Sur la perpendiculaire en $A$, reporter $2$ unités $\to A'$. Sur la perpendiculaire en $B$ (côté \textbf{opposé}), reporter $1$ unité $\to B''$. $(A'B'') \cap (AB) = Q$.
        \end{itemize}
        \item $\Omega$ est le milieu de $[PQ]$. $R = \Omega P$.
        \item Tracer le cercle centre $\Omega$ rayon $R$.
    \end{enumerate}
    \textbf{Points remarquables sur $(AB)$ (calcul) :}
    Résolution $|x+1| = 2|x-3|$ :
    \begin{itemize}
        \item Cas $x>3$ : $x+1=2x-6 \Rightarrow x=7$ (Point $Q$).
        \item Cas $-1<x<3$ : $x+1=-2x+6 \Rightarrow 3x=5 \Rightarrow x=5/3$ (Point $P$).
    \end{itemize}
    $\Omega$ milieu de $[PQ]$ : $\frac{7 + 5/3}{2} = \frac{26/3}{2} = 13/3$. $R = \frac{7 - 5/3}{2} = \frac{16/3}{2} = 8/3$. \textbf{Ok.}
\end{enumerate}
\end{exoresolu}

\begin{popfigure}
\begin{tikzpicture}[scale=0.7, >=Stealth]
\begin{axis}[
    axis equal, axis lines=middle,
    xmin=-3, xmax=9, ymin=-5, ymax=5,
    xtick={-1,1.666,3,4.333,7}, ytick={-4,0,4},
    xlabel=$x$, ylabel=$y$,
    width=0.9\linewidth, height=6cm,
    grid=major, grid style={popline},
]
% Points A, B
\addplot[only marks, mark=*, mark size=2.5, popdark] coordinates {(-1,0) (3,0)};
\node[below left] at (axis cs:-1,0) {$A$};
\node[below right] at (axis cs:3,0) {$B$};
% Points P, Q
\addplot[only marks, mark=*, mark size=2, popgreen] coordinates {(1.666,0) (7,0)};
\node[below] at (axis cs:1.666,0) {$P$};
\node[below] at (axis cs:7,0) {$Q$};
% Center Omega
\addplot[only marks, mark=*, mark size=2.5, poporange] coordinates {(4.333,0)};
\node[above right] at (axis cs:4.333,0) {$\Omega$};
% Circle
\addplot[popblue, thick, domain=0:360, samples=100, variable=t] ({4.333+2.666*cos(t)}, {2.666*sin(t)});
\node[popblue, above] at (axis cs:4.333,3) {$\mathcal{E}: MA/MB=2$};
% Construction lines for P (Thales)
\coordinate (Ap) at (axis cs:-1,2);
\coordinate (Bp) at (axis cs:3,1);
\draw[popmuted, thin] (Ap) -- (Bp);
% Construction lines for Q (Thales external)
\coordinate (Bpp) at (axis cs:3,-1);
\draw[popmuted, thin] (Ap) -- (Bpp);
\end{axis}
\end{tikzpicture}
\legende{Cercle d'Apollonius pour $k=2$. Centre $\Omega$ milieu de $[PQ]$, $P,Q$ points de rapport $2$ sur $(AB)$.}
\end{popfigure}

\begin{attention}[Les 5 erreurs qui coûtent des points au Probatoire]
\begin{enumerate}
    \item \textbf{Oublier la condition $\sum \alpha_i \neq 0$ pour l'existence du barycentre.}
    Si la somme est nulle, écrire « $G$ est le barycentre... » est une \textbf{erreur grave} (0 point à la question). Il faut écrire : « La somme des coefficients est nulle, le barycentre n'existe pas, le vecteur $\sum \alpha_i \overrightarrow{MA_i}$ est constant. »
    
    \item \textbf{Confondre rapport signé et rapport de longueurs (valeur absolue).}
    $\frac{\overrightarrow{GA}}{\overrightarrow{GB}} = -\frac{\beta}{\alpha}$ (signé, colinéarité).
    $\frac{GA}{GB} = \left| \frac{\beta}{\alpha} \right|$ (longueurs, toujours positif).
    Au Probatoire, préciser « rapport signé » ou « rapport des distances » selon la question.
    
    \item \textbf{Mauvaise réduction de $\sum \alpha_i \overrightarrow{MA_i}$ quand $\sum \alpha_i = 0$.}
    Ne pas écrire « $= 0 \cdot \overrightarrow{MG}$ » (incohérent, $G$ n'existe pas).
    \textbf{Bonne rédaction :} « Comme $\sum \alpha_i = 0$, l'expression ne dépend pas de $M$. En prenant $M=A$, on trouve $\vec{v} = \dots$ »
    
    \item \textbf{Oublier la discussion sur $k$ pour les lignes de niveau (ensemble vide, point, cercle/droite).}
    \begin{itemize}
        \item $MA^2+MB^2=k$ : ne pas dire « c'est un cercle centre $I$ » sans vérifier $k \ge AB^2/2$.
        \item $MA^2-MB^2=k$ : c'est \textbf{toujours une droite} (pas de discussion sur l'existence, mais discussion pour la position de $H$).
        \item $\overrightarrow{MA}\cdot\overrightarrow{MB}=k$ : discuter $k \ge -AB^2/4$.
        \item $MA/MB=k$ : distinguer $k=1$ (droite médiatrice) et $k\neq 1$ (cercle d'Apollonius).
    \end{itemize}
    
    \item \textbf{Erreurs de construction du barycentre (sens des coefficients sur la droite auxiliaire).}
    \begin{itemize}
        \item Méthode standard (utilisée ici) : Sur $(Ax)$, $\overrightarrow{AA'} = \beta \vec{u}$ et $\overrightarrow{AB'} = (\alpha+\beta)\vec{u}$.
        \item Coefficients \textbf{même signe} $\rightarrow$ $A'$ et $B'$ du même côté de $A$ $\rightarrow$ $G \in [AB]$.
        \item Coefficients \textbf{signes contraires} $\rightarrow$ $A'$ et $B'$ de part et d'autre de $A$ (ou sens opposés) $\rightarrow$ $G \notin [AB]$.
    \end{itemize}
    Inverser les sens conduit à un point faux. Vérifier toujours la position de $G$ par rapport à $[AB]$ via le signe de $\frac{\beta}{\alpha}$ (rapport signé).
\end{enumerate}
\end{attention}

\begin{savaistu}[Le barycentre dans la vraie vie : Téléphonie mobile et GPS]
Au Cameroun, les antennes relais \textbf{MTN} et \textbf{Orange} couvrent le territoire. Votre téléphone reçoit des signaux de plusieurs antennes.
\begin{itemize}
    \item \textbf{Localisation par trilatération (GPS) :} Le récepteur calcule sa position comme barycentre des positions des satellites, pondéré par l'inverse des distances (ou la qualité du signal). C'est exactement la formule vectorielle $\overrightarrow{MG} = \frac{\sum \alpha_i \overrightarrow{MA_i}}{\sum \alpha_i}$.
    \item \textbf{Électrification rurale (EDC/ENEO) :} Pour placer un transformateur alimentant trois villages $A, B, C$, on cherche le point $G$ minimisant la somme des carrés des distances pondérées par la population : $G$ est l'isobarycentre (centre de gravité) si populations égales, ou barycentre pondéré par le nombre d'habitants.
    \item \textbf{Cartographie (MINTP) :} Le centre d'un rond-point ou d'un échangeur est souvent conçu comme barycentre des axes routiers entrants pour optimiser les flux.
\end{itemize}
Les lignes de niveau $MA^2+MB^2=k$ modélisent les zones d'intensité de signal constante (somme de puissances), et $MA/MB=k$ les frontières de "handover" (basculement d'une antenne à l'autre).
\end{savaistu}

\newpage

\coursec{Exercices}

\courssub{Je vérifie mes connaissances}

\begin{exercice}[QCM : le barycentre en un coup d'œil]
Pour chaque question, une seule réponse est exacte. Recopier la lettre correspondante sur la copie.
\begin{enumerate}
\item Le barycentre $G$ du système $\{(A,2)\,;\,(B,2)\}$ est :
\begin{enumerate}
\item le point $A$ ; \item le milieu du segment $[AB]$ ; \item le point $B$ ; \item un point extérieur au segment $[AB]$.
\end{enumerate}
\item Si $G = \mathrm{bar}\{(A,3)\,;\,(B,1)\}$, alors le vecteur $\overrightarrow{AG}$ est égal à :
\begin{enumerate}
\item $\dfrac{3}{4}\overrightarrow{AB}$ ; \item $\dfrac{1}{4}\overrightarrow{AB}$ ; \item $\dfrac{1}{3}\overrightarrow{AB}$ ; \item $3\overrightarrow{AB}$.
\end{enumerate}
\item Le barycentre du système $\{(A,5)\,;\,(B,-5)\}$ :
\begin{enumerate}
\item est le milieu de $[AB]$ ; \item existe et est extérieur à $[AB]$ ; \item n'existe pas ; \item est confondu avec $A$.
\end{enumerate}
\item Si $G = \mathrm{bar}\{(A,1)\,;\,(B,2)\,;\,(C,3)\}$ et $H = \mathrm{bar}\{(B,2)\,;\,(C,3)\}$, alors :
\begin{enumerate}
\item $G$ est le milieu de $[AH]$ ; \item $G \in (AH)$ ; \item $G = H$ ; \item $A$, $G$, $H$ ne sont jamais alignés.
\end{enumerate}
\end{enumerate}
\end{exercice}

\begin{exercice}[Vrai ou faux, en justifiant]
Dire si chacune des affirmations suivantes est vraie ou fausse, en justifiant soigneusement la réponse (une démonstration ou un contre-exemple est attendu).
\begin{enumerate}
\item « Si $\alpha + \beta = 0$, alors le barycentre de $\{(A,\alpha)\,;\,(B,\beta)\}$ est le milieu de $[AB]$. »
\item « L'ensemble des points $M$ tels que $\dfrac{MA}{MB} = 1$ est un cercle. »
\item « Pour tout réel $k$, l'ensemble des points $M$ tels que $MA^2 + MB^2 = k$ est toujours un cercle. »
\item « Le barycentre d'un système de points pondérés ne change pas si l'on multiplie tous les coefficients par $-3$. »
\end{enumerate}
\end{exercice}

\begin{exercice}[Définitions et réductions immédiates]
\begin{enumerate}
\item Compléter : le point $G$ est le barycentre du système $\{(A,\alpha)\,;\,(B,\beta)\}$ avec $\alpha + \beta \neq 0$ si et seulement si $\alpha\overrightarrow{GA} + \beta\overrightarrow{GB} = \ldots$
\item Soit $G = \mathrm{bar}\{(A,1)\,;\,(B,-2)\}$. Exprimer $\overrightarrow{AG}$ en fonction de $\overrightarrow{AB}$, puis placer $G$ sur une figure.
\item Soient $A$, $B$, $C$ trois points et $G = \mathrm{bar}\{(A,2)\,;\,(B,1)\,;\,(C,1)\}$. Pour tout point $M$ du plan, réduire la somme vectorielle $2\overrightarrow{MA} + \overrightarrow{MB} + \overrightarrow{MC}$.
\item Réduire, pour tout point $M$, la somme $\overrightarrow{MA} - \overrightarrow{MB} + 2\overrightarrow{MC}$ où $C = \mathrm{bar}\{(A,1)\,;\,(B,-1)\}$ n'existe pas : que constate-t-on sur la somme des coefficients du système définissant $C$ ?
\end{enumerate}
\end{exercice}

\begin{exercice}[Lignes de niveau : reconnaître la nature]
Soient $A$ et $B$ deux points distincts du plan, $I$ le milieu de $[AB]$, avec $AB = 8\ \mathrm{cm}$. Sans faire de calcul détaillé, donner la nature géométrique attendue pour chacun des ensembles suivants (cercle, droite, médiatrice, segment, ensemble vide, point) :
\begin{enumerate}
\item $\mathcal{E}_1 = \{M \ ;\ MA^2 + MB^2 = 50\}$ ;
\item $\mathcal{E}_2 = \{M \ ;\ MA^2 - MB^2 = 0\}$ ;
\item $\mathcal{E}_3 = \{M \ ;\ \overrightarrow{MA} \cdot \overrightarrow{MB} = 0\}$ ;
\item $\mathcal{E}_4 = \left\{M \ ;\ \dfrac{MA}{MB} = 3\right\}$ ;
\item $\mathcal{E}_5 = \{M \ ;\ MA^2 + MB^2 = 20\}$.
\end{enumerate}
\end{exercice}

\courssub{Je m'entraîne}

\begin{exercice}[Barycentres de deux points]
Sur une droite graduée, on considère les points $A$ et $B$ d'abscisses respectives $-2$ et $6$.
\begin{enumerate}
\item Déterminer l'abscisse du barycentre $G = \mathrm{bar}\{(A,3)\,;\,(B,1)\}$ et construire $G$.
\item Déterminer l'abscisse du barycentre $H = \mathrm{bar}\{(A,1)\,;\,(B,-3)\}$ et construire $H$. Que remarque-t-on sur la position de $H$ par rapport au segment $[AB]$ ?
\item Déterminer les réels $\alpha$ et $\beta$ tels que $A = \mathrm{bar}\{(G,\alpha)\,;\,(B,\beta)\}$ avec $\alpha + \beta = 4$.
\end{enumerate}
\end{exercice}

\begin{exercice}[Barycentre de trois points et barycentres partiels]
Soit $ABC$ un triangle. On pose $G = \mathrm{bar}\{(A,2)\,;\,(B,1)\,;\,(C,3)\}$.
\begin{enumerate}
\item On note $I = \mathrm{bar}\{(B,1)\,;\,(C,3)\}$. Montrer que $G$ est le barycentre de $\{(A,2)\,;\,(I,4)\}$, puis construire $G$.
\item Construire $G$ par une seconde méthode, en utilisant cette fois le barycentre partiel $J = \mathrm{bar}\{(A,2)\,;\,(B,1)\}$.
\item Démontrer que les droites $(AG)$, $(BG)$ et $(CG)$ coupent respectivement les droites $(BC)$, $(CA)$ et $(AB)$ en des points que l'on précisera comme barycentres de deux points.
\end{enumerate}
\end{exercice}

\begin{exercice}[Alignement et concours par les barycentres]
Soit $ABC$ un triangle. On considère les points :
\[ P = \mathrm{bar}\{(A,2)\,;\,(B,1)\}, \qquad Q = \mathrm{bar}\{(B,1)\,;\,(C,2)\}, \qquad R = \mathrm{bar}\{(A,1)\,;\,(C,1)\}. \]
\begin{enumerate}
\item Construire les points $P$, $Q$ et $R$ sur une figure soignée.
\item Montrer que le point $G = \mathrm{bar}\{(A,2)\,;\,(B,1)\,;\,(C,2)\}$ appartient aux trois droites $(AQ)$, $(BR)$ et $(CP)$.
\item En déduire que les droites $(AQ)$, $(BR)$ et $(CP)$ sont concourantes en $G$.
\item Montrer que les points $P$, $Q$ et le point $S = \mathrm{bar}\{(A,2)\,;\,(C,-2)\}$ sont alignés (on pourra chercher un barycentre commun).
\end{enumerate}
\end{exercice}

\begin{exercice}[Premières lignes de niveau]
Soient $A$ et $B$ deux points tels que $AB = 6\ \mathrm{cm}$, et $I$ le milieu de $[AB]$.
\begin{enumerate}
\item Montrer que pour tout point $M$ : $MA^2 + MB^2 = 2MI^2 + \dfrac{AB^2}{2}$.
\item Déterminer et construire l'ensemble $\mathcal{C}$ des points $M$ tels que $MA^2 + MB^2 = 26$.
\item Montrer que pour tout point $M$ : $MA^2 - MB^2 = 2\overrightarrow{IM} \cdot \overrightarrow{BA}$ (on pourra écrire $MA^2 - MB^2 = (\overrightarrow{MA} - \overrightarrow{MB})\cdot(\overrightarrow{MA} + \overrightarrow{MB})$).
\item Déterminer et construire l'ensemble $\mathcal{D}$ des points $M$ tels que $MA^2 - MB^2 = 12$.
\end{enumerate}
\end{exercice}

\courssub{Je me prépare au Probatoire}

\begin{exercice}[Type Probatoire : le chantier de l'école de Foumbot]
Pour implanter une cuve d'eau dans la cour d'une école de Foumbot, un technicien modélise la cour par un plan muni d'un repère. Trois poteaux sont plantés aux points $A$, $B$ et $C$, non alignés. La cuve doit être placée au point $G = \mathrm{bar}\{(A,2)\,;\,(B,-1)\,;\,(C,3)\}$.

\textbf{Partie A : deux constructions de $G$}
\begin{enumerate}
\item Justifier que le point $G$ existe.
\item On pose $H = \mathrm{bar}\{(B,-1)\,;\,(C,3)\}$.
\begin{enumerate}
\item Montrer que $\overrightarrow{BH} = \dfrac{3}{2}\overrightarrow{BC}$.
\item Montrer que $G = \mathrm{bar}\{(A,2)\,;\,(H,2)\}$. En déduire la position de $G$ sur le segment $[AH]$ et construire $G$.
\end{enumerate}
\item On pose maintenant $K = \mathrm{bar}\{(A,2)\,;\,(B,-1)\}$.
\begin{enumerate}
\item Montrer que $\overrightarrow{AK} = -\overrightarrow{AB}$.
\item En utilisant le barycentre partiel $K$, retrouver une construction de $G$ et vérifier la cohérence avec la question 2.
\end{enumerate}
\end{enumerate}

\textbf{Partie B : une droite remarquable}

On note $I$ le milieu de $[AC]$.
\begin{enumerate}
\item Montrer que $I = \mathrm{bar}\{(A,2)\,;\,(C,2)\}$.
\item Montrer que les points $B$, $G$ et $I$ sont alignés (on pourra écrire $G$ comme barycentre de $\{(B,-1)\,;\,(I,4)\}$).
\item En déduire que $G$ appartient à une droite remarquable du triangle $ABC$, que l'on nommera.
\end{enumerate}
\end{exercice}

\begin{popfigure}
\begin{tikzpicture}[scale=0.8, >=Stealth]
% Triangle ABC
\coordinate (A) at (0,0);
\coordinate (B) at (5,1);
\coordinate (C) at (2,4);
% Points
\coordinate (H) at ($(B)!1.5!(C)$); % BH = 3/2 BC
\coordinate (G) at ($(A)!0.5!(H)$); % Midpoint AH
\coordinate (K) at ($(A)!-1!(B)$); % AK = -AB
\coordinate (I) at ($(A)!0.5!(C)$); % Midpoint AC
% Draw
\draw[thick] (A) -- (B) -- (C) -- cycle;
\draw[dashed, popblue] (A) -- (H);
\draw[dashed, poporange] (K) -- (C);
\draw[dashed, popgreen] (B) -- (I);
% Points
\foreach \p/\label/\pos in {A/A/below left, B/B/below right, C/C/above, H/H/right, G/G/above left, K/K/left, I/I/left}
  \fill[popdark] (\p) circle (2pt) node[\pos] {$\label$};
% Labels
\node[popblue, above] at ($(A)!0.5!(H)$) {Construction 1};
\node[poporange, below] at ($(K)!0.5!(C)$) {Construction 2};
\node[popgreen, right] at ($(B)!0.5!(I)$) {Médiane};
\end{tikzpicture}
\legende{Figure pour l'exercice « Chantier de Foumbot » : deux constructions de $G$ et la médiane $(BI)$.}
\end{popfigure}

\begin{exercice}[Type Probatoire : synthèse sur les lignes de niveau]
Sur une carte de la région du Centre, deux relais MTN sont modélisés par deux points $A$ et $B$ tels que $AB = 6\ \mathrm{cm}$ (sur la carte). On note $I$ le milieu de $[AB]$. Un technicien cherche les positions $M$ d'un répéteur vérifiant diverses contraintes.

\textbf{Partie A : contrainte $MA^2 + MB^2 = 26$}
\begin{enumerate}
\item Démontrer que pour tout point $M$ du plan : $MA^2 + MB^2 = 2MI^2 + 18$.
\item En déduire que $MA^2 + MB^2 = 26$ équivaut à $MI = 2$.
\item Caractériser et construire l'ensemble $\mathcal{C}_1$ des points $M$ solutions.
\end{enumerate}

\textbf{Partie B : contrainte $MA^2 - MB^2 = 12$}
\begin{enumerate}
\item Démontrer que pour tout point $M$ : $MA^2 - MB^2 = 2\overrightarrow{AB} \cdot \overrightarrow{IM}$.
\item Soit $H$ le projeté orthogonal de $M$ sur la droite $(AB)$. Montrer que $MA^2 - MB^2 = 12$ équivaut à $\overrightarrow{AB} \cdot \overrightarrow{IH} = 6$, puis que $H$ est un point fixe de $(AB)$ que l'on placera (on donnera la valeur de $\overrightarrow{IH}$ en centimètres, en orientant la droite de $A$ vers $B$).
\item Caractériser et construire l'ensemble $\mathcal{C}_2$ des points $M$ solutions.
\end{enumerate}

\textbf{Partie C : contrainte $\overrightarrow{MA} \cdot \overrightarrow{MB} = 7$}
\begin{enumerate}
\item Démontrer que pour tout point $M$ : $\overrightarrow{MA} \cdot \overrightarrow{MB} = MI^2 - \dfrac{AB^2}{4}$.
\item En déduire la nature et les éléments caractéristiques de l'ensemble $\mathcal{C}_3$ des points $M$ tels que $\overrightarrow{MA} \cdot \overrightarrow{MB} = 7$, puis construire $\mathcal{C}_3$.
\end{enumerate}

\textbf{Partie D : contrainte $\dfrac{MA}{MB} = 2$}
\begin{enumerate}
\item Montrer que $\dfrac{MA}{MB} = 2$ équivaut à $MA^2 - 4MB^2 = 0$.
\item On introduit les points $G_1 = \mathrm{bar}\{(A,1)\,;\,(B,2)\}$ et $G_2 = \mathrm{bar}\{(A,1)\,;\,(B,-2)\}$. Montrer que $MA^2 - 4MB^2 = 0$ équivaut à $\overrightarrow{MG_1} \cdot \overrightarrow{MG_2} = 0$.
\item En déduire la nature de l'ensemble $\mathcal{C}_4$ des points $M$ tels que $\dfrac{MA}{MB} = 2$, préciser son centre et son rayon, puis le construire sur la même figure que $\mathcal{C}_1$, $\mathcal{C}_2$ et $\mathcal{C}_3$.
\end{enumerate}
\end{exercice}

\begin{popfigure}
\begin{tikzpicture}[scale=0.7, >=Stealth]
% Base AB
\coordinate (A) at (-3,0);
\coordinate (B) at (3,0);
\coordinate (I) at (0,0);
% Circles
\draw[popblue, thick] (I) circle (2); % C1 radius 2
\draw[poporange, thick] (I) circle (4); % C3 radius 4
% Line C2
\coordinate (H) at (1,0); % IH = 1 cm towards B
\draw[popgreen, thick, dashed] (H) -- ++(0,3) (H) -- ++(0,-3);
% Circle C4 (Apollonius)
% G1 = bar{(A,1),(B,2)} -> AG1/G1B = 2/1 -> G1 at 1/3 from A? Wait. A=-3, B=3. G1 = (A+2B)/3 = (-3+6)/3 = 1.
% G2 = bar{(A,1),(B,-2)} -> (A-2B)/(-1) = 2B-A = 6-(-3)=9.
% Center O = (1+9)/2 = 5. Radius = 4.
\coordinate (G1) at (1,0);
\coordinate (G2) at (9,0);
\coordinate (O) at (5,0);
\draw[poppurple, thick] (O) circle (4);
% Points
\foreach \p/\label/\pos in {A/A/left, B/B/right, I/I/below, H/H/below, G1/G_1/below, G2/G_2/below, O/O/above}
  \fill[popdark] (\p) circle (2pt) node[\pos] {$\label$};
% Axis
\draw[<->] (-4.5,0) -- (10.5,0) node[right] {$(AB)$};
% Legends
\node[popblue] at (-3,3.5) {$\mathcal{C}_1$};
\node[poporange] at (-3,4.5) {$\mathcal{C}_3$};
\node[popgreen] at (1.5,3) {$\mathcal{C}_2$};
\node[poppurple] at (9,4.5) {$\mathcal{C}_4$};
\end{tikzpicture}
\legende{Figure de synthèse pour l'exercice « Relais MTN » : les quatre lignes de niveau $\mathcal{C}_1$, $\mathcal{C}_2$, $\mathcal{C}_3$, $\mathcal{C}_4$.}
\end{popfigure}

\begin{exercice}[Type Probatoire : quadrilatère et isobarycentre de quatre points]
Un lotissement de la ville de Bafoussam a la forme d'un quadrilatère $ABCD$ quelconque. On note $P$, $Q$, $R$, $S$ les milieux respectifs des côtés $[AB]$, $[BC]$, $[CD]$, $[DA]$, et $U$, $V$ les milieux respectifs des diagonales $[AC]$ et $[BD]$. On note enfin $G$ l'isobarycentre des quatre points $A$, $B$, $C$, $D$.

\textbf{Partie A : concours de trois segments}
\begin{enumerate}
\item Justifier l'existence de $G$ et rappeler la relation vectorielle qui le caractérise.
\item En utilisant les barycentres partiels $P$ et $R$, montrer que $G$ est le milieu du segment $[PR]$.
\item Montrer de même que $G$ est le milieu du segment $[QS]$ et le milieu du segment $[UV]$.
\item En déduire que les segments $[PR]$, $[QS]$ et $[UV]$ sont concourants en leur milieu commun.
\item Quelle est la nature du quadrilatère $PQRS$ ? Justifier.
\end{enumerate}

\textbf{Partie B : une égalité vectorielle et son exploitation}

Soit $M$ un point quelconque du plan.
\begin{enumerate}
\item Montrer que $\overrightarrow{MA} + \overrightarrow{MB} + \overrightarrow{MC} + \overrightarrow{MD} = 4\overrightarrow{MG}$.
\item On suppose dans cette question que $ABCD$ est un parallélogramme. Montrer que $G$ est le point d'intersection des diagonales $[AC]$ et $[BD]$.
\item Déterminer l'ensemble $\mathcal{E}$ des points $M$ du plan tels que :
\[ \left\| \overrightarrow{MA} + \overrightarrow{MB} + \overrightarrow{MC} + \overrightarrow{MD} \right\| = 4\,AB. \]
(On exprimera le résultat à l'aide du point $G$.)
\end{enumerate}
\end{exercice}

\begin{popfigure}
\begin{tikzpicture}[scale=0.8, >=Stealth]
% Quadrilateral ABCD
\coordinate (A) at (0,0);
\coordinate (B) at (5,1);
\coordinate (C) at (4,4);
\coordinate (D) at (1,3);
% Midpoints
\coordinate (P) at ($(A)!0.5!(B)$);
\coordinate (Q) at ($(B)!0.5!(C)$);
\coordinate (R) at ($(C)!0.5!(D)$);
\coordinate (S) at ($(D)!0.5!(A)$);
\coordinate (U) at ($(A)!0.5!(C)$);
\coordinate (V) at ($(B)!0.5!(D)$);
\coordinate (G) at ($(P)!0.5!(R)$); % or (U)!0.5!(V)
% Draw
\draw[thick] (A) -- (B) -- (C) -- (D) -- cycle;
\draw[popblue, dashed] (P) -- (R);
\draw[poporange, dashed] (Q) -- (S);
\draw[popgreen, dashed] (U) -- (V);
\draw[poppurple, thick] (P) -- (Q) -- (R) -- (S) -- cycle; % Varignon parallelogram
% Points
\foreach \p/\label/\pos in {A/A/below left, B/B/below right, C/C/above right, D/D/above left, P/P/below, Q/Q/right, R/R/above, S/S/left, U/U/above left, V/V/below right, G/G/above}
  \fill[popdark] (\p) circle (2pt) node[\pos] {$\label$};
\end{tikzpicture}
\legende{Quadrilatère $ABCD$, milieux $P,Q,R,S,U,V$ et isobarycentre $G$ (concours de $[PR]$, $[QS]$, $[UV]$). $PQRS$ est un parallélogramme (Varignon).}
\end{popfigure}

\begin{aretenir}[Formules clés pour les lignes de niveau (Probatoire)]
Soient $A$, $B$ deux points distincts, $I$ le milieu de $[AB]$, $AB = d$.
\begin{itemize}
\item \textbf{Médiane (Leibniz) :} $MA^2 + MB^2 = 2MI^2 + \dfrac{d^2}{2}$.
\item \textbf{Différence des carrés :} $MA^2 - MB^2 = 2\,\overrightarrow{AB} \cdot \overrightarrow{IM} = 2\,\overrightarrow{BA} \cdot \overrightarrow{MI}$.
\item \textbf{Produit scalaire :} $\overrightarrow{MA} \cdot \overrightarrow{MB} = MI^2 - \dfrac{d^2}{4}$.
\item \textbf{Rapport des distances (Apollonius) :} $\dfrac{MA}{MB} = k \ (k>0, k\neq 1) \iff M$ appartient au cercle de diamètre $[G_1G_2]$ où $G_1 = \mathrm{bar}\{(A,1);(B,k)\}$ et $G_2 = \mathrm{bar}\{(A,1);(B,-k)\}$.
\end{itemize}
\end{aretenir}

\begin{methode}[Réduire une somme vectorielle $\sum \alpha_i \overrightarrow{MA_i}$]
\begin{enumerate}
\item Calculer la somme des coefficients $S = \sum \alpha_i$.
\item \textbf{Si $S \neq 0$ :} Le barycentre $G$ existe. Écrire $\overrightarrow{MA_i} = \overrightarrow{MG} + \overrightarrow{GA_i}$. La somme devient $S\overrightarrow{MG} + \sum \alpha_i \overrightarrow{GA_i} = S\overrightarrow{MG}$ (car $\sum \alpha_i \overrightarrow{GA_i} = \vec{0}$ par définition de $G$).
\item \textbf{Si $S = 0$ :} La somme ne dépend pas de $M$. Choisir un origine commode (souvent l'un des $A_i$) pour calculer le vecteur constant $\sum \alpha_i \overrightarrow{OA_i}$.
\end{enumerate}
\end{methode}

\begin{attention}[Pièges classiques au Probatoire]
\begin{itemize}
\item \textbf{Somme des coefficients nulle :} Ne pas écrire « le barycentre n'existe pas » sans vérifier si les points sont confondus. Pour deux points distincts, somme nulle $\implies$ pas de barycentre.
\item \textbf{Orientation des vecteurs :} Dans $MA^2 - MB^2 = 2\overrightarrow{AB}\cdot\overrightarrow{IM}$, l'ordre compte. Vérifier avec un cas simple ($M=B$ : $BA^2 - 0 = 2\overrightarrow{AB}\cdot\overrightarrow{IB} = 2\overrightarrow{AB}\cdot(-\frac{1}{2}\overrightarrow{AB}) = -AB^2$ ? Non. $M=B \implies MB=0, MA=AB$. $AB^2 = 2\overrightarrow{AB}\cdot\overrightarrow{IB}$. $\overrightarrow{IB} = \frac{1}{2}\overrightarrow{AB}$. RHS $= 2\overrightarrow{AB}\cdot\frac{1}{2}\overrightarrow{AB} = AB^2$. Correct.)
\item \textbf{Ensemble vide :} Pour $MA^2+MB^2=k$, le cercle existe seulement si $k \ge \frac{AB^2}{2}$. Pour $\overrightarrow{MA}\cdot\overrightarrow{MB}=k$, le cercle existe seulement si $k \ge -\frac{AB^2}{4}$.
\item \textbf{Barycentre partiel :} Le poids du barycentre partiel est la \textbf{somme} des poids qui le composent.
\end{itemize}
\end{attention}

\begin{savaistu}[Le cercle d'Apollonius et les télécommunications]
Le cercle d'Apollonius ($\frac{MA}{MB}=k$) trouve une application concrète dans la localisation par différence de temps d'arrivée (TDOA) utilisée par les réseaux MTN/Orange au Cameroun. Si un téléphone reçoit un signal de deux antennes $A$ et $B$ avec des intensités proportionnelles aux distances (ou des temps de propagation différents), l'ensemble des positions possibles du téléphone est un arc de cercle d'Apollonius. Avec trois antennes, l'intersection de deux tels cercles donne la position précise (trilatération), principe de base du GPS et de la géolocalisation GSM.
\end{savaistu}

\newpage

\coursec{Corrigés des exercices}

\begin{corrige}[Exercice 1 — QCM : le barycentre en un coup d'œil]
\begin{enumerate}
\item \textbf{Réponse b.} Les coefficients sont égaux : $G$ est l'isobarycentre de $A$ et $B$, c'est-à-dire le milieu de $[AB]$.
\item \textbf{Réponse b.} $G = \mathrm{bar}\{(A,3)\,;\,(B,1)\}$ donne $3\overrightarrow{GA}+\overrightarrow{GB}=\vec{0}$, d'où $3\overrightarrow{GA}+\overrightarrow{GA}+\overrightarrow{AB}=\vec{0}$, soit $4\overrightarrow{GA}=-\overrightarrow{AB}$, c'est-à-dire $\overrightarrow{AG}=\dfrac{1}{4}\overrightarrow{AB}$. (Le coefficient de $B$, qui vaut $1$, se retrouve au numérateur : $G$ est proche de $A$, le point le plus lourd.)
\item \textbf{Réponse c.} La somme des coefficients est $5+(-5)=0$ et $A\neq B$ : le barycentre n'existe pas.
\item \textbf{Réponse b.} Par associativité du barycentre, $G = \mathrm{bar}\{(A,1)\,;\,(H,2+3)\} = \mathrm{bar}\{(A,1)\,;\,(H,5)\}$, donc $G\in(AH)$. Ce n'est pas le milieu de $[AH]$ car les coefficients $1$ et $5$ sont différents.
\end{enumerate}
\end{corrige}

\begin{corrige}[Exercice 2 — Vrai ou faux, en justifiant]
\begin{enumerate}
\item \textbf{Faux.} Si $\alpha+\beta=0$ avec $\alpha\neq 0$ (donc $\beta=-\alpha\neq 0$) et $A\neq B$, le barycentre n'existe pas. Si $\alpha=\beta=0$, le barycentre n'est pas défini. Le milieu de $[AB]$ correspond à $\alpha=\beta\neq 0$, donc $\alpha+\beta=2\alpha\neq 0$.
\item \textbf{Faux.} $\dfrac{MA}{MB}=1\iff MA=MB$ : l'ensemble est la \textbf{médiatrice} du segment $[AB]$, qui est une droite, pas un cercle. (Le cercle d'Apollonius n'apparaît que pour $k>0$ et $k\neq 1$.)
\item \textbf{Faux.} Avec $I$ milieu de $[AB]$, l'identité de la médiane donne $MA^2+MB^2=2MI^2+\dfrac{AB^2}{2}$. L'équation $MA^2+MB^2=k$ équivaut à $MI^2=\dfrac{2k-AB^2}{4}$. Si $k<\dfrac{AB^2}{2}$, l'ensemble est \textbf{vide} ; si $k=\dfrac{AB^2}{2}$, il est réduit au point $I$. Ce n'est un cercle que si $k>\dfrac{AB^2}{2}$.
\item \textbf{Vrai.} Le barycentre est homogène : $\mathrm{bar}\{(A_i,\alpha_i)\}=\mathrm{bar}\{(A_i,\lambda\alpha_i)\}$ pour tout $\lambda\neq 0$, car $\sum\lambda\alpha_i\overrightarrow{GA_i}=\lambda\sum\alpha_i\overrightarrow{GA_i}=\vec{0}$ et $\lambda\sum\alpha_i\neq 0$. Avec $\lambda=-3$, le barycentre est inchangé.
\end{enumerate}
\end{corrige}

\begin{corrige}[Exercice 3 — Définitions et réductions immédiates]
\begin{enumerate}
\item $\alpha\overrightarrow{GA}+\beta\overrightarrow{GB}=\vec{0}$ (le vecteur nul).
\item $G=\mathrm{bar}\{(A,1)\,;\,(B,-2)\}$ : $\overrightarrow{GA}-2\overrightarrow{GB}=\vec{0}$. En introduisant $A$ dans $\overrightarrow{GB}$ : $\overrightarrow{GA}-2(\overrightarrow{GA}+\overrightarrow{AB})=\vec{0}$, soit $-\overrightarrow{GA}-2\overrightarrow{AB}=\vec{0}$, d'où $\overrightarrow{AG}=-2\overrightarrow{AB}$. Le point $G$ est sur la droite $(AB)$, du côté opposé à $B$ par rapport à $A$, avec $AG=2AB$ (coefficients de signes contraires : $G$ est extérieur au segment $[AB]$).
\item La somme des coefficients est $2+1+1=4\neq 0$, donc le barycentre $G$ du système existe et, par la propriété fondamentale du barycentre, pour tout point $M$ :
\[ 2\overrightarrow{MA}+\overrightarrow{MB}+\overrightarrow{MC}=4\overrightarrow{MG}. \]
\item La somme des coefficients est $1-1+2=2\neq 0$ : la somme se réduit. Notons $\Gamma=\mathrm{bar}\{(A,1)\,;\,(B,-1)\,;\,(C,2)\}$, qui existe car la somme \emph{totale} des coefficients vaut $2\neq 0$. Alors, pour tout point $M$ :
\[ \overrightarrow{MA}-\overrightarrow{MB}+2\overrightarrow{MC}=2\overrightarrow{M\Gamma}. \]
\textbf{Remarque importante :} le sous-système $\{(A,1)\,;\,(B,-1)\}$ a une somme de coefficients nulle, donc il n'admet pas de barycentre partiel ; mais cela n'empêche pas le système complet d'avoir un barycentre, car seule la somme \emph{totale} compte. On peut d'ailleurs réduire directement : pour tout point $M$, $\overrightarrow{MA}-\overrightarrow{MB}=\overrightarrow{BA}$ (vecteur constant), donc
\[ \overrightarrow{MA}-\overrightarrow{MB}+2\overrightarrow{MC}=\overrightarrow{BA}+2\overrightarrow{MC}. \]
\end{enumerate}
\end{corrige}

\begin{corrige}[Exercice 4 — Lignes de niveau : reconnaître la nature]
On a $AB=8$ cm, donc $\dfrac{AB^2}{2}=32$ et $\dfrac{AB^2}{4}=16$. On note $I$ le milieu de $[AB]$.
\begin{enumerate}
\item $\mathcal{E}_1$ : $MA^2+MB^2=50\iff 2MI^2+32=50\iff MI^2=9$. Comme $50>32=\dfrac{AB^2}{2}$, c'est un \textbf{cercle} de centre $I$ et de rayon $3$ cm.
\item $\mathcal{E}_2$ : $MA^2-MB^2=0\iff MA=MB$ (car $MA,MB\ge 0$) : c'est la \textbf{médiatrice} de $[AB]$ (cas particulier de droite).
\item $\mathcal{E}_3$ : $\overrightarrow{MA}\cdot\overrightarrow{MB}=0\iff MI^2-\dfrac{AB^2}{4}=0\iff MI=4$ : \textbf{cercle de diamètre $[AB]$} (centre $I$, rayon $4$ cm).
\item $\mathcal{E}_4$ : $\dfrac{MA}{MB}=3$ avec $3>0$ et $3\neq 1$ : \textbf{cercle d'Apollonius}, de diamètre $[G_1G_2]$ avec $G_1=\mathrm{bar}\{(A,1)\,;\,(B,3)\}$ et $G_2=\mathrm{bar}\{(A,1)\,;\,(B,-3)\}$ (les deux barycentres existent : $1+3=4\neq 0$ et $1-3=-2\neq 0$).
\item $\mathcal{E}_5$ : $MA^2+MB^2=20$ avec $20<32=\dfrac{AB^2}{2}$ : \textbf{ensemble vide}.
\end{enumerate}
\end{corrige}

\begin{corrige}[Exercice 5 — Barycentres de deux points]
\begin{enumerate}
\item L'abscisse du barycentre est le barycentre des abscisses :
\[ x_G=\dfrac{3\times(-2)+1\times 6}{3+1}=\dfrac{-6+6}{4}=0. \]
$G$ est à l'origine de la droite graduée. Vérification : $\overrightarrow{AG}=\dfrac{1}{4}\overrightarrow{AB}$, et $\frac14(6-(-2))=2$, donc $x_G=-2+2=0$.
\item $x_H=\dfrac{1\times(-2)+(-3)\times 6}{1-3}=\dfrac{-2-18}{-2}=\dfrac{-20}{-2}=10$.
Vérification : $\overrightarrow{AH}=\dfrac{-3}{-2}\overrightarrow{AB}=\dfrac32\overrightarrow{AB}$, soit $x_H=-2+\frac32\times 8=10$.
\textbf{Remarque :} les coefficients $1$ et $-3$ sont de signes contraires, donc $H$ est \textbf{extérieur au segment} $[AB]$ (ici au-delà de $B$, car $|-3|>|1|$).
\item $A=\mathrm{bar}\{(G,\alpha)\,;\,(B,\beta)\}$ avec $\alpha+\beta=4$ signifie $\alpha\overrightarrow{AG}+\beta\overrightarrow{AB}=\vec{0}$. Or $\overrightarrow{AG}=\dfrac14\overrightarrow{AB}$, donc $\dfrac{\alpha}{4}+\beta=0$, soit $\alpha=-4\beta$. Avec $\alpha+\beta=4$ : $-4\beta+\beta=4$, d'où $\beta=-\dfrac{4}{3}$ et $\alpha=\dfrac{16}{3}$. Vérification : $\dfrac{16}{3}-\dfrac{4}{3}=\dfrac{12}{3}=4$. (Tout couple proportionnel conviendrait si la condition $\alpha+\beta=4$ n'était pas imposée.)
\end{enumerate}
\end{corrige}

\begin{corrige}[Exercice 6 — Barycentre de trois points et barycentres partiels]
\begin{enumerate}
\item $I=\mathrm{bar}\{(B,1)\,;\,(C,3)\}$ existe car $1+3=4\neq 0$, et $\overrightarrow{BI}=\dfrac{3}{4}\overrightarrow{BC}$. Par le théorème d'associativité (barycentre partiel), en regroupant $B$ et $C$ :
\[ G=\mathrm{bar}\{(A,2)\,;\,(I,1+3)\}=\mathrm{bar}\{(A,2)\,;\,(I,4)\}. \]
Donc $\overrightarrow{AG}=\dfrac{4}{2+4}\overrightarrow{AI}=\dfrac{2}{3}\overrightarrow{AI}$. Construction : placer $I$ sur $[BC]$ avec $BI=\frac34 BC$, puis $G$ sur $[AI]$ avec $AG=\frac23 AI$.
\item $J=\mathrm{bar}\{(A,2)\,;\,(B,1)\}$ : $\overrightarrow{AJ}=\dfrac{1}{3}\overrightarrow{AB}$. Alors $G=\mathrm{bar}\{(J,2+1)\,;\,(C,3)\}=\mathrm{bar}\{(J,3)\,;\,(C,3)\}$ : $G$ est le \textbf{milieu de $[JC]$}. Construction : placer $J$ au tiers de $[AB]$ depuis $A$, puis prendre le milieu de $[JC]$. Les deux méthodes donnent le même point (unicité du barycentre).
\item La droite $(AG)$ : comme $G=\mathrm{bar}\{(A,2)\,;\,(I,4)\}$, $G\in(AI)$, donc $(AG)=(AI)$ coupe $(BC)$ en $I=\mathrm{bar}\{(B,1)\,;\,(C,3)\}$.
La droite $(BG)$ : regroupons $A$ et $C$ : soit $L=\mathrm{bar}\{(A,2)\,;\,(C,3)\}$ (somme $5\neq 0$). Alors $G=\mathrm{bar}\{(B,1)\,;\,(L,5)\}$, donc $(BG)$ coupe $(CA)$ en $L=\mathrm{bar}\{(A,2)\,;\,(C,3)\}$.
La droite $(CG)$ : avec $J=\mathrm{bar}\{(A,2)\,;\,(B,1)\}$, $G=\mathrm{bar}\{(J,3)\,;\,(C,3)\}$, donc $(CG)$ coupe $(AB)$ en $J=\mathrm{bar}\{(A,2)\,;\,(B,1)\}$.
\end{enumerate}
\end{corrige}

\begin{corrige}[Exercice 7 — Alignement et concours par les barycentres]
\begin{enumerate}
\item Construction : $\overrightarrow{AP}=\dfrac13\overrightarrow{AB}$ ; $\overrightarrow{BQ}=\dfrac{2}{3}\overrightarrow{BC}$ ; $R$ milieu de $[AC]$.
\item La somme des coefficients de $G$ est $2+1+2=5\neq 0$ : $G$ existe.
\begin{itemize}
\item Regroupons $B$ et $C$ : $Q=\mathrm{bar}\{(B,1)\,;\,(C,2)\}$, de poids total $3$. Donc $G=\mathrm{bar}\{(A,2)\,;\,(Q,3)\}$, d'où $G\in(AQ)$.
\item Regroupons $A$ et $C$ : $R=\mathrm{bar}\{(A,1)\,;\,(C,1)\}=\mathrm{bar}\{(A,2)\,;\,(C,2)\}$ (homogénéité), de poids total $4$. Donc $G=\mathrm{bar}\{(B,1)\,;\,(R,4)\}$, d'où $G\in(BR)$.
\item Regroupons $A$ et $B$ : $P=\mathrm{bar}\{(A,2)\,;\,(B,1)\}$, de poids total $3$. Donc $G=\mathrm{bar}\{(P,3)\,;\,(C,2)\}$, d'où $G\in(CP)$.
\end{itemize}
\item Les trois droites $(AQ)$, $(BR)$ et $(CP)$ passent toutes par $G$ : elles sont \textbf{concourantes en $G$}.
\item \textbf{Le point $S=\mathrm{bar}\{(A,2)\,;\,(C,-2)\}$ n'existe pas} : la somme des coefficients est $2+(-2)=0$ et $A\neq C$. C'est le piège de l'énoncé : au Probatoire, on vérifie \emph{toujours} la somme des coefficients avant toute construction.

Ce résultat a une interprétation géométrique précise : montrons que $(PQ)$ est \textbf{parallèle à $(AC)$}, ce qui explique que le « point d'intersection » de $(PQ)$ avec $(AC)$ n'existe pas. Choisissons une origine $O$ quelconque. De $P=\mathrm{bar}\{(A,2)\,;\,(B,1)\}$ on tire $3\overrightarrow{OP}=2\overrightarrow{OA}+\overrightarrow{OB}$, et de $Q=\mathrm{bar}\{(B,1)\,;\,(C,2)\}$ on tire $3\overrightarrow{OQ}=\overrightarrow{OB}+2\overrightarrow{OC}$. Par soustraction :
\[ 3\bigl(\overrightarrow{OP}-\overrightarrow{OQ}\bigr)=2\overrightarrow{OA}-2\overrightarrow{OC}
\quad\Longrightarrow\quad 3\overrightarrow{QP}=2\overrightarrow{CA}. \]
Les vecteurs $\overrightarrow{QP}$ et $\overrightarrow{CA}$ sont colinéaires : les droites $(PQ)$ et $(AC)$ sont parallèles. Un point $S$ de $(AC)$ aligné avec $P$ et $Q$ devrait être « à l'infini » : c'est exactement ce que traduit algébriquement la somme nulle $2+(-2)=0$.
\end{enumerate}
\end{corrige}

\begin{corrige}[Exercice 8 — Premières lignes de niveau]
\begin{enumerate}
\item $I$ milieu de $[AB]$, donc $\overrightarrow{IA}+\overrightarrow{IB}=\vec{0}$ et $IA=IB=\dfrac{AB}{2}$. Pour tout point $M$ :
\begin{align*}
MA^2+MB^2&=\|\overrightarrow{MI}+\overrightarrow{IA}\|^2+\|\overrightarrow{MI}+\overrightarrow{IB}\|^2\\
&=2MI^2+2\overrightarrow{MI}\cdot(\overrightarrow{IA}+\overrightarrow{IB})+IA^2+IB^2\\
&=2MI^2+0+\dfrac{AB^2}{4}+\dfrac{AB^2}{4}=2MI^2+\dfrac{AB^2}{2}.
\end{align*}
C'est l'\textbf{identité de la médiane}.
\item Avec $AB=6$ : $MA^2+MB^2=26\iff 2MI^2+18=26\iff MI^2=4\iff MI=2$. L'ensemble $\mathcal{C}$ est le \textbf{cercle de centre $I$ et de rayon $2$ cm}. Construction : cercle centré au milieu de $[AB]$, passant par les points situés à $2$ cm de $I$.
\item $MA^2-MB^2=(\overrightarrow{MA}-\overrightarrow{MB})\cdot(\overrightarrow{MA}+\overrightarrow{MB})$. Or $\overrightarrow{MA}-\overrightarrow{MB}=\overrightarrow{BA}$ et $\overrightarrow{MA}+\overrightarrow{MB}=2\overrightarrow{MI}$ (car $I$ est le milieu de $[AB]$). Donc :
\[ MA^2-MB^2=2\overrightarrow{MI}\cdot\overrightarrow{BA}=2\overrightarrow{IM}\cdot\overrightarrow{AB}. \]
\item $MA^2-MB^2=12\iff 2\overrightarrow{IM}\cdot\overrightarrow{AB}=12\iff \overrightarrow{IM}\cdot\overrightarrow{AB}=6$. Soit $H$ le projeté orthogonal de $M$ sur $(AB)$ : comme $\overrightarrow{HM}\perp\overrightarrow{AB}$, on a $\overrightarrow{IM}\cdot\overrightarrow{AB}=\overrightarrow{IH}\cdot\overrightarrow{AB}$. En orientant $(AB)$ de $A$ vers $B$ ($\overrightarrow{AB}=6$) : $6\,\overline{IH}=6$, soit $\overline{IH}=1$ : $H$ est le point \textbf{fixe} de $(AB)$ situé à $1$ cm de $I$ du côté de $B$. L'ensemble $\mathcal{D}$ est la \textbf{droite perpendiculaire à $(AB)$ passant par $H$}.
\end{enumerate}
\end{corrige}

\begin{corrige}[Exercice 9 — Type Probatoire : le chantier de l'école de Foumbot]
\textbf{Partie A}
\begin{enumerate}
\item La somme des coefficients est $2+(-1)+3=4\neq 0$ : le barycentre $G$ existe.
\item \begin{enumerate}
\item $H=\mathrm{bar}\{(B,-1)\,;\,(C,3)\}$ : $-\overrightarrow{HB}+3\overrightarrow{HC}=\vec{0}$. En introduisant $B$ dans $\overrightarrow{HC}$ : $-\overrightarrow{HB}+3\overrightarrow{HB}+3\overrightarrow{BC}=\vec{0}$, soit $2\overrightarrow{HB}=-3\overrightarrow{BC}$, d'où $\overrightarrow{BH}=\dfrac{3}{2}\overrightarrow{BC}$.
\item Par associativité : $G=\mathrm{bar}\{(A,2)\,;\,(H,-1+3)\}=\mathrm{bar}\{(A,2)\,;\,(H,2)\}$. Les coefficients sont égaux : $G$ est le \textbf{milieu du segment $[AH]$}. Construction : placer $H$ sur la demi-droite $[BC)$ avec $BH=\frac32 BC$, puis prendre le milieu de $[AH]$.
\end{enumerate}
\item \begin{enumerate}
\item $K=\mathrm{bar}\{(A,2)\,;\,(B,-1)\}$ : $2\overrightarrow{KA}-\overrightarrow{KB}=\vec{0}$. En introduisant $A$ dans $\overrightarrow{KB}$ : $2\overrightarrow{KA}-\overrightarrow{KA}-\overrightarrow{AB}=\vec{0}$, soit $\overrightarrow{KA}=\overrightarrow{AB}$, d'où $\overrightarrow{AK}=-\overrightarrow{AB}$ : $K$ est le symétrique de $B$ par rapport à $A$.
\item Par associativité : $G=\mathrm{bar}\{(K,2-1)\,;\,(C,3)\}=\mathrm{bar}\{(K,1)\,;\,(C,3)\}$, donc $\overrightarrow{KG}=\dfrac{3}{4}\overrightarrow{KC}$. Cohérence : les deux constructions aboutissent au même point $G$ (unicité du barycentre), comme le montre la figure.
\end{enumerate}
\end{enumerate}
\textbf{Partie B}
\begin{enumerate}
\item $I$ milieu de $[AC]$ : $I=\mathrm{bar}\{(A,1)\,;\,(C,1)\}=\mathrm{bar}\{(A,2)\,;\,(C,2)\}$ par homogénéité (multiplication des coefficients par $2$).
\item Regroupons $A$ et $C$ : $G=\mathrm{bar}\{(B,-1)\,;\,(I,2+2)\}=\mathrm{bar}\{(B,-1)\,;\,(I,4)\}$ (somme $3\neq 0$). Donc $G\in(BI)$ : les points $B$, $G$, $I$ sont alignés.
\item $I$ étant le milieu de $[AC]$, la droite $(BI)$ est la \textbf{médiane issue de $B$} du triangle $ABC$. Donc $G$ appartient à la médiane issue de $B$.
\end{enumerate}
\textbf{Barème indicatif (10 pts) :} A1 : 1 pt ; A2 : 2,5 pts ; A3 : 2 pts ; B1 : 1 pt ; B2 : 2 pts ; B3 : 1,5 pt.
\textbf{Points de vigilance :} justifier l'existence de chaque barycentre par le calcul de la somme des coefficients \emph{avant} de l'utiliser ; le poids d'un barycentre partiel est la somme des poids ; en B2, citer explicitement le théorème d'associativité.
\end{corrige}

\begin{corrige}[Exercice 10 — Type Probatoire : synthèse sur les lignes de niveau]
\textbf{Partie A}
\begin{enumerate}
\item Identité de la médiane (démontrée à l'exercice 8) : $MA^2+MB^2=2MI^2+\dfrac{AB^2}{2}=2MI^2+\dfrac{36}{2}=2MI^2+18$.
\item $MA^2+MB^2=26\iff 2MI^2+18=26\iff MI^2=4\iff MI=2$.
\item $\mathcal{C}_1$ est le \textbf{cercle de centre $I$ et de rayon $2$ cm}.
\end{enumerate}
\textbf{Partie B}
\begin{enumerate}
\item $MA^2-MB^2=(\overrightarrow{MA}-\overrightarrow{MB})\cdot(\overrightarrow{MA}+\overrightarrow{MB})=\overrightarrow{BA}\cdot 2\overrightarrow{MI}=2\overrightarrow{AB}\cdot\overrightarrow{IM}$.
\item Soit $H$ le projeté orthogonal de $M$ sur $(AB)$ : $\overrightarrow{IM}=\overrightarrow{IH}+\overrightarrow{HM}$ avec $\overrightarrow{HM}\perp\overrightarrow{AB}$, donc $\overrightarrow{AB}\cdot\overrightarrow{IM}=\overrightarrow{AB}\cdot\overrightarrow{IH}$. Ainsi $MA^2-MB^2=12\iff 2\,\overrightarrow{AB}\cdot\overrightarrow{IH}=12\iff \overrightarrow{AB}\cdot\overrightarrow{IH}=6$. En orientant $(AB)$ de $A$ vers $B$ ($\overrightarrow{AB}=6$) : $6\,\overline{IH}=6$, soit $\overline{IH}=1$ cm. $H$ est le point \textbf{fixe} de $(AB)$ situé à $1$ cm de $I$ du côté de $B$.
\item $\mathcal{C}_2$ est la \textbf{droite perpendiculaire à $(AB)$ en $H$}.
\end{enumerate}
\textbf{Partie C}
\begin{enumerate}
\item $\overrightarrow{MA}\cdot\overrightarrow{MB}=(\overrightarrow{MI}+\overrightarrow{IA})\cdot(\overrightarrow{MI}+\overrightarrow{IB})=MI^2+\overrightarrow{MI}\cdot(\overrightarrow{IA}+\overrightarrow{IB})+\overrightarrow{IA}\cdot\overrightarrow{IB}=MI^2+0-IA^2=MI^2-\dfrac{AB^2}{4}$, car $\overrightarrow{IA}+\overrightarrow{IB}=\vec{0}$ et $\overrightarrow{IA}\cdot\overrightarrow{IB}=-IA^2$ (vecteurs opposés).
\item $\overrightarrow{MA}\cdot\overrightarrow{MB}=7\iff MI^2=7+\dfrac{36}{4}=7+9=16\iff MI=4$. $\mathcal{C}_3$ est le \textbf{cercle de centre $I$ et de rayon $4$ cm}.
\end{enumerate}
\textbf{Partie D}
\begin{enumerate}
\item Comme $MB>0$ et $MA\ge 0$ : $\dfrac{MA}{MB}=2\iff MA=2MB\iff MA^2=4MB^2\iff MA^2-4MB^2=0$ (l'élévation au carré est licite car les deux membres sont positifs).
\item $G_1=\mathrm{bar}\{(A,1)\,;\,(B,2)\}$ existe ($1+2=3\neq 0$) et $\overrightarrow{MA}+2\overrightarrow{MB}=3\overrightarrow{MG_1}$. $G_2=\mathrm{bar}\{(A,1)\,;\,(B,-2)\}$ existe ($1-2=-1\neq 0$) et $\overrightarrow{MA}-2\overrightarrow{MB}=-\overrightarrow{MG_2}$. Alors :
\[ MA^2-4MB^2=(\overrightarrow{MA}-2\overrightarrow{MB})\cdot(\overrightarrow{MA}+2\overrightarrow{MB})=(-\overrightarrow{MG_2})\cdot(3\overrightarrow{MG_1})=-3\,\overrightarrow{MG_1}\cdot\overrightarrow{MG_2}. \]
Donc $MA^2-4MB^2=0\iff\overrightarrow{MG_1}\cdot\overrightarrow{MG_2}=0$.
\item $\overrightarrow{MG_1}\cdot\overrightarrow{MG_2}=0$ signifie que $M$ voit le segment $[G_1G_2]$ sous un angle droit : $\mathcal{C}_4$ est le \textbf{cercle de diamètre $[G_1G_2]$} (cercle d'Apollonius). Calculs : $\overrightarrow{AG_1}=\dfrac{2}{3}\overrightarrow{AB}$, donc $G_1$ est à $4$ cm de $A$ vers $B$ (soit à $1$ cm de $I$ vers $B$) ; $\overrightarrow{AG_2}=\dfrac{-2}{-1}\overrightarrow{AB}=2\overrightarrow{AB}$, donc $G_2$ est à $12$ cm de $A$ (soit à $9$ cm de $I$ vers $B$). Le centre $O$ est le milieu de $[G_1G_2]$, situé à $5$ cm de $I$ vers $B$, et le rayon vaut $\dfrac{G_1G_2}{2}=\dfrac{8}{2}=4$ cm.
\end{enumerate}
\textbf{Barème indicatif (10 pts) :} A : 2 pts ; B : 3 pts ; C : 2 pts ; D : 3 pts.
\textbf{Points de vigilance :} citer $I$ milieu de $[AB]$ avant d'utiliser les identités ; en B, justifier que le projeté $H$ est fixe (indépendant de $M$) ; en D1, préciser que $MA,MB\ge 0$ permet l'élévation au carré ; en D2, vérifier que les deux barycentres existent ($1+2=3\neq 0$ et $1-2=-1\neq 0$).
\end{corrige}

\begin{corrige}[Exercice 11 — Type Probatoire : quadrilatère et isobarycentre de quatre points]
\textbf{Partie A}
\begin{enumerate}
\item La somme des coefficients est $1+1+1+1=4\neq 0$ : $G$ existe. Il est caractérisé par :
\[ \overrightarrow{GA}+\overrightarrow{GB}+\overrightarrow{GC}+\overrightarrow{GD}=\vec{0}. \]
\item $P$ milieu de $[AB]$ : $P=\mathrm{bar}\{(A,1)\,;\,(B,1)\}$, de poids total $2$. $R$ milieu de $[CD]$ : $R=\mathrm{bar}\{(C,1)\,;\,(D,1)\}$, de poids total $2$. Par associativité : $G=\mathrm{bar}\{(P,2)\,;\,(R,2)\}$ : $G$ est le \textbf{milieu de $[PR]$}.
\item De même : $Q=\mathrm{bar}\{(B,1)\,;\,(C,1)\}$ et $S=\mathrm{bar}\{(D,1)\,;\,(A,1)\}$ donnent $G=\mathrm{bar}\{(Q,2)\,;\,(S,2)\}$ : $G$ milieu de $[QS]$. Et $U=\mathrm{bar}\{(A,1)\,;\,(C,1)\}$, $V=\mathrm{bar}\{(B,1)\,;\,(D,1)\}$ donnent $G=\mathrm{bar}\{(U,2)\,;\,(V,2)\}$ : $G$ milieu de $[UV]$.
\item Les trois segments $[PR]$, $[QS]$ et $[UV]$ ont le même milieu $G$ : ils sont \textbf{concourants en $G$, leur milieu commun}.
\item Les diagonales $[PR]$ et $[QS]$ du quadrilatère $PQRS$ se coupent en leur milieu commun $G$ : $PQRS$ est un \textbf{parallélogramme} (théorème de Varignon).
\end{enumerate}
\textbf{Partie B}
\begin{enumerate}
\item Par la propriété fondamentale du barycentre (somme des coefficients $4\neq 0$), pour tout point $M$ :
\[ \overrightarrow{MA}+\overrightarrow{MB}+\overrightarrow{MC}+\overrightarrow{MD}=4\overrightarrow{MG}. \]
(Détail : $\sum\overrightarrow{MA_i}=\sum(\overrightarrow{MG}+\overrightarrow{GA_i})=4\overrightarrow{MG}+\vec{0}=4\overrightarrow{MG}$.)
\item Si $ABCD$ est un parallélogramme, ses diagonales $[AC]$ et $[BD]$ ont le même milieu : $U=V$. Or $G$ est le milieu de $[UV]$ (question A3), donc $G=U=V$ : $G$ est le \textbf{point d'intersection des diagonales} de $ABCD$.
\item D'après B1 : $\left\|\overrightarrow{MA}+\overrightarrow{MB}+\overrightarrow{MC}+\overrightarrow{MD}\right\|=\|4\overrightarrow{MG}\|=4\,MG$ (car $4>0$). L'équation équivaut donc à $4MG=4AB$, soit $MG=AB$. L'ensemble $\mathcal{E}$ est le \textbf{cercle de centre $G$ et de rayon $AB$}.
\end{enumerate}
\textbf{Barème indicatif (10 pts) :} A1 : 1 pt ; A2 : 1,5 pt ; A3 : 2 pts ; A4 : 1 pt ; A5 : 1 pt ; B1 : 1,5 pt ; B2 : 1 pt ; B3 : 1 pt.
\textbf{Points de vigilance :} le poids d'un barycentre partiel est la somme des poids (ici $2$, pas $1$) ; en A5, citer la caractérisation du parallélogramme par les diagonales qui se coupent en leur milieu ; en B3, ne pas oublier que $\|4\overrightarrow{MG}\|=4MG$ car $4>0$.
\end{corrige}

\newpage

\coursec{Fiche bilan}

\begin{popfigure}
\begin{center}\tcbox[colback=popink,colframe=popink,coltext=white,arc=9pt,boxsep=4pt,left=6pt,right=6pt]{\bfseries\small Barycentres et lignes de niveau}\end{center}
\vspace{-2pt}
\begin{tcbraster}[raster columns=3,raster equal height=rows,raster column skip=6pt,raster row skip=6pt,enhanced,blankest]
\begin{tcolorbox}[enhanced,colback=popblue,colframe=popblue,coltext=white,boxrule=0.9pt,arc=3pt,left=4pt,right=4pt,top=3pt,bottom=3pt,boxsep=1pt,halign=left]\bfseries\scriptsize Barycentre de 2 points $\alpha\overrightarrow{GA}+\beta\overrightarrow{GB}=\vec 0$\end{tcolorbox}
\begin{tcolorbox}[enhanced,colback=poporange,colframe=poporange,coltext=white,boxrule=0.9pt,arc=3pt,left=4pt,right=4pt,top=3pt,bottom=3pt,boxsep=1pt,halign=left]\bfseries\scriptsize 3 ou 4 points barycentres partiels\end{tcolorbox}
\begin{tcolorbox}[enhanced,colback=popgreen,colframe=popgreen,coltext=white,boxrule=0.9pt,arc=3pt,left=4pt,right=4pt,top=3pt,bottom=3pt,boxsep=1pt,halign=left]\bfseries\scriptsize Propri\'et\'es homog\'en\'eit\'e, associativit\'e\end{tcolorbox}
\begin{tcolorbox}[enhanced,colback=poppurple,colframe=poppurple,coltext=white,boxrule=0.9pt,arc=3pt,left=4pt,right=4pt,top=3pt,bottom=3pt,boxsep=1pt,halign=left]\bfseries\scriptsize R\'eduction de $\sum\alpha_i\overrightarrow{MA_i}$\end{tcolorbox}
\begin{tcolorbox}[enhanced,colback=popgold,colframe=popgold,coltext=white,boxrule=0.9pt,arc=3pt,left=4pt,right=4pt,top=3pt,bottom=3pt,boxsep=1pt,halign=left]\bfseries\scriptsize Lignes de niveau $MA^2\pm MB^2=k$\end{tcolorbox}
\end{tcbraster}

\legende{Carte mentale de la le\c{c}on : les cinq id\'ees ma\^itresses.}
\end{popfigure}

\begin{aretenir}[L'essentiel de la le\c{c}on]
\textbf{D\'efinitions cl\'es.}
\begin{itemize}
\item $G$ est le \cle{barycentre} de $\{(A,\alpha),(B,\beta)\}$ (avec $\alpha+\beta\neq 0$) si $\alpha\overrightarrow{GA}+\beta\overrightarrow{GB}=\vec 0$, c'est-\`a-dire $\overrightarrow{AG}=\dfrac{\beta}{\alpha+\beta}\overrightarrow{AB}$.
\item G\'en\'eralisation : $G$ est le barycentre de $\{(A_i,\alpha_i)\}_{1\le i\le n}$ (avec $\sum_{i=1}^{n}\alpha_i\neq 0$) si $\sum_{i=1}^{n}\alpha_i\overrightarrow{GA_i}=\vec 0$.
\item L'\cle{isobarycentre} (tous les coefficients \'egaux) de $A$ et $B$ est le milieu de $[AB]$ ; celui de $A,B,C$ est le centre de gravit\'e du triangle $ABC$ (point de concours des m\'edianes).
\end{itemize}
\textbf{Formules et th\'eor\`emes \`a conna\^itre par c\oe ur.}
\begin{center}
\begin{tabularx}{\linewidth}{|l|X|}\hline
\rowcolor{popblueL} R\'esultat & \'Enonc\'e \\ \hline
Position de $G$ & $G\in(AB)$ ; lorsque $\alpha$ et $\beta$ sont de m\^eme signe, $G\in[AB]$ et $G$ est le plus proche du point affect\'e du plus grand coefficient en valeur absolue. \\ \hline
Homog\'en\'eit\'e & Multiplier tous les coefficients par un m\^eme r\'eel non nul ne change pas le barycentre. \\ \hline
Associativit\'e & On peut remplacer plusieurs points par leur barycentre partiel, affect\'e de la somme de leurs coefficients (si celle-ci est non nulle). \\ \hline
R\'eduction & Si $\sum\alpha_i\neq 0$ : $\sum_{i=1}^{n}\alpha_i\overrightarrow{MA_i}=\left(\sum_{i=1}^{n}\alpha_i\right)\overrightarrow{MG}$ ; si $\sum\alpha_i=0$ : la somme est un vecteur constant (ind\'ependant de $M$). \\ \hline
Coordonn\'ees & $x_G=\dfrac{\sum\alpha_i x_i}{\sum\alpha_i}$, \quad $y_G=\dfrac{\sum\alpha_i y_i}{\sum\alpha_i}$. \\ \hline
$MA^2+MB^2=k$ & Ensemble de centre $I$ milieu de $[AB]$ : cercle de rayon $\sqrt{\dfrac{k}{2}-\dfrac{AB^2}{4}}$ si $\dfrac{k}{2}-\dfrac{AB^2}{4}>0$ ; r\'eduit \`a $\{I\}$ si nul ; vide si n\'egatif. \\ \hline
$MA^2-MB^2=k$ & Droite perpendiculaire \`a $(AB)$, passant par le point $H$ de $(AB)$ tel que $\overrightarrow{IH}\cdot\overrightarrow{AB}=\dfrac{k}{2}$. \\ \hline
$\overrightarrow{MA}\cdot\overrightarrow{MB}=k$ & Cercle de centre $I$ milieu de $[AB]$ et de rayon $\sqrt{k+\dfrac{AB^2}{4}}$ (si positif) ; cas $k=0$ : cercle de diam\`etre $[AB]$. \\ \hline
$\dfrac{MA}{MB}=k$ ($k>0$) & $k=1$ : m\'ediatrice de $[AB]$ ; $k\neq 1$ : cercle (dit d'Apollonius) de diam\`etre $[G_1G_2]$, o\`u $G_1$ et $G_2$ sont les barycentres de $(A,k),(B,-1)$ et $(A,k),(B,1)$. \\ \hline
\end{tabularx}
\end{center}
\textbf{M\'ethodes express.}
\begin{itemize}
\item \emph{Construire $G$} : \'ecrire $\overrightarrow{AG}=\dfrac{\beta}{\alpha+\beta}\overrightarrow{AB}$ puis reporter la fraction de $\overrightarrow{AB}$ \`a partir de $A$.
\item \emph{Alignement} : montrer qu'un point est le barycentre de deux autres, ou exprimer un m\^eme point comme barycentre de deux syst\`emes diff\'erents.
\item \emph{Concours de droites} : montrer par associativit\'e que le m\^eme barycentre appartient \`a chacune des droites.
\item \emph{Ligne de niveau} : introduire le milieu $I$ de $[AB]$, d\'ecomposer $\overrightarrow{MA}=\overrightarrow{MI}+\overrightarrow{IA}$ et $\overrightarrow{MB}=\overrightarrow{MI}+\overrightarrow{IB}$, utiliser $\overrightarrow{IB}=-\overrightarrow{IA}$, puis conclure selon le signe de la quantit\'e obtenue.
\end{itemize}
\end{aretenir}

\begin{attention}[V\'erifications indispensables]
\begin{itemize}
\item Avant toute conclusion sur l'existence d'un barycentre, v\'erifier que la \cle{somme des coefficients est non nulle}.
\item Dans la formule $\overrightarrow{AG}=\dfrac{\beta}{\alpha+\beta}\overrightarrow{AB}$, c'est le coefficient de $B$ qui appara\^it au num\'erateur : ne pas inverser $\alpha$ et $\beta$.
\item Pour $\dfrac{MA}{MB}=k$, la condition est $k>0$ (un rapport de distances est positif) ; le cas $k=1$ donne une droite, pas un cercle.
\item Pour $MA^2+MB^2=k$, toujours discuter l'existence selon le signe de $\dfrac{k}{2}-\dfrac{AB^2}{4}$.
\end{itemize}
\end{attention}

\begin{aretenir}[Checklist avant le Probatoire]
\begin{itemize}
\item[$\square$] Je sais d\'efinir le barycentre par $\alpha\overrightarrow{GA}+\beta\overrightarrow{GB}=\vec 0$.
\item[$\square$] Je v\'erifie toujours que la somme des coefficients est non nulle avant de conclure \`a l'existence.
\item[$\square$] Je sais passer de la d\'efinition \`a $\overrightarrow{AG}=\dfrac{\beta}{\alpha+\beta}\overrightarrow{AB}$ pour construire $G$.
\item[$\square$] Je sais que l'isobarycentre de deux points est le milieu du segment.
\item[$\square$] J'utilise l'homog\'en\'eit\'e pour simplifier les coefficients (ex. diviser tous les coefficients par 2).
\item[$\square$] Je ma\^itrise la technique des barycentres partiels pour 3 ou 4 points.
\item[$\square$] Je sais prouver un alignement ou un concours de droites par le barycentre.
\item[$\square$] Je sais r\'eduire $\sum\alpha_i\overrightarrow{MA_i}$ selon que $\sum\alpha_i$ est nulle ou non.
\item[$\square$] Je connais les identit\'es $MA^2+MB^2=2MI^2+\dfrac{AB^2}{2}$ et $MA^2-MB^2=2\,\overrightarrow{IM}\cdot\overrightarrow{AB}$, o\`u $I$ est le milieu de $[AB]$.
\item[$\square$] Je sais que $\overrightarrow{MA}\cdot\overrightarrow{MB}=0$ caract\'erise le cercle de diam\`etre $[AB]$.
\item[$\square$] Je sais que $\dfrac{MA}{MB}=1$ donne la m\'ediatrice de $[AB]$.
\item[$\square$] Pour une ligne de niveau, je discute l'existence selon la valeur de $k$ (ensemble vide, point, cercle, droite).
\end{itemize}
\end{aretenir}

\findecours

\end{document}
